\documentclass[11pt]{article}
\usepackage[margin=1in]{geometry}
\usepackage{amsmath,amssymb,amsthm,mathtools}
\usepackage{enumitem}
\usepackage{microtype}
\usepackage[hidelinks]{hyperref}
\numberwithin{equation}{section}
\newtheorem{theorem}{Theorem}[section]
\newtheorem{proposition}[theorem]{Proposition}
\newtheorem{lemma}[theorem]{Lemma}
\newtheorem{corollary}[theorem]{Corollary}
\theoremstyle{definition}
\newtheorem{definition}[theorem]{Definition}
\theoremstyle{remark}
\newtheorem{remark}[theorem]{Remark}
\newcommand{\E}{\mathbb E}
\newcommand{\Pp}{\mathbb P}
\newcommand{\Var}{\operatorname{Var}}
\newcommand{\TV}{\mathrm{TV}}
\newcommand{\Unif}{\mathrm{Unif}}
\newcommand{\dTV}{d_{\TV}}
\newcommand{\nuUnif}{\nu_{\Unif}}
\newcommand{\1}{\mathbf 1}
\newcommand{\Z}{\mathbb Z}
\newcommand{\cA}{\mathcal A}
\newcommand{\cI}{\mathcal I}
\newcommand{\cT}{\mathcal T}
\newcommand{\cU}{\mathcal U}
\newcommand{\Hh}{\mathsf H}
\DeclareMathOperator{\supp}{supp}
\DeclareMathOperator{\dist}{dist}
\DeclareMathOperator{\Cov}{Cov}
\title{Time-Inhomogeneous Riffle Shuffling}
\author{}
\date{}
\makeatletter
\renewcommand{\maketitle}{\begin{center}{\LARGE\@title\par}\end{center}\vspace{1em}}
\makeatother
\begin{document}
\maketitle
\begin{abstract}
We study riffle shuffles with deterministic pile sizes whose distribution may
vary on the time scale $\log N$. We prove uniform upper and lower bounds in
total variation in terms of two quantities: the size of the largest word
class and a variational expression describing cold spots. A fixed negative
margin gives mixing, whereas a fixed positive margin gives distance tending
to one, under the stated local nondegeneracy conditions. The estimates are
first obtained directly from the finite sequence of pile sizes. They then
transfer to time-dependent profiles under scalar integrability and local
empirical approximation. The upper bound also applies to random pile sizes
by averaging. The proofs use exact sampling with prescribed column counts,
conditional exploration of the shuffle graph, and an ascent test on
localized collision sets.
\end{abstract}

\section{Introduction}

Riffle shuffling randomizes a deck of cards by cutting it into ordered
piles and combining them by a uniformly random interleaving. The pile
sizes may be deterministic or random. In~\cite{SSW}, cutoff was established
for general sequences of pile sizes whose empirical distribution converges
to a fixed nondegenerate probability measure on every macroscopic time
interval (of length $\Omega(\log N)$). In particular, the leading-order mixing time
depends only on this measure. This includes exact bisections, uniformly
random cuts, and multi-partite shuffles with general distributions of the
pile sizes.

It is natural to ask what happens when the distribution of the cuts
changes on the mixing time scale. For example, a dealer may first make
nearly balanced cuts and then switch to strongly asymmetric cuts, or use
different distributions during successive intervals of length
$\Theta(\log N)$. Such a process is no longer described by a single
limiting distribution. To study this setting, we allow the empirical
distribution near time $t\log N$ to be described by a probability measure
$\mu_t$ which depends on $t$.

We prove upper and lower bounds for these time-inhomogeneous riffle
shuffles in terms of two explicit quantities associated with the profile
$(\mu_t)$. The first records the size of the largest class of cards which
have followed the same sequence of piles. The second describes the
``cold spots'' where cards retain an unusually strong correlation with
their initial ordering. Namely, under the approximation and local
nondegeneracy conditions stated below, a fixed negative margin gives
convergence to uniform measure, while a fixed positive margin gives total
variation distance tending to one. The quantities involve costs accumulated
over subintervals of the shuffling history, accounting for the different
times at which cold spots can form. We first establish the estimates
directly from the finite sequence of pile sizes and then pass to the
profile $(\mu_t)$.

\subsection{Main results}

As our estimates hold for each fixed sequence of pile sizes, we first
consider deterministic cuts. Fix $k\ge2$ and $N\ge2$. A shuffle with
nonnegative integer pile sizes $(n_0,\ldots,n_{k-1})$, where $\sum_i n_i=N$,
takes the top $n_0$ cards to form the first pile, the next $n_1$ cards to form the
second pile, and so on. One then chooses uniformly among the
$\binom{N}{n_0,\ldots,n_{k-1}}$ interleavings which preserve the order
within each pile.

We study a sequence of such shuffles with prescribed pile sizes
$(n_0^{(j)},\ldots,n_{k-1}^{(j)})$ at step $j$, independently choosing
the interleaving at each step. Write
$p^{(j)}=(n_0^{(j)},\ldots,n_{k-1}^{(j)})/N$ for the corresponding
proportions. Starting from the identity deck, let $P_N^K$ denote the law
of the inverse of the chronological composition of the first $K$
shuffles. The stationary measure is the uniform distribution on
permutations of $N$ cards, denoted by $\nuUnif$. Our interest is in the
total variation distance $\dTV(P_N^K,\nuUnif)$. Since inversion preserves
this distance, the same estimates hold for the forward shuffle.
(All logarithms below are natural.)

To state the results, we introduce the costs associated with a single
shuffle. For $p\in D_k=\{p\in[0,1]^k:\sum_i p_i=1\}$, set
\[
 a_p=\log\frac1{\max_i p_i},\qquad
 \psi_p(\theta)=-\log\sum_i p_i^\theta.
\]
For a time profile $(\mu_t)$ of probability measures on $D_k$, let
$A(r,s)=\int_r^s\E_{\mu_t}a_p\,dt$ and
$B_\theta(r,s)=\int_r^s\E_{\mu_t}\psi_p(\theta)\,dt$ be the costs
accumulated over the logarithmic-time interval $[r,s]$. Define
\begin{align*}
 \Lambda_0(T)&=1-A(0,T),\\
 \Lambda_1(T)&=\sup_{0\le s\le u\le T}\inf_{\theta\in[3,4]}
 \left\{\frac{\theta+3}{2}\bigl(1-A(0,s)\bigr)
       -B_\theta(s,u)-2B_2(u,T)\right\}.
\end{align*}

Here $\Lambda_0$ corresponds to the largest class of equal strings, and
$\Lambda_1$ to the cold spots discussed above. The supremum over $s,u$
allows the latter obstruction to arise during any part of the shuffling
history.

We require a local nondegeneracy condition on the cuts. Call a column
$g$-good if two distinct piles each contain at least $gN$ cards. The
array is locally $(g,\gamma,\rho)$-good on $[r,T]$ if every interval
$[a,a+\rho]\subseteq[r,T]$ contains at least $\gamma\rho\log N$ good
column times $j/\log N$. Thus a positive proportion of the shuffles
in each such interval must split off two piles of linear size.

To compare the cuts with $(\mu_t)$, we use a regular finite partition
of $D_k$. The array is almost like the profile if, on every
$\rho$-grid interval, its cell frequencies approximate the integrated
profile masses, with error at most $\varphi$ times the interval length
plus $1/\log N$. The precise definitions of regularity, scalar
integrability, and almost-likeness are given in
\eqref{eq:partition-regularity}, \eqref{eq:scalar-integrability}, and
\eqref{eq:almost-like}, respectively; the same partition is used in
all three conditions.

We now state the profile bounds. If both margins are bounded above by a
fixed negative constant, then the deck is asymptotically uniform. If
either margin is bounded below by a fixed positive constant, then the
total variation distance tends to one. The proofs are given in
Sections~\ref{sec:upper-bound} and~\ref{sec:lower-bound}.

\begin{theorem}[Buffered profile upper bound]% Lean: formalization/upper_bound_v2/RiffleUpperV2/ProfileUpper/Buffered.lean :: bufferedOfEmpirical
% Lean: formalization/upper_bound_v2/RiffleUpperV2/ProfileUpper/Main.lean :: bufferedUpper
\label{thm:upper-profile}
Fix $k\ge2$, $T>0$ and $\eta>0$. There exists $\chi_0>0$ such that,
for every $0<\chi\le\chi_0$ and every regular Borel partition $\mathcal P$ of
mesh $\chi$ with regularity constant $1$, the following holds.
For every $0<g<1/k$ and $\gamma>0$, there exists $\rho_0>0$ such
that each $0<\rho\le\rho_0$ admits $\varphi_0>0$. Fix
$0<\varphi\le\varphi_0$. Let $\mu$ satisfy the scalar-integrability
condition for this partition on $[0,T]$, and suppose
\[
 \max\{\Lambda_0(T),\Lambda_1(T)\}\le-\eta.
\]
For every $\varepsilon>0$, all sufficiently large $N$ and every array
with $K=\lfloor T\log N\rfloor$ which is
$(\mathcal P,\rho,\varphi)$-almost like $\mu$ on $[0,T]$ and locally
$(g,\gamma,\rho)$-good on $[0,T]$ satisfy
\[
 \dTV(P_N^K,\nuUnif)\le\varepsilon.
\]
The partition and all meshes are fixed before $\mu$ and $N$.
\end{theorem}

\begin{theorem}[Macroscopic lower bound]% Lean: formalization/lower_bound_v2/RiffleLowerV2/MacroscopicLower/Main.lean :: RiffleLowerV2.MacroscopicLower.macroscopicIntegrable; RiffleLowerV2.MacroscopicLower.macroscopicIntegrable_tendsto
\label{thm:macroscopic-lower}
Fix $k\ge2$, $T>0$ and $\eta>0$. There are $0<\chi<1/k$, a $\chi$-regular finite Borel
partition $\mathcal R$ with regularity constant $1$, and $\rho_0>0$ such that
\[
 2\rho_0\log k\le\frac1{100}\min\{\eta/200000,1/1000\}.
\]
For every $0<\rho\le\rho_0$ there exists $\varphi_0>0$ with the following
property. Fix $0<\varphi\le\varphi_0$, $\chi_g,\gamma>0$, and a time
profile $\mu$ satisfying the scalar integrability condition~\eqref{eq:scalar-integrability} on $[0,T]$ for $\mathcal R$. If
\[
 \max\{\Lambda_0(T),\Lambda_1(T)\}\ge\eta,
\]
then, for every $\varepsilon>0$, there is $N_0\ge2$ such that all
$N\ge N_0$ and all cut arrays with $K=\lfloor T\log N\rfloor$ which are
$(\mathcal R,\rho,\varphi)$-almost like $\mu$ on $[0,T]$ and satisfy the local
good-column condition with parameters $\chi_g,\gamma,\rho$ on $[0,T]$
have $\dTV(P_N^K,\nuUnif)\ge1-\varepsilon$. In particular the distance
tends to one for every sequence of arrays eventually satisfying these
fixed hypotheses, even if the applicable empirical alternative varies
with $N$.
\end{theorem}

In fact, the main estimates can be stated directly in terms of the
finite sequence of cuts, without first choosing a limiting profile.
These estimates also give the uniformity needed to deduce
Theorems~\ref{thm:upper-profile} and~\ref{thm:macroscopic-lower}.
Put $S_N=K/\log N$ and let $p_N(t)=p^{(j)}$ on
$((j-1)/\log N,j/\log N)$. Define the empirical costs by
\[
 A^{(N)}(r,s)=\int_r^s a_{p_N(t)}\,dt,\qquad
 B^{(N)}_\theta(r,s)=\int_r^s\psi_{p_N(t)}(\theta)\,dt
 \quad(0\le r\le s\le S_N),
\]
and
\begin{align*}
 \beta^{(N)}(s)&=1-A^{(N)}(0,s),\\
 \Lambda_1^{(N)}&=\sup_{0\le s\le u\le S_N}\inf_{\theta\in[3,4]}
 \left\{\frac{\theta+3}{2}\beta^{(N)}(s)
       -B^{(N)}_\theta(s,u)-2B^{(N)}_2(u,S_N)\right\}.
\end{align*}
We say the array has empirical margin at least $\eta$ for the upper bound if
$\beta^{(N)}(S_N)\le-\eta$ and the infimum in this display
is at most $-\eta$ for every grid pair
$s=i/\log N$, $u=j/\log N$, $0\le i\le j\le K$.
(For the lower bound, $\Lambda_1^{(N)}$ uses all real endpoints as in
the displayed definition.) With this notation, the finite-array bounds
are as follows.

\begin{theorem}[Empirical upper bound]% Lean: formalization/upper_bound_v2/RiffleUpperV2/EmpiricalUpper/Main.lean :: empiricalUpper
% Lean: formalization/upper_bound_v2/RiffleUpperV2/EmpiricalUpper/ExplorationTV.lean :: concreteExplorationTV
\label{thm:upper-empirical}
Fix $k\ge2$ and $T,\eta_0,g_0,\gamma_0>0$. There exists $\rho_0>0$
such that, for every fixed $0<\rho\le\rho_0$, there exist $c>0$ and
$N_0\ge2$ with the following property. For $N\ge N_0$ and
$K\le T\log N$, suppose the deterministic array on $N$ labels has
empirical margin at least $\eta_0$ and satisfies the local
$(g_0,\gamma_0,\rho)$-good condition on $[0,K/\log N]$. Then
\begin{equation}\label{eq:upper-empirical-tv}
 \dTV(P_N^K,\nuUnif)\le N^{-c}.
\end{equation}
The choices of $c,N_0$ precede the cut array and the horizon $K$.
\end{theorem}

\begin{theorem}[Empirical lower bound]% Lean: formalization/lower_bound_v2/RiffleLowerV2/EmpiricalLower/Main.lean :: RiffleLowerV2.EmpiricalLower.empiricalCombined
\label{thm:empirical-lower}
Fix $k\ge2$, $T>0$, positive margins $\eta_0,\eta_1$ and arbitrary real parameters
$\chi_g,\gamma,\rho$. For every $\varepsilon>0$, there is $N_0\ge2$ such
that all $N\ge N_0$ and all deterministic cut arrays with $K=\lfloor T\log N\rfloor$ have
$\dTV(P_N^K,\nuUnif)\ge1-\varepsilon$ whenever either
\[
 \beta^{(N)}(S_N)\ge\eta_0,
\]
or all of the following hold:
\[
 \chi_g,\gamma,\rho>0,\qquad
 2\rho\log k\le\frac1{100}\min\{\eta_1/100000,1/1000\},\qquad
 \Lambda_1^{(N)}\ge\eta_1,
\]
and the cut array satisfies the local good-column condition with
parameters $\chi_g,\gamma,\rho$ on $[0,T]$.
\end{theorem}

\subsection{Consequences and Discussion}

The two margins describe different ways in which the initial order can
persist. For $\Lambda_0$, a single class of equal strings is large
enough to produce a long increasing run. The cold spots detected by
$\Lambda_1$ may instead contain many smaller classes, whose combined
contribution can be distinguished from a uniformly random deck. Thus the
cold spot picture continues to describe the obstruction to mixing when
the distribution of the cuts changes with time.

The profile assumptions concern only the scalar quantities used in the
proof: the cell probabilities of the selected partition, the maximal
cost, and the power-sum costs for $\theta\in[2,4]$. Integrability of
these functions suffices for the approximation argument. In particular,
we do not need to assume that $(\mu_t)$ is a measurable probability
kernel.

We also obtain an upper bound on the mixing time. To state it, suppose
the cuts for each $N$ are given by one infinite sequence, so that the
processes at different horizons are consistent. The following corollary
specifies the approximation condition as the spatial and temporal
meshes tend to zero.

\begin{corollary}[Critical-time upper bound]% Lean: formalization/upper_bound_v2/RiffleUpperV2/Shared/Profile.lean :: ArbitrarilyFineApproximation; criticalTime; mixingTime
% Lean: formalization/upper_bound_v2/RiffleUpperV2/ProfileUpper/Critical.lean :: criticalOfBuffered; criticalTimeUpperOfBuffered; mixingTimeUpperOfConvergence
\label{cor:upper-critical}
Suppose a consistent infinite sequence of cuts is specified for each
$N$, and these arrays admit arbitrarily fine approximations to $\mu$:
for each $T>0$ and each mesh bound $\chi_*>0$ there are
$0<\chi\le\chi_*$, a $\chi$-regular partition with constant $1$ satisfying
scalar integrability, and $g,\gamma,\bar\rho>0$, with $g<1/k$,
such that for every $0<\rho\le\bar\rho$ and every $\varphi>0$ the
arrays at horizon $\lfloor T\log N\rfloor$ are eventually almost
like $\mu$ and locally good with these parameters. For every $t\ge0$ with
\[
 t>\inf\{T>0:\max\{\Lambda_0(T),\Lambda_1(T)\}<0\},
 \qquad \inf\varnothing=+\infty,
\]
we have $\dTV(P_N^{\lfloor t\log N\rfloor},\nuUnif)\to0$.
Moreover, for every $0<\varepsilon<1$,
\[
 \limsup_{N\to\infty}\frac{t_{\rm mix}^{(N)}(\varepsilon)}{\log N}
       \le\inf\{T>0:\max\{\Lambda_0(T),\Lambda_1(T)\}<0\}.
\]
Here $t_{\rm mix}^{(N)}(\varepsilon)$ is the infimum of the integers
$m\ge0$ for which $\dTV(P_N^m,\nuUnif)\le\varepsilon$.
Both infima are interpreted in $[0,+\infty]$, with empty infimum
$+\infty$.
\end{corollary}

On the other hand, Theorem~\ref{thm:macroscopic-lower} gives distance
tending to one at every time with a strictly positive margin and the
specified approximation conditions. Both conclusions retain a fixed
buffer from zero, as required by the finite estimates.

Finally, uniformity over the deterministic cuts yields an upper bound
for random pile sizes. Namely, condition on the entire sequence of pile
sizes and apply the empirical bound whenever the margin and local-good
conditions hold. Convexity of total variation bounds the averaged
distance by the deterministic error plus the probability that these
conditions fail. Hence they need only hold with probability tending to
one, with common parameters. The pile sizes at different times may be
dependent. We give the precise statement in
Corollary~\ref{cor:upper-random}.

\subsection{Proof Techniques and Challenges}

The proofs split into lower and upper bounds, following the cold spot
and truncated exponential moment approach in~\cite{SSW}. We discuss
these separately, beginning with the lower bound to explain the two
obstructions appearing in the main results.

\paragraph{Lower Bound via Cold Spot Hypothesis Testing.}
To describe the obstruction to mixing, identify each card with a string
in $\{0,1,\ldots,k-1\}^K$ recording its pile labels. As explained in
Lemma~\ref{lem:matrix-graph}, the inverse shuffled permutation is obtained
by sorting these strings lexicographically, with equal strings retaining
their initial order. A class of equal strings therefore produces an
increasing run. A word obtained by choosing a maximal pile in each column
has mean occupancy $N^{\beta^{(N)}(S_N)}$; when this mean grows as a
positive power of $N$, concentration
gives a run which is unlikely to occur in a uniformly random permutation.
This proves the lower bound corresponding to $\Lambda_0$.

The more delicate obstruction comes from cold spots. When the maximal-word
argument does not already apply, a fixed positive lower bound on
$\Lambda_1^{(N)}$ supplies times $s\le u$ and a tilt
$\theta\in[3,4]$ with enough slack to construct a family of strings
whose equal-word classes have approximately deterministic locations,
relative to a common prefix. To construct this family, we approximate
the actual column proportions by a finite collection of representatives
and prescribe digit frequencies proportional to their $\theta$th powers
between $s$ and $u$. We use a second tilted class in the tail and insert
a short bridge containing two good columns. The bridge supplies strings
on either side of the selected family, so fluctuations of the endpoints
do not move its cards outside the intended cold set.

There are two probabilistic issues in this construction. First, the
column counts are prescribed, so different cards do not receive their
digits independently. We use exact falling-factorial moments to estimate
the number of selected strings which occur at least twice. The estimates
are uniform even when the mean occupancy of a string is very small or
very large (see Lemma~\ref{lem:two-category-poisson}). Second, the common
prefix has a random location in the deck. Successive hypergeometric
thinning controls the positions relative to this prefix, giving a
translate of a deterministic cold set. We then count ascents over all
translates to obtain a hypothesis test. The translation is chosen from
the shuffle matrix before the independent uniform permutation in the
component-sorting representation is sampled. Concentration for component
sorting and a uniform bound for the translate scan complete the lower
bound.

\paragraph{Upper Bound via Truncated Exponential Moments.}
For the upper bound, join adjacent positions whenever their sorted
strings agree. This gives a random subgraph $G$ of the path on $N$
vertices. The inverse shuffle has the same law as an independent uniform
permutation sorted within the components of $G$. A truncated second
moment argument then reduces the problem to controlling an exponential
moment of the number of shared edges of two independent shuffle graphs.
As in~\cite{SSW}, we first estimate the expected number of shared edges
and then prove a conditional version of this estimate during an
exploration.

For the first moment, we may assume the initial shuffle is good, so
that each of its piles leaves a positive fraction of the cards outside
it. These outside cards provide an unexposed population for the
conditioning arguments. First, we study the fluctuation in the location
of a prefix block. By choosing which coordinates to reveal, its
endpoints become translates of a hypergeometric variable whose variance
can be bounded below. Anti-concentration for this variable spreads the
edge probabilities over sufficiently many positions. Second, we divide
the strings into prefix blocks whose sizes are controlled by their
endpoint fluctuations. Extremely small pile proportions require
care: we record the positions and digits of these exceptional choices
separately, and group the remaining choices by their accumulated costs.
The empirical margin bounds the contribution of each class, while
localization controls the number of blocks contributing at one position.
This gives the first-moment estimate in
Proposition~\ref{prop:upper-first}.

Having established this estimate, we explore the two graphs from both
directions within each first-column pile, giving $2k$ exploration
processes. The goal is to bound the expected number of remaining shared
edges conditionally on the words already revealed. We decompose the
unrevealed strings into disjoint prefix blocks and refine them so that
each sufficiently large block is followed by a good column. After
revealing the row sets and induced column counts, the internal law is
again a product of uniform fixed-count column laws. Concentration shows
that the induced proportions remain close enough to the original ones
for the first-moment estimate to apply. For small blocks, we instead
estimate their sizes and collision counts directly by counting
completions of the partially exposed matrix. An adapted moment argument
then converts these conditional bounds into the required exponential
estimate.

The rest of the paper is organized as follows. Section~2 introduces
the notation, profile approximation, and shuffle graph representation.
Section~3 proves the upper bound from the moment estimates and deduces
its consequences. Sections~4 and~5 establish the first-moment and
conditional exploration estimates, respectively. Finally, Section~6
proves the lower bound.


\section{Preliminaries}
\label{sec:preliminaries}

\subsection{Pile sizes and logarithmic time}

Write $[N]=\{1,\ldots,N\}$ and $[k]_0=\{0,\ldots,k-1\}$. Throughout,
$k\ge2$, $N\ge2$, and all logarithms are natural. A cut array consists of
nonnegative integers $n_i^{(j)}$, $1\le j\le K$, $i\in[k]_0$, satisfying
$\sum_i n_i^{(j)}=N$. Put
\[
 D_k=\left\{p\in[0,1]^k:\sum_i p_i=1\right\},\qquad
 p^{(j)}=\frac1N(n_0^{(j)},\ldots,n_{k-1}^{(j)}),\qquad
 S_N=\frac K{\log N}.
\]
Column $j$ is placed at time $j/\log N$. Its contribution to an interval
with arbitrary endpoints is defined by linear interpolation. More
precisely, for a function $f$ on $D_k$, set
\begin{equation}% Lean: formalization/lower_bound_v2/RiffleLowerV2/Shared/Empirical.lean :: columnFraction;empiricalIntegral;gridSum
\label{eq:empirical-integral}
 \int_r^s f(p_N(t))\,dt
 =\frac1{\log N}\sum_{j=1}^K
 \left[\bigl((s\log N-j+1)\vee0\bigr)\wedge1
       -\bigl((r\log N-j+1)\vee0\bigr)\wedge1\right]f(p^{(j)}).
\end{equation}
Equivalently, $p_N(t)=p^{(j)}$ on
$((j-1)/\log N,j/\log N)$; the integrand is extended by zero outside
$[0,S_N]$. The explicit formula specifies the convention at every
endpoint. In particular, on grid intervals this is exactly the normalized
sum of the corresponding column costs.

For $T\ge0$ we use $K=\lfloor T\log N\rfloor$. Then
\begin{equation}% Lean: formalization/time_profile/RiffleTimeProfile/Base.lean :: horizonDepth_error
\label{eq:horizon-rounding}
 0\le T-S_N<\frac1{\log N}.
\end{equation}
The distinction between $T$ and $S_N$ will be retained in the estimates.

\subsection{Scalar costs and margins}
\label{subsec:scalar-costs}

For $p\in D_k$ and $\theta\ge1$, write
\[
 p_{\max}=\max_i p_i,\qquad a_p=\log\frac1{p_{\max}},\qquad
 \phi_p(\theta)=\sum_i p_i^\theta,\qquad
 \psi_p(\theta)=-\log\phi_p(\theta).
\]
Zero coordinates contribute zero to the power sum. The tilted
probabilities, cost, and entropy are
\begin{align}% Lean: formalization/time_profile/RiffleTimeProfile/Base.lean :: tiltedWeight;tiltedCost;shannonEntropy;tiltedEntropy
\label{eq:tilted-costs}
 p_i^{[\theta]}&=\frac{p_i^\theta}{\phi_p(\theta)},\qquad
 \sum_i p_i^{[\theta]}=1,\nonumber\\
 I_p(\theta)&=\sum_{i:p_i>0}p_i^{[\theta]}\log\frac1{p_i},\qquad
 H_p(\theta)=-\sum_{i:p_i^{[\theta]}>0}p_i^{[\theta]}\log p_i^{[\theta]}.
\end{align}
Thus no logarithm of a zero coordinate is used in a nonzero summand.

Define the empirical costs by
\begin{align*}
 A^{(N)}(r,s)&=\int_r^s a_{p_N(t)}\,dt,&
 B^{(N)}_\theta(r,s)&=\int_r^s\psi_{p_N(t)}(\theta)\,dt,\\
 J^{(N)}_\theta(r,s)&=\int_r^s I_{p_N(t)}(\theta)\,dt,&
 H^{(N)}_\theta(r,s)&=\int_r^s H_{p_N(t)}(\theta)\,dt.
\end{align*}
The two empirical margins are
\begin{align}% Lean: formalization/lower_bound_v2/RiffleLowerV2/Shared/Empirical.lean :: A;B;J;H;beta;F;lambda1
\label{eq:empirical-margins}
 \beta^{(N)}(s)&=1-A^{(N)}(0,s),\nonumber\\
 \Lambda_1^{(N)}&=\sup_{0\le s\le u\le S_N}\inf_{\theta\in[3,4]}
 \left\{\frac{\theta+3}{2}\beta^{(N)}(s)
       -B^{(N)}_\theta(s,u)-2B^{(N)}_2(u,S_N)\right\}.
\end{align}

A time profile is a family $(\mu_t)_{t\in\mathbb R}$ of probability
measures on $D_k$. No regularity in $t$ is implicit in this definition.
Whenever the scalar functions concerned are integrable, write
\[
 A(r,s)=\int_r^s\int_{D_k}a_p\,d\mu_t(p)\,dt,\qquad
 B_\theta(r,s)=\int_r^s\int_{D_k}\psi_p(\theta)\,d\mu_t(p)\,dt.
\]
Write $\beta(s)=1-A(0,s)$, so that $\Lambda_0(T)=\beta(T)$.
Recall the margins $\Lambda_0(s)=\beta(s)$ and
\begin{equation}% Lean: formalization/lower_bound_v2/RiffleLowerV2/Shared/Empirical.lean :: macroBeta;macroF;macroLambda1
\label{eq:macro-margin}
 \Lambda_1(T)=\sup_{0\le s\le u\le T}\inf_{\theta\in[3,4]}
 \left\{\frac{\theta+3}{2}\Lambda_0(s)
             -B_\theta(s,u)-2B_2(u,T)\right\}.
\end{equation}
The smaller expressions arising in the first-moment argument are
\begin{align}% Lean: formalization/upper_bound_v2/RiffleUpperV2/ProfileUpper/Dominance.lean :: macroF2;macroLambda2;macroLambda3
\label{eq:auxiliary-margins}
 \Lambda_2(T)&=\sup_{0\le s\le u\le T}\inf_{\theta\in[3,4]}
 \left\{\frac{\theta-1}{2}\Lambda_0(s)
             -B_\theta(s,u)-2B_2(u,T)\right\},\nonumber\\
 \Lambda_3(T)&=\sup_{0\le s\le T}\{\Lambda_0(s)-B_2(s,T)\}.
\end{align}
The supremum and infimum conventions in these definitions are literal;
their boundedness follows from the next lemma.

\begin{lemma}[One-column bounds]% Lean: formalization/time_profile/RiffleTimeProfile/Analysis.lean :: powerSum_pos;powerSum_le_one;maxCost_bounds;psi_bounds
% Lean: formalization/lower_bound_v2/RiffleLowerV2/ProfileTransfer/Spatial.lean :: continuous_maxCost;jointPsiContinuity;integrable_maxCost;integrable_psi
\label{lem:scalar-bounds}
For $p\in D_k$ and $2\le\theta\le4$,
\[
 0\le a_p\le\log k,\qquad
 0<\phi_p(\theta)\le1,\qquad
 0\le\psi_p(\theta)\le4\log k.
\]
Moreover, for every $\theta\ge1$,
\begin{equation}% Lean: formalization/upper_bound_v2/RiffleUpperV2/ProfileUpper/Dominance.lean :: basicOneColumn;profileDominance;integratedDominance
\label{eq:one-column-dominance}
 (\theta-1)a_p\le\psi_p(\theta).
\end{equation}
The function $a_p$ is continuous on $D_k$, and
$(p,\theta)\mapsto\psi_p(\theta)$ is continuous on $D_k\times[2,4]$.
\end{lemma}

\begin{proof}
We have $1/k\le p_{\max}\le1$, giving the bounds on $a_p$.
At least one coordinate is positive, so $\phi_p(\theta)>0$.
Since $p_i^\theta\le p_i$, the power sum is at most one. Choosing a
coordinate with $p_i\ge1/k$ gives
$\phi_p(\theta)\ge p_i^\theta\ge k^{-\theta}$ and hence the stated
upper bound on $\psi_p(\theta)$. Also,
\[
 \phi_p(\theta)=\sum_i p_i^{\theta-1}p_i
 \le p_{\max}^{\theta-1}\sum_i p_i=p_{\max}^{\theta-1}.
\]
Taking minus logarithms proves~\eqref{eq:one-column-dominance}; when
$\theta=1$ both sides are zero. Continuity of the maximum gives
continuity of $a_p$. Each summand $p_i^\theta$ is continuous on the
displayed domain, including at $p_i=0$, and the sum stays positive.
Composition with $-\log$ proves the last assertion.
\end{proof}

The same bounds hold after integration against $\mu_t$. In particular,
all spatial integrals above exist. Temporal integrability remains an
explicit hypothesis.

\begin{lemma}[Reduction and monotonicity of the margins]
% Lean: formalization/upper_bound_v2/RiffleUpperV2/ProfileUpper/Dominance.lean :: lambda2_le_beta;lambda3_le_beta;dominanceReduction
% Lean: formalization/upper_bound_v2/RiffleUpperV2/ProfileUpper/AllMargins.lean :: lambda3_eq_beta;allMarginMonotonicity
\label{lem:margin-dominance}
Suppose $T\ge0$ and the functions
$t\mapsto\int a_p\,d\mu_t(p)$ and
$t\mapsto\int\psi_p(\theta)\,d\mu_t(p)$, $\theta\in[2,4]$, are
integrable on $[0,T]$. Then
\begin{equation}\label{eq:dominance-reduction}
 \Lambda_2(T)\le\Lambda_0(T),\qquad
 \Lambda_3(T)=\Lambda_0(T),\qquad
 \max_{0\le i\le3}\Lambda_i(T)=\max\{\Lambda_0(T),\Lambda_1(T)\}.
\end{equation}
Every $\Lambda_i$, $0\le i\le3$, is nonincreasing on $[0,T]$.
For $0\le S\le T$,
\begin{equation}% Lean: formalization/upper_bound_v2/RiffleUpperV2/ProfileUpper/Analytic.lean :: macroLambda_horizon;macroF_bound;macroF_clip_le;macroF_horizon_eq
\label{eq:lambda-horizon}
 0\le\Lambda_1(S)-\Lambda_1(T)\le8\log k\,(T-S).
\end{equation}
\end{lemma}

\begin{proof}
Integrating~\eqref{eq:one-column-dominance} gives
\[
 (\theta-1)A(r,s)\le B_\theta(r,s),\qquad 0\le r\le s\le T.
\]
Thus $\Lambda_0(s)-B_2(s,T)\le\Lambda_0(T)$; equality is obtained at
$s=T$, proving the assertion about $\Lambda_3$. In the infimum defining
$\Lambda_2$, use $\theta=3$. Its value is at most
\[
 \Lambda_0(s)-B_3(s,u)-2B_2(u,T)
 \le\Lambda_0(s)-2A(s,u)-2A(u,T)
 \le\Lambda_0(T).
\]
Taking the supremum proves~\eqref{eq:dominance-reduction}.

Nonnegativity of $a_p$ makes $\Lambda_0$ nonincreasing. To compare
$\Lambda_1(T)$ with $\Lambda_1(S)$, take an admissible pair $(s,u)$ at
time $T$ and replace it by $(s\wedge S,u\wedge S)$. All three intervals
of integration are shortened. Since their integrands are nonnegative
and $(\theta+3)/2\ge0$, the expression at time $T$ is no larger than
the expression at time $S$ for this clipped pair. Infima and then
suprema give $\Lambda_1(T)\le\Lambda_1(S)$. The same argument, with
$(\theta-1)/2$ in place of $(\theta+3)/2$, proves monotonicity of
$\Lambda_2$. The equality already proved handles $\Lambda_3$.

For the other direction in~\eqref{eq:lambda-horizon}, a pair admissible
at time $S$ is also admissible at time $T$, and the two corresponding
expressions differ by exactly $2B_2(S,T)$, independently of $\theta$.
Consequently
\[
 \Lambda_1(S)\le\Lambda_1(T)+2B_2(S,T)
 \le\Lambda_1(T)+8\log k\,(T-S).
\]
All uses of infima and suprema are legitimate: the expressions have
absolute value at most $4+20T\log k$ on the indicated domains, by
Lemma~\ref{lem:scalar-bounds}.
\end{proof}

\subsection{Local approximation of a time profile}
\label{subsec:profile-approximation}

A finite Borel partition $\mathcal P=\{D_1,\ldots,D_m\}$ of $D_k$
includes a representative $p^{[i]}\in D_i$ for every cell. We call it
$\chi$-regular if, for all $p,p'\in D_i$ and $\theta\in[2,4]$,
\begin{align}% Lean: formalization/time_profile/RiffleTimeProfile/Base.lean :: BorelPartition;BorelPartition.IsRegular
\label{eq:partition-regularity}
 |a_p-a_{p'}|&\le\chi,&
 |\psi_p(\theta)-\psi_{p'}(\theta)|&\le\chi,\nonumber\\
 p^{[i]}_l>2\chi&\ \Longrightarrow\
 p_l>\chi\ \text{ and }\left|\log\frac1{p_l}-\log\frac1{p^{[i]}_l}\right|\le\chi,
 \nonumber\\
 p^{[i]}_l\le2\chi&\ \Longrightarrow\ p_l\le3\chi.
\end{align}

\begin{lemma}[Regular partitions]% Lean: formalization/time_profile/RiffleTimeProfile/Analysis.lean :: exists_regular_partition
\label{lem:regular-partition}
For every $0<\chi<1/k$, there is a finite $\chi$-regular Borel
partition of $D_k$.
\end{lemma}

\begin{proof}
Besides the two continuous functions in Lemma~\ref{lem:scalar-bounds},
consider the coordinate functions
$p\mapsto\log(1/(p_l\vee\chi))$. Compactness gives a common spatial
modulus of continuity, uniform in $\theta\in[2,4]$ and in the finitely
many coordinates. Partition $[0,1]^k$ into half-open coordinate boxes
of sufficiently small side length, and intersect with $D_k$, keeping
the nonempty intersections. Choose a representative in each cell.
The side length can be chosen below $\chi$, with oscillation of every
displayed continuous function at most $\chi$ on a box. If
$p^{[i]}_l>2\chi$, coordinate closeness gives $p_l>\chi$, so the
truncated logarithms are the actual logarithms. If
$p^{[i]}_l\le2\chi$, coordinate closeness gives $p_l\le3\chi$.
These are precisely~\eqref{eq:partition-regularity}.
\end{proof}

For a chosen partition, the scalar integrability condition on $[0,T]$
means that each of
\begin{equation}% Lean: formalization/lower_bound_v2/RiffleLowerV2/Shared/Empirical.lean :: ScalarIntegrable
\label{eq:scalar-integrability}
 t\longmapsto\mu_t(D_i),\qquad
 t\longmapsto\int a_p\,d\mu_t(p),\qquad
 t\longmapsto\int\psi_p(\theta)\,d\mu_t(p)
 \quad(1\le i\le m,\ 2\le\theta\le4)
\end{equation}
is Lebesgue integrable on $[0,T]$.

For $\rho,\varphi>0$, an array with $K=\lfloor T\log N\rfloor$ is
$(\mathcal P,\rho,\varphi)$-almost like $(\mu_t)$ if, for every
$a\in\rho\mathbb Z$ with $0\le a<b=\min(a+\rho,T)$ and every cell,
\begin{equation}% Lean: formalization/time_profile/RiffleTimeProfile/Base.lean :: AlmostLike;intervalCount
\label{eq:almost-like}
 \left|\frac{\#\{j:a\le j/\log N\le b,\ p^{(j)}\in D_i\}}{\log N}
       -\int_a^b\mu_t(D_i)\,dt\right|
 \le\varphi(b-a)+\frac1{\log N}.
\end{equation}
Closed endpoints in this counting condition are intentional. A column
is $g$-good if two distinct piles each contain at least $gN$ cards.
The local good-column condition on $[r,T]$, with parameters
$g,\gamma,\rho>0$, is
\begin{equation}% Lean: formalization/time_profile/RiffleTimeProfile/Base.lean :: GoodColumn;LocalGood
\label{eq:local-good}
 \#\{j:a\le j/\log N\le a+\rho,\ p^{(j)}\text{ is }g\text{-good}\}
 \ge\gamma\rho\log N
 \quad(r\le a,\ a+\rho\le T).
\end{equation}
Here $a$ ranges over all real numbers in the indicated range, not only
the grid $\rho\mathbb Z$.

\begin{lemma}[Uniform scalar approximation]% Lean: formalization/lower_bound_v2/RiffleLowerV2/ProfileTransfer/Quadrature.lean :: sampledQuadrature;quadrature_budget;genericQuadrature
% Lean: formalization/lower_bound_v2/RiffleLowerV2/ProfileTransfer/Compat/GridQuadrature.lean :: representative_grid_block;one_sided_quadrature
% Lean: formalization/lower_bound_v2/RiffleLowerV2/ProfileTransfer/Endpoint.lean :: empiricalIntegral_sample_error
\label{lem:scalar-quadrature}
Let $T,\rho>0$, $\varphi,\omega,M\ge0$, and let
$\mathcal P=\{D_1,\ldots,D_m\}$ be a finite Borel partition with
representatives. Suppose $f:D_k\to[0,M]$ is continuous and has
oscillation at most $\omega$ on each cell. Assume integrability of
$t\mapsto\mu_t(D_i)$ and $t\mapsto\int f\,d\mu_t$ on $[0,T]$.
If~\eqref{eq:almost-like} holds, then, uniformly for
$0\le r\le s\le S_N$,
\begin{align}\label{eq:quadrature-bound}
 \left|\int_r^s f(p_N(t))\,dt-
              \int_r^s\int f\,d\mu_t\,dt\right|
 \le{}&2T\omega+Mm\varphi T+4M\rho\nonumber\\
 &+\frac{4Mm(T/\rho+2)+6M}{\log N}.
\end{align}
\end{lemma}

\begin{proof}
We first use the samples with $r<j/\log N\le s$. Passing from the
closed counting intervals in~\eqref{eq:almost-like} to these
half-open intervals removes at most one column. The error for the
count in each cell of a full time interval of length $\rho$ is
therefore at most $\varphi\rho+2/\log N$. The total number of samples
in an interval of length $\rho$, divided by $\log N$, differs from
$\rho$ by at most $1/\log N$.

Put $w=\min(\omega,M)$. On a full grid interval, replace both the
spatial observable and the sampled observable by their cell
representatives. The first replacement costs at most $w\rho$.
The cell-count estimate contributes at most
$M/\log N+Mm(\varphi\rho+2/\log N)$; the additional
$M/\log N$ also bounds a possible discrepancy in the total cell
mass. Replacing the representative samples back by actual samples
costs at most $w(\rho+1/\log N)$. Hence the integral is bounded above
by the sampled sum, normalized by $\log N$, plus
\[
 (2w+Mm\varphi)\rho+\frac{w+M+2Mm}{\log N}.
\]
Sum this inequality over the full grid intervals contained in
$[r,s]$. There are at most $\lfloor T/\rho\rfloor+1$ of them.
The two remaining portions have total length at most $2\rho$,
and their integrals are at most $2M\rho$. By nonnegativity, any
additional samples can be retained on the right. We obtain a
one-sided error bounded by
\[
 2wT+Mm\varphi T+2M\rho+
 \frac{(w+M+2Mm)(\lfloor T/\rho\rfloor+1)}{\log N}.
\]
Apply the same argument to $M-f$. Subtracting its integral and sum
from those of the constant $M$ gives the reverse inequality, with
an additional error at most $M/\log N$. This subtraction is valid
by the assumed temporal integrability.

Finally, the interpolated integral~\eqref{eq:empirical-integral}
differs from the sampled sum by at most $2M/\log N$: only the two
endpoint columns can be partial. Thus the absolute error is at most
the last display plus $3M/\log N$. Since $m\ge1$, $w\le M$,
$w\le\omega$, and $\lfloor T/\rho\rfloor\le T/\rho$, this is
bounded by~\eqref{eq:quadrature-bound}.
\end{proof}

\begin{lemma}[Transfer of the margins]% Lean: formalization/lower_bound_v2/RiffleLowerV2/ProfileTransfer/Margin.lean :: common_F_error;common_lambda_error;marginTransfer
\label{lem:margin-transfer}
Let $T>0$, and assume~\eqref{eq:scalar-integrability} for a fixed
partition. Suppose an array with $K=\lfloor T\log N\rfloor$ satisfies,
for every $0\le r\le s\le S_N$ and $2\le\theta\le4$,
\begin{equation}\label{eq:interval-transfer}
 |A^{(N)}(r,s)-A(r,s)|\le e,\qquad
 |B^{(N)}_\theta(r,s)-B_\theta(r,s)|\le e,
 \qquad e\ge0.
\end{equation}
Then
\begin{align}\label{eq:margin-transfer}
 |\beta^{(N)}(S_N)-\Lambda_0(T)|&\le e+\frac{\log k}{\log N},\nonumber\\
 |\Lambda_1^{(N)}-\Lambda_1(T)|&\le\frac{13}{2}e+
                                          \frac{14\log k}{\log N}.
\end{align}
\end{lemma}

\begin{proof}
The first estimate follows by splitting $A(0,T)$ at $S_N$ and
using~\eqref{eq:horizon-rounding} and
$0\le A(S_N,T)\le(T-S_N)\log k$. For the second, first compare
the variational expressions on the common triangle
$0\le s\le u\le S_N$. The coefficient $(\theta+3)/2$ lies in
$[3,7/2]$, so~\eqref{eq:interval-transfer} bounds their difference
by $(7/2+1+2)e=13e/2$, uniformly in $s,u,\theta$. Taking an
infimum and then a supremum preserves this absolute bound.
Lemma~\ref{lem:margin-dominance} bounds the remaining change of
horizon by $8\log k\,(T-S_N)\le8\log k/\log N$. Enlarging $8$
to $14$ gives exactly~\eqref{eq:margin-transfer}.
\end{proof}

\begin{proposition}[Uniform profile transfer]% Lean: formalization/lower_bound_v2/RiffleLowerV2/ProfileTransfer/Main.lean :: profileTransfer
% Lean: formalization/lower_bound_v2/RiffleLowerV2/ProfileTransfer/Parameters.lean :: transferChi_spec;transferRho_spec;transferPhi_spec;transferThreshold_spec;transfer_error_budget
\label{prop:profile-transfer}
Fix $T,e>0$ and an arbitrary positive upper bound on $\rho$.
There exist $0<\chi<1/k$, a finite $\chi$-regular partition
$\mathcal P$, and $\rho_0>0$ below this upper bound such that, for
every $0<\rho\le\rho_0$, there is $\varphi_0>0$ with the following
property. For each $0<\varphi\le\varphi_0$ and each profile satisfying
\eqref{eq:scalar-integrability} for this same partition, all
sufficiently large $N$ and all arrays satisfying~\eqref{eq:almost-like}
obey both~\eqref{eq:interval-transfer} and~\eqref{eq:margin-transfer}.
The partition is chosen before the profile and before $N$; the
large-$N$ threshold is uniform over the cut arrays.
\end{proposition}

\begin{proof}
Choose
\[
 \chi=\min\left\{\frac1{2k},\frac e{8(T+1)}\right\}
\]
and a regular partition by Lemma~\ref{lem:regular-partition}. In
Lemma~\ref{lem:scalar-quadrature}, apply the same bound $M=4\log k$
to $a_p$ and to all $\psi_p(\theta)$, $2\le\theta\le4$, with
$\omega=\chi$. Choose $\rho_0$ no larger than the prescribed upper
bound and $e/(16(4\log k+1))$, and choose
\[
 \varphi_0=\frac e{4(4mT\log k+1)}.
\]
The first three terms on the right of~\eqref{eq:quadrature-bound}
are each at most $e/4$. After these choices, increase $N$ until its
last term is at most $e/4$. This proves the interval estimates
simultaneously for every endpoint pair and every $\theta\in[2,4]$.
Lemma~\ref{lem:margin-transfer} now applies. Every choice preceding
$N$ is independent of the cut array, as required.
\end{proof}

\begin{corollary}[Transfer for a fixed partition]% Lean: formalization/upper_bound_v2/RiffleUpperV2/Margins/FixedPartition.lean :: fixedPartitionQuadrature;anyFixedPartitionTransfer
\label{cor:fixed-partition-transfer}
Fix $T,e>0$. There is $\chi_0>0$ such that, for every
$0<\chi\le\chi_0$ and every fixed $\chi$-regular partition, the
conclusion of Proposition~\ref{prop:profile-transfer} holds with that
partition. More precisely, after any positive upper bound on $\rho$
is prescribed, one can choose $\rho_0$ below it, then $\varphi_0$
for each $0<\rho\le\rho_0$, and finally a common large-$N$
threshold for each $0<\varphi\le\varphi_0$. This threshold is
uniform over both profiles satisfying~\eqref{eq:scalar-integrability}
and cut arrays satisfying~\eqref{eq:almost-like}.
\end{corollary}

\begin{proof}
Take $\chi_0$ to be half the value of $\chi$ chosen in the proof of
Proposition~\ref{prop:profile-transfer}. The spatial contribution
$2T\chi$ is then at most $e/4$ for any $\chi\le\chi_0$. For the
given partition, repeat the choices of $\rho_0$, $\varphi_0$, and
$N$ in that proof. Only the number of cells enters the remaining
error terms; none involves the profile or the array.
\end{proof}

For the critical-time statement we use a consistent infinite array for
each $N$, so that finite horizons are obtained by restriction. Such
arrays admit arbitrarily fine approximation to $(\mu_t)$ if, for every
$T>0$ and every $\chi_*>0$, there are $0<\chi\le\chi_*$, a
$\chi$-regular partition satisfying~\eqref{eq:scalar-integrability},
and $g\in(0,1/k)$, $\gamma>0$, $\bar\rho>0$ such that, for every
fixed $0<\rho\le\bar\rho$ and $\varphi>0$, all sufficiently large
$N$ satisfy both~\eqref{eq:almost-like} and~\eqref{eq:local-good}
on $[0,T]$. This specifies the order of quantifiers used below.



\subsection{Shuffle matrices}
For each column $j$, choose uniformly an assignment of the symbols
$[k]_0$ to the $N$ row labels, with exactly $n_i^{(j)}$ occurrences of
symbol $i$. Choose these assignments independently across columns.
The resulting $N\times K$ matrix is uniform among all matrices with
the prescribed column counts. Write its row words as $w_1,\ldots,w_N$.
Stable sorting orders these words lexicographically and orders equal
words by their original row labels. Its permutation maps sorted
positions to original labels. Write $s_1\le_{\rm lex}\cdots\le_{\rm lex}s_N$
for the sorted words, and join $i$ to $i+1$ when $s_i=s_{i+1}$.
This defines the shuffle graph $G$.

Let $P_N^K$ denote the law of the inverse of the chronological
composition of the $K$ forward shuffles, started from the identity
deck. Thus the one-step inverse operations are applied in reverse
chronological order. Write $\nuUnif$ for uniform measure on
$\mathfrak S_N$, and
\[
 \dTV(P,Q)=\frac12\sum_{\sigma\in\mathfrak S_N}|P(\sigma)-Q(\sigma)|.
\]
For $H\subseteq\mathbb Z$, let $E(H)$ be its internal edges in the
integer path and let $\partial H$ consist of the edges with exactly
one endpoint in $H$. The latter includes edges at the ends of the
deck. For a graph $G$, $E(G)$ denotes its edge set; we also identify
$G$ with this set when taking intersections. For a permutation $\pi$,
write $\pi_G$ for the permutation obtained by sorting its values
increasingly inside every component of $G$.

\begin{lemma}[Matrix and graph representations]% Lean: formalization/lower_bound_v2/RiffleLowerV2/Model/Physical.lean :: physicalColumnShear;physicalPullbackEquiv;stableSort_physicalPullback;physicalMatrixBridge
% Lean: formalization/lower_bound_v2/RiffleLowerV2/Model/GraphLaw.lean :: matrix_multiset_expect;graphRepresentation
% Lean: formalization/lower_bound_v2/RiffleLowerV2/Model/FiniteLaws.lean :: columnIndependence;inversionTV
\label{lem:matrix-graph}
Stable sorting the fixed-count matrix has law $P_N^K$. If $G$ is
sampled from this matrix and $\pi$ is a uniform permutation independent
of the matrix, then $\pi_G$ also has law $P_N^K$. More generally,
for every real function $F$ of the sorted words and a permutation,
\[
 \E F\bigl((s_i)_{i=1}^N,\operatorname{sort}(w_1,\ldots,w_N)\bigr)
 =\E\E_\pi F\bigl((s_i)_{i=1}^N,\pi_G\bigr).
\]
The total-variation distances of the forward and inverse shuffle
laws from $\nuUnif$ are equal.
\end{lemma}
\begin{proof}
A one-step interleaving is specified by the word recording which
pile supplies each successive card. These words are uniform with
their prescribed counts. Its inverse gathers equal symbols in
increasing symbol order and preserves their previous order, which
is exactly stable sorting by that symbol.

For several steps, the symbols are initially assigned to the current
positions at each inverse operation. Process the inverse operations
from column $K$ back to column $1$, pulling each symbol assignment
back to the original row labels by the permutation already determined
by the later columns. This reindexes only that column and preserves
its counts. The reindexing permutation depends only on later columns,
so it is recovered when these later columns are known. The operation
therefore defines a bijection of the finite product of column-assignment
spaces, with inverse obtained by reversing each reindexing. It
preserves the uniform law on that product.

To identify the resulting permutation, use induction on the list of
columns. Sorting by a leading symbol, while preserving the order
already determined by the suffix, is the same as sorting by the
whole word. When two whole words agree, both procedures preserve
their original label order. The empty list gives the identity.
Consequently the pullback bijection takes the physical inverse
permutation to stable sorting of the full matrix. This proves the
first assertion, including the chronological convention.

Next relabel all rows by an independent uniform permutation. Every
fixed relabeling preserves the column counts and the uniform matrix
law, and it leaves the sorted words unchanged. Conditional on the
matrix, the relabeled original labels appear in a uniform order in
the sorted positions. Stable sorting restores increasing label
order precisely inside the equal-word components of $G$. Thus the
stable sort of the relabeled matrix is obtained by taking a uniform
permutation and sorting it inside those components. Averaging first
over the relabeling and then over the matrix proves the displayed
identity and the graph-mixture assertion.

Finally inversion is a bijection of $\mathfrak S_N$ preserving
uniform measure. Reindexing the finite sum defining total variation
by this bijection proves the last assertion.
\end{proof}


% UPPER PRELIMINARIES BEGIN
\subsection{The shuffle graph}
Use the path graph and componentwise sorting representation of
Lemma~\ref{lem:matrix-graph}. Its components are the nonempty
equal-word intervals, including isolated vertices. For a fixed path
graph $G$, denote the componentwise-sorted uniform law by $P_G$.
Thus $P_N^K=\E P_G$.

% Lean: GraphOverlap/Density.lean: graphLaw_density, overlapIdentity;
% GraphOverlap/RunProduct.lean: boundedComponentOverlap_strong, boundedComponentOverlap.
\begin{lemma}% Lean: formalization/upper_bound_v2/RiffleUpperV2/GraphOverlap/Density.lean :: graphLaw_density; overlapIdentity
% Lean: formalization/upper_bound_v2/RiffleUpperV2/GraphOverlap/RunProduct.lean :: boundedComponentOverlap_strong; boundedComponentOverlap
\label{lem:upper-overlap}
The density of $P_G$ with respect to $\nuUnif$ is
\[
 f_G(\sigma)=\mathbf1_{\{\sigma\text{ is increasing on every component of }G\}}
                   \prod_{I\in\operatorname{comp}(G)}|I|!.
\]
For any two path graphs $G,H$,
\begin{equation}\label{eq:upper-overlap}
 \E_{\nuUnif}f_G f_H
 =\frac{\displaystyle\prod_{I\in\operatorname{comp}(G)}|I|!
                    \prod_{J\in\operatorname{comp}(H)}|J|!}
          {\displaystyle\prod_{I\in\operatorname{comp}(G\cup H)}|I|!}.
\end{equation}
If the components of $H$ have at most $B$ vertices, this is at most
$B^{|G\cap H|}$. Consequently, for $B\ge2$, it is also at most
$\binom B{\lfloor B/2\rfloor}^{|G\cap H|}$.
\end{lemma}
\begin{proof}
A permutation increasing on every component has exactly
$\prod_I|I|!$ preimages under componentwise sorting; a permutation
which is not increasing has none. This proves the density formula.
The simultaneous increasing condition for $G$ and $H$ is the increasing
condition for $G\cup H$, whose uniform probability is
$1/\prod_{I\in\operatorname{comp}(G\cup H)}|I|!$. This proves
\eqref{eq:upper-overlap}.

For a vertex $i$, let $r_G(i)$ be its position, counted from the left,
in its component of $G$. Then $\prod_i r_G(i)=\prod_I|I|!$ and
$r_G(i),r_H(i)\le r_{G\cup H}(i)$. Unless the edge immediately to
the left of $i$ belongs to both graphs, at least one of $r_G(i)$ and
$r_H(i)$ is one. At the other vertices, $r_H(i)\le B$. Thus
\[
 r_G(i)r_H(i)\le r_{G\cup H}(i)
       B^{\mathbf1_{\{(i-1,i)\in G\cap H\}}}.
\]
Multiplying over the vertices proves the claimed bound. Finally
$B=\binom B1\le\binom B{\lfloor B/2\rfloor}$.
\end{proof}

% Lean: GraphOverlap/Truncation.lean: truncatedTV, truncatedVariance_eq.
\begin{lemma}% Lean: formalization/upper_bound_v2/RiffleUpperV2/GraphOverlap/Truncation.lean :: truncatedTV; truncatedVariance_eq; shuffleTruncatedRadicand_nonneg
\label{lem:upper-truncation}
Let $S$ be any event determined by one word matrix, and let $G,G'$ be
the graphs of two independent matrices with the same deterministic
column counts. Then
\begin{equation}\label{eq:upper-truncation}
 \dTV(P_N^K,\nuUnif)
 \le\frac12\sqrt{\E\left[
   \bigl(\E_{\nuUnif}f_Gf_{G'}-1\bigr)
                      \mathbf1_S\mathbf1_{S'}\right]}+\Pp(S^c).
\end{equation}
The expression under the square root is nonnegative.
\end{lemma}
\begin{proof}
Write $h=\E[f_G\mathbf1_S]$ and $r=\E[f_G\mathbf1_{S^c}]$,
so the shuffle density is $h+r$. Since $f_G$ has mean one,
$\E_{\nuUnif}h=\Pp(S)$ and $\E_{\nuUnif}r=\Pp(S^c)$.
The triangle inequality gives
\[
 \E_{\nuUnif}|h+r-1|
 \le\E_{\nuUnif}|h-\Pp(S)|+2\Pp(S^c).
\]
Apply Cauchy--Schwarz to the first term. Independence and finite
summation identify its squared upper bound as
\[
 \E_{\nuUnif}(h-\Pp(S))^2
 =\E[(\E_{\nuUnif}f_Gf_{G'}-1)\mathbf1_S\mathbf1_{S'}].
\]
This also proves the nonnegativity assertion.
\end{proof}

% Lean: GraphOverlap/FactorialTail.lean; GraphOverlap/Corollaries.lean:
% boundedComponentTail_factorial, exists_boundedComponentTail.
\begin{lemma}% Lean: formalization/upper_bound_v2/RiffleUpperV2/GraphOverlap/Corollaries.lean :: boundedComponentTail_factorial; exists_boundedComponentTail
\label{lem:upper-components}
Put $\mu_*=N\prod_{j=1}^K\max_d p_d^{(j)}$.
For every integer $B\ge1$,
\begin{equation}\label{eq:upper-components}
 \Pp(G\text{ has a component with more than }B\text{ vertices})
       \le\frac{N\mu_*^B}{(B+1)!}.
\end{equation}
In particular, if $\mu_*\le N^{-\eta}$ and $\eta B\ge A+1$,
the right side is at most $N^{-A}/(B+1)!$.
\end{lemma}
\begin{proof}
A component with more than $B$ vertices supplies $B+1$ distinct rows
having the same full word. The event is impossible if $B+1>N$,
so suppose $B+1\le N$.
For any specified set of $B+1$ rows, its probability
of agreeing in column $j$ is
\[
 \sum_d\frac{(Np_d^{(j)})_{B+1}}{(N)_{B+1}}
 \le\sum_d(p_d^{(j)})^{B+1}
 \le(\max_d p_d^{(j)})^B.
\]
Here $(m)_r=m(m-1)\cdots(m-r+1)$, and the first inequality follows
factor by factor; impossible samples have probability zero. Multiply
over columns and sum over the $\binom N{B+1}$ sets of rows.
Since $\binom N{B+1}\le N^{B+1}/(B+1)!$, this gives the result.
\end{proof}
% UPPER PRELIMINARIES END

% UPPER MAIN BEGIN
\section{Upper bound}\label{sec:upper-bound}\label{sec:upper}
\subsection{Empirical costs and an initial good shuffle}
We first give an upper bound requiring no limiting time profile.
The original deck size remains the logarithmic scale when the argument
is applied to a smaller population. Thus let $2\le n\le N$ and let
the current columns be indexed by $a<j\le K$, with $a<K\le T\log N$.
Their proportions, relative to $n$, are denoted by $p^{(j)}$. At integer
endpoints set
\begin{align*}
 A^{(N)}\left(\frac{s}{\log N},\frac{u}{\log N}\right)
       &=\frac1{\log N}\sum_{j=s+1}^u a_{p^{(j)}},\\
 B^{(N)}_\theta\left(\frac{s}{\log N},\frac{u}{\log N}\right)
       &=\frac1{\log N}\sum_{j=s+1}^u\psi_{p^{(j)}}(\theta),\\
 \beta^{(N)}\left(\frac{s}{\log N}\right)
 &=\frac{\log n}{\log N}
       -A^{(N)}\left(\frac{a}{\log N},\frac{s}{\log N}\right).
\end{align*}
We say that the empirical margin is at least $\eta$ if the following
inequalities hold for every integer pair $a\le s\le u\le K$:
\begin{equation}% Lean: formalization/upper_bound_v2/RiffleUpperV2/Shared/Empirical.lean :: empiricalA; empiricalB; empiricalBeta; empiricalExpression; empiricalInf; EmpiricalNegativeMargin
\label{eq:upper-empirical-margin}
 \begin{split}
 \beta^{(N)}\left(\frac K{\log N}\right)&\le-\eta,\\
 \inf_{3\le\theta\le4}\left\{
       \frac{\theta+3}{2}\beta^{(N)}\left(\frac s{\log N}\right)
       -B^{(N)}_\theta\left(\frac s{\log N},\frac u{\log N}\right)
       -2B^{(N)}_2\left(\frac u{\log N},\frac K{\log N}\right)\right\}
       &\le-\eta.
 \end{split}
\end{equation}
Only grid endpoints are required in this condition; the interpolation
between them is not an additional hypothesis. Zero pile proportions
are allowed. A column is $g$-good if two distinct coordinates of its
proportion vector are at least $g$.

% Lean: EmpiricalUpper/Restart.lean: terminal_horizon, restartMargin,
% restartHypotheses; Exploration/WindowClusters.lean: integerWindow_cluster.
\begin{lemma}% Lean: formalization/upper_bound_v2/RiffleUpperV2/EmpiricalUpper/Restart.lean :: restartHypotheses; terminal_horizon; restartMargin; restartTVContraction
% Lean: formalization/upper_bound_v2/RiffleUpperV2/Exploration/WindowClusters.lean :: integerWindow_cluster
\label{lem:upper-restart}
Fix $k\ge2$, $\eta_0,\eta,g,\gamma,\rho>0$, with
$4\eta\le\eta_0$, $\gamma\le1$ and
$\rho\log k\le\min\{\eta/7,1/4\}$.
For all sufficiently large $N$, an array on $N$ labels satisfying
\eqref{eq:upper-empirical-margin} with margin $\eta_0$ and the local
$(g,\gamma,\rho)$-good condition on $[0,K/\log N]$ has a suffix
with the following properties. Its first column is $g$-good; two more
$g$-good columns occur within the next $\lceil6/\gamma\rceil$ columns;
its empirical margin, after resetting its initial index to zero, is
at least $2\eta$; and every consecutive window of
$\lceil\rho\log N\rceil+1$ columns contains at least
$\gamma\rho\log N-2$ good columns. Its maximal full-word mean is at
most $N^{-2\eta}$. The total variation distance for the original array
is no larger than that for this suffix.
\end{lemma}
\begin{proof}
Since $a_p\le\log k$, the terminal margin implies
$(1+\eta_0)\log N\le K\log k$. The initial window of
$\lceil\rho\log N\rceil+1$ columns therefore fits for large $N$.
Passing from a closed time window to the integer window loses at most
two columns. Among at least $\gamma\rho\log N-2$ good indices in
that window there are three columns, say $r+1<j+1<h+1$, with
$h-r\le\lceil6/\gamma\rceil$. Indeed, otherwise every third good
index would be separated by more than that amount, contradicting
the number of good indices once $\gamma\rho\log N>8$.
Delete the $r$ columns preceding column $r+1$.

Deleting $r$ columns increases $\beta^{(N)}(s/\log N)$ at subsequent times
by at most $r\log k/\log N$, and increases every expression in
\eqref{eq:upper-empirical-margin} by at most
$7r\log k/(2\log N)$. Here $r\le\rho\log N+2$; hence the total
loss is at most $\eta/2+7\log k/\log N\le\eta$ for large $N$.
The retained margin is thus at least $2\eta$. Every retained integer
window is an original integer window, so its good-column bound is
inherited. The identity
$\mu_*=N^{\beta^{(N)}(K/\log N)}$ gives the asserted maximal mean.
Finally, the deleted shuffles act by averaging translates of the
retained permutation law. Uniform measure is invariant under these
translations, and total variation contracts under averaging. This
proves the last assertion.
\end{proof}

% Lean: Contracts/V1/Empirical.lean: E; EmpiricalUpper/Main.lean: empiricalUpper.

\begin{proof}[Proof of Theorem~\ref{thm:upper-empirical}]
Put $\eta=\min\{\eta_0/4,1\}$,
$g=\min\{g_0,1/(2k)\}$ and $\gamma=\min\{\gamma_0,1\}$.
The local condition remains true with these smaller parameters.
Choose $\rho_0$ small enough for Proposition~\ref{prop:upper-conditional}
and so that $\rho_0\log k\le\min\{\eta/7,3\eta/160,1/4\}$.
Lemma~\ref{lem:upper-restart} gives a suffix with margin $2\eta$,
the required good cluster and integer-window condition, and maximal
mean at most $N^{-2\eta}$.

Proposition~\ref{prop:upper-conditional} and the construction in
Section~\ref{sec:upper-exponential} supply $2k$ adapted explorations,
each with conditional expected remainder at most $CN^{-3\eta/80}$.
The single-matrix regularity event used there satisfies
\[
 \Pp(\mathcal R^c)\le
 AN^A e^{-c_H N^{\eta/30}}
 +2N\sqrt{AN^A e^{-c_H N^{\eta/30}}}
 +\frac{k}{2}N^{\eta/10-2\eta}.
\]
Choose an integer $B\ge2$ with $\eta B\ge2$ and intersect
$\mathcal R$ with the event that all graph components have size at
most $B$. Lemma~\ref{lem:upper-components} adds at most
$N^{-1}/(B+1)!$ to the failure probability. Apply
Proposition~\ref{prop:upper-exponential} with $d=2k$ and
$b=\binom B{\lfloor B/2\rfloor}$. All its hypotheses concern the
actual pair histories constructed below, and its denominator tends
to one. The resulting bound is a fixed multiple of
$N^{-3\eta/160}+N^{-19\eta/10}+N^{-1}$, plus stretched exponential
terms. After decreasing the positive exponent and increasing $N_0$,
this is at most $N^{-c}$. The contraction in
Lemma~\ref{lem:upper-restart} completes the proof.
\end{proof}

\subsection{A limiting profile}
% Lean: ProfileUpper/Buffered.lean: bufferedOfEmpirical;
% ProfileUpper/Main.lean: bufferedUpper; Margins/ProfileMargin.lean.

\begin{proof}[Proof of Theorem~\ref{thm:upper-profile}]
Use Corollary~\ref{cor:fixed-partition-transfer} with accuracy
$\eta/100$, first fixing $\chi$, then $\rho$, and finally $\varphi$.
The comparison is uniform over $\theta\in[2,4]$ and over interval
endpoints. Restricting to grid endpoints is exact; changing the
horizon from $T$ to $\lfloor T\log N\rfloor/\log N$ costs at most
$14\log k/\log N$. Consequently the two profile inequalities
give \eqref{eq:upper-empirical-margin} with margin $\eta/2$ for large
$N$. Choose the upper bound on $\rho$ also to satisfy
Theorem~\ref{thm:upper-empirical} with this margin and the given
$g,\gamma$. Since $K/\log N\le T$, the local condition on $[0,T]$
implies the one on $[0,K/\log N]$. The empirical theorem gives
$N^{-c}$, which is at most $\varepsilon$ for sufficiently large $N$.
All thresholds are uniform over the arrays in the statement.
\end{proof}

% Lean: Shared/Profile.lean: ArbitrarilyFineApproximation, criticalTime;
% ProfileUpper/Critical.lean: criticalOfBuffered, marginMonotonicity (Analytic.lean).

\begin{proof}[Proof of Corollary~\ref{cor:upper-critical}]
If $t$ exceeds the displayed infimum, its definition supplies $0<s<t$ at which
the reduced margin is negative. The reduced margin is nonincreasing
in the horizon, so it is negative at $t$ as well. Apply
Theorem~\ref{thm:upper-profile}, selecting the partition beneath its
mesh bound, then $\rho$ beneath both allowed bounds, and finally an
allowed positive $\varphi$. Arbitrarily fine approximation supplies
the eventual deterministic hypotheses, proving convergence.
For every fixed $t$ exceeding the infimum, convergence implies eventually
$t_{\rm mix}^{(N)}(\varepsilon)\le\lfloor t\log N\rfloor$.
Divide by $\log N$, take the upper limit, and take the infimum over
these $t$. If the infimum is $+\infty$, the upper-limit assertion
holds in the extended order.
\end{proof}

\subsection{Random pile sizes}
% Lean: Contracts/V1/Profile.lean: REmpiricalBound, REmpirical, RProfile;
% ProfileUpper/Mixture.lean, ProfileUpper/Random.lean, ProfileUpper/Main.lean: randomUpper.
\begin{corollary}% Lean: formalization/upper_bound_v2/RiffleUpperV2/ProfileUpper/Main.lean :: randomEmpiricalBound; randomEmpiricalUpper; randomProfileUpper; randomUpper
\label{cor:upper-random}
Under the parameter choices of Theorem~\ref{thm:upper-empirical}, let
$Q_N$ be any law on deterministic arrays of length $K\le T\log N$.
For all sufficiently large $N$,
\[
 \dTV\left(\int P_N^K(\mathbf n)\,Q_N(d\mathbf n),
                         \nuUnif\right)
 \le Q_N\{\text{the empirical margin or local-good condition fails}\}
                        +N^{-c}.
\]
Thus the annealed distance tends to zero if these two conditions hold
with probability tending to one, with one common deterministic $\rho$.
The same conclusion holds under the parameter choices of
Theorem~\ref{thm:upper-profile} if almost-likeness and local goodness
each hold with probability tending to one.
\end{corollary}
\begin{proof}
Convexity of total variation bounds the mixture distance by the
$Q_N$-average of the conditional distances. On the good arrays use
$N^{-c}$, and on the other arrays use the bound one. This proves the
finite inequality and its empirical consequence. In the profile
case, for any fixed positive error use
Theorem~\ref{thm:upper-profile} on the intersection of the two good
events. Its complement has probability tending to zero by the union
bound. Letting the fixed error tend to zero proves the claim.
\end{proof}
% UPPER MAIN END

% UPPER FIRST MOMENT BEGIN
\section{The first moment}\label{sec:upper-first}
We prove the uniform estimate used for the whole deck and, later, for
the random populations left inside unexplored cylinders.

% Lean: Contracts/V1/Empirical.lean: F; FirstMoment/Empirical.lean: empiricalFirstMoment;
% FirstMoment/UniformGraph.lean: empiricalGraphDecay.
\begin{proposition}% Lean: formalization/upper_bound_v2/RiffleUpperV2/FirstMoment/Empirical.lean :: empiricalFirstMoment
% Lean: formalization/upper_bound_v2/RiffleUpperV2/FirstMoment/UniformGraph.lean :: empiricalGraphDecay
% Lean: formalization/upper_bound_v2/RiffleUpperV2/Localization/EndpointWindow.lean :: endpointWindow
% Lean: formalization/upper_bound_v2/RiffleUpperV2/Localization/RelevantParents.lean :: relevant_parent_in_window; relevantParents_finite
% Lean: formalization/upper_bound_v2/RiffleUpperV2/FirstMoment/GraphAssembly.lean :: actualGraphAssembly
\label{prop:upper-first}
Fix $k\ge2$, $T,\eta,g>0$, with $g<1/k$. There is $N_0\ge2$ such
that, for every $N\ge N_0$, $2\le n\le N$, and
$a<K\le T\log N$, the following holds. Suppose that the first of
the fixed-count columns on $n$ labels is $g$-good, and that
\eqref{eq:upper-empirical-margin} holds. For independent shuffle
graphs $G,G'$ of these columns,
\begin{equation}\label{eq:upper-first}
 \E|G\cap G'|\le N^{-\eta/8}.
\end{equation}
\end{proposition}

\subsection{Interval fluctuations}
Condition on the assignment of the first column. A nonempty pile
$R$ has size $M\le(1-g)n$. The remaining columns are still uniform
fixed-count columns on all $n$ labels. The labels outside $R$ will
provide the unused population needed to smooth interval endpoints.
Permutation invariance of the rows shows that the graph law within
$R$ depends only on $M$ and its deterministic rank offset.

For a legal prefix $x$ following this first digit, write $\mu_x$ for
its mean count in $R$, $\ell_x=\log_N\mu_x$, and
$e_x=-B^{(N)}_2(M_x/\log N,K/\log N)$, where $M_x$ is its absolute final column.
Its descendants occupy an interval with endpoints $\iota_x,\tau_x$,
and its internal equal-word edges form a graph $G_x$.
All counts and fluctuations in this paragraph exclude the already
fixed first-column assignment.

Fix $0<\chi<1/k$, and in each subsequent column let $l_j,r_j$ be the
smallest and largest digits of mass at least $\chi$. For a full word,
its left boundary stops at its first digit strictly greater than $l_j$,
and its right boundary at its first digit strictly less than $r_j$;
if no crossing occurs, use the full word. Let $v$ be the longer
boundary. If either boundary has full length, call the pair terminal;
its block is the single full word. Otherwise the pair has two genuine
crossings. For every column through $v$, consider the rows whose first
difference from $v$ occurs there and is smaller, respectively larger.
Let $L_v,R_v$ be the largest means of these pieces, and put
\[
 F_v=\tfrac12\log_N\min\{L_v,R_v\},\qquad V_v=\ell_v.
\]
Both side means are positive. Their elementary comparison is
\begin{equation}\label{eq:upper-side-scale}
 \min\{L_v,R_v\}\ge\chi\mu_v,\qquad
 M\le(2|v|+\chi^{-1})\max\{L_v,R_v\}.
\end{equation}
Indeed, each genuine crossing contributes a side piece of mean at
least $\chi\mu_v$. The at most $2|v|$ first-difference pieces together
with the prefix cylinder partition $R$ in expectation, giving the
second inequality. In particular,
$2F_v\ge V_v-\log(1/\chi)/\log N$.

% Lean: Hypergeometric/AntiConcentration.lean: pointMass_le_four_div_sqrt;
% Hypergeometric/Cylinder.lean: anchorCylinderTail;
% lower_bound_v2/Thinning/Tails.lean: bernstein_of_transform, half_mean_of_transform.
\begin{lemma}% Lean: formalization/upper_bound_v2/RiffleUpperV2/Hypergeometric/AntiConcentration.lean :: HypergeomParams.pointMass_le_four_div_sqrt
% Lean: formalization/lower_bound_v2/RiffleLowerV2/Thinning/Tails.lean :: bernstein_of_transform; half_mean_of_transform
\label{lem:upper-atoms}
If $Y$ is hypergeometric with population $n$, marked count $r$, and
sample size $m$, then
\[
 \E Y=\frac{rm}{n},\qquad
 \Var(Y)=\frac{rm(n-r)(n-m)}{n^2(n-1)}\quad(n\ge2).
\]
If $\sigma^2=\Var(Y)>0$, then $\sup_z\Pp(Y=z)\le4/\sigma$.
A prefix count in any deterministic set of rows, with mean $\mu$,
satisfies, for $t>0$,
\[
 \Pp(|X-\mu|\ge t)\le2\exp\left\{-\frac{t^2}{4(\mu+t)}\right\},
 \qquad
\Pp(X\le\mu/2)\le\exp\left\{-\frac{1-\log2}{2}\mu\right\}.
\]
\end{lemma}
\begin{proof}
The moment formulas follow by expanding a sum of marked-position
indicators and using the one- and two-position sampling probabilities.
For the atom estimate let $p_q=\Pp(Y=q)$ and define, within this proof,
$b(x)=(r-x)(m-x)/n$ and
$d(x)=x(n-r-m+x)/n$. The binomial coefficient formula gives detailed
balance $b(q)p_q=d(q+1)p_{q+1}$, while
$d(x)-b(x)=x-\E Y$. Telescoping therefore yields
\[
 b(q)p_q=\sum_{j\le q}(\E Y-j)p_j
       \le\E|Y-\E Y|\le\sigma.
\]
Choose a mode $q$. Comparing its mass with the preceding mass by
detailed balance gives $q\le\E Y+2$. Since
$0\le r+m-q-\E Y\le2n$, expanding the quadratic $b$ gives
$b(q)\ge b(\E Y)-4$. Also
$b(\E Y)=\sigma^2(n-1)/n\ge\sigma^2/2$ when $n\ge2$.
If $\sigma^2\ge16$, then $b(q)\ge\sigma^2/4$, and the displayed
inequality gives $p_q\le4/\sigma$. If $\sigma^2<16$, use
$p_q\le1\le4/\sigma$.

For the prefix bound, expose successive independent uniform subsets
corresponding to its digits. The falling-factorial expansion of the
exponential transform gives the sampling bound for a nonnegative
exponent. Apply the same expansion to the complement sample for a
negative exponent, and then iterate the hypergeometric thinning.
This gives $\E e^{tX}\le\exp\{\mu(e^t-1)\}$ for every real $t$.
The centered transform and $e^t-1-t\le t^2$ for $|t|\le1$,
with $|t|$ chosen as the deviation divided by $2(\mu+\text{deviation})$,
give the two-sided bound. For the lower half-mean bound use
$t=-\log2$ in the uncentered transform. An arbitrary deterministic
row set has the same count law as a uniform row set of that size by
permutation invariance, so the argument applies to the anchor $R$.
\end{proof}

% Lean: Anchor/Smoothing.lean: endpointSmoothing, endpointShiftHypergeometric,
% smoothingFailure, endpointVariance; Anchor/ConditionalAtoms.lean.
\begin{lemma}% Lean: formalization/upper_bound_v2/RiffleUpperV2/Anchor/Smoothing.lean :: endpointSmoothing; endpointShiftHypergeometric; smoothingFailure; endpointVariance
% Lean: formalization/upper_bound_v2/RiffleUpperV2/Anchor/EqualMeans.lean :: equalCylinderMeans; smoothingFailure_equal
% Lean: formalization/upper_bound_v2/RiffleUpperV2/Anchor/EqualVariance.lean :: endpointVariance_equal
% Lean: formalization/upper_bound_v2/RiffleUpperV2/Anchor/ConditionalAtoms.lean :: conditionalEndpointAtoms
\label{lem:upper-smoothing}
For a genuine boundary pair with join $v$ and any legal continuation
$x$, there is a finite reveal containing the labeled prefix-$x$ rows
and all their continuations such that both interval endpoints are
known translates of the same hypergeometric variable. There is a
good event measurable in that reveal whose complement has probability
at most
\[
 3\exp\left\{-\frac{1-\log2}{2}g\min(L_v,R_v)\right\}.
\]
On this event the conditional variance is at least
\[
 \frac{g\min(L_v,R_v)}{16(2|v|+\chi^{-1})}.
\]
Thus every conditional endpoint atom has mass at most
$16\sqrt{(2|v|+\chi^{-1})/(g\min(L_v,R_v))}$.
\end{lemma}
\begin{proof}
Choose columns attaining $L_v$ and $R_v$. The same columns attain
these side means for every continuation $x$, since their
first-difference pieces are unchanged. It suffices to describe the
case $L_v\le R_v$; reflection interchanges the two sides, and the
complement of a hypergeometric sample has the same variance.

Suppose the left maximizing column is later than the right maximizing
column $j$. Expose every column except $j$. Let $T$ be the rows of
$R$ matching $x$ before $j$, and let $B\subseteq T$ be the rows whose
full suffix through the end of $x$ is lexicographically less than
the corresponding suffix of $x$. Expose column $j$ on
$A=T\setminus B$. Set $p=\sum_{d\le x_j}p_d^{(j)}$ and let $a$
be the number of exposed digits in $A$ at most $x_j$. The number
of such digits in $B$ is conditionally
$\operatorname{Hyp}(np-a,n-|A|,|B|)$. All other contributions to the
lower endpoint are known; the upper endpoint differs by the known
prefix count. Every prefix-$x$ row lies in $A$, so its continuation
is indeed known.

The set $B$ contains the left maximizing cylinder with column $j$
omitted; its mean is $L_v/p_{x_j}^{(j)}$. Require this count and
the remaining marked and unmarked counts to be at least half their
respective means. Lemma~\ref{lem:upper-atoms} bounds the union of
their failures by the asserted three lower-tail terms: each relevant
mean is at least $gL_v$. On this event, the marked and unmarked
fractions are at least $p/2$ and $(1-p)/2$, and the sample-complement
fraction is at least $g$, because $A\cup B\subseteq R$.
The variance is consequently at least
$g|B|p(1-p)/4\ge gL_vR_v/(8M)$.

If the left maximizing column $j$ is earlier, take $B$ to consist of
the rows with larger full suffix through the end of $x$,
and put $p=\sum_{d<x_j}p_d^{(j)}$. The same reveal leaves a
hypergeometric variable with marked count $np$ less the exposed
marked count. Now $B$ contains the right maximizing cylinder with
column $j$ omitted. The identical half-mean estimates again give
variance at least $gL_vR_v/(8M)$ and the same failure bound.

If the maximizing columns agree at $j$, expose all other columns and
the entire global set with digit $x_j$ in column $j$. Write
$p=\sum_{d<x_j}p_d^{(j)}$, $q=\sum_{d>x_j}p_d^{(j)}$, and let
$T$ be the rows of $R$ matching before $j$. On the unrevealed
$n(p+q)$ positions, exactly $np$ are marked. With $B$ the rows of
$T$ outside the exposed equal-digit set, the remaining variable is
$\operatorname{Hyp}(np,n(p+q),|B|)$.
Put $a=M\prod_{i<j}p_{x_i}^{(i)}$, so $L_v=ap$ and $R_v=aq$.
Require $|B|\ge a(p+q)/2$ and at least $(n-M)(p+q)/2$
non-equal digits outside $R$. These are two unconditional cylinder
counts with the indicated means. Each mean is at least $gL_v$,
so their failure probability is at most
$2\exp\{-(1-\log2)gL_v/2\}$. On this event the hypergeometric
mean is at least $L_v/2$, its unmarked fraction is at least $1/2$,
and its sample-complement fraction is at least $g/2$. Its variance
is at least $gL_v/8$.
In the unequal cases use \eqref{eq:upper-side-scale}; in the equal
case use $2|v|+\chi^{-1}\ge1$. The claimed common variance bound
follows. Applying the explicit atom estimate $4/\sigma$ proves the
last assertion.
\end{proof}

\subsection{Digit costs and stable intervals}
Fix $\delta>0$. Within each nonterminal boundary join, stop a branch
at its first prefix $x$ for which $\ell_x-F_v<\delta$; an immediate
stop is allowed. The terminal margin ensures that every branch stops:
at full depth, $\ell_x\le-\eta$ and
$\ell_x-F_v\le-\eta/2+\log(1/\chi)/(2\log N)<\delta$ for large $N$.
A proper stop has a parent with exponent at least $F_v+\delta$.
It also satisfies $F_v\ge\delta-\log(1/\chi)/\log N$.

For an extension from $v$ to $x$, let $S_x$ and $Q_x$ be its
normalized logarithmic costs on digits of mass at least $\chi$ and
less than $\chi$, respectively. Then
\begin{equation}\label{eq:upper-cost-split}
 \ell_x=V_v-S_x-Q_x,\qquad S_x,Q_x\ge0.
\end{equation}
Retain boundary words with at most
$(\eta+7)\log N/\log(1/\chi)$ exceptional positions before their
crossings, and retain extensions with at most
$(\log(1/\chi))^{-1/2}\log N$ exceptional positions.
Group the retained prefixes by their first pile, boundary words,
length, exceptional positions and digits, good-cost bin of width
$\eta/256$, and immediate/proper and large/small status. Here large
means $\ell_x\ge2\delta$. Terminal classes are singletons.

% Lean: FirstMoment/Count.lean, RetainedClasses.lean, Discard.lean,
% Entropy.lean, ClassExponents.lean, UniformClassDensity.lean.
\begin{lemma}% Lean: formalization/upper_bound_v2/RiffleUpperV2/FirstMoment/RetainedClasses.lean :: retainedClasses_uniform_count
% Lean: formalization/upper_bound_v2/RiffleUpperV2/FirstMoment/Discard.lean :: discardEnvelope_power_bound
% Lean: formalization/upper_bound_v2/RiffleUpperV2/FirstMoment/ClassExponents.lean :: retainedClass_originalEntropy; retainedClass_largeExponent; retainedClass_smallProperExponent; retainedClass_smallImmediateExponent; retainedClass_terminalExponent
% Lean: formalization/upper_bound_v2/RiffleUpperV2/FirstMoment/UniformClassDensity.lean :: retainedClassDensity_eventually
\label{lem:upper-classes}
Take $\delta=\eta/256$. For sufficiently small fixed $\chi>0$ and
all sufficiently large $N$, the number of retained classes is at most
$N^{\eta/64}$ and the expected number of discarded rows in each
first pile is at most $N^{-(\eta+4)}$.
For a retained class $\mathcal C$, write
$H_{\mathcal C}=\log_N|\mathcal C|$. Every nonterminal class obeys
\begin{equation}\label{eq:upper-class-entropy}
 H_{\mathcal C}\le\theta S_x-B^{(N)}_\theta(M_v/\log N,M_x/\log N)
       +\frac{\eta}{64}
       +\frac{4\log k}{\sqrt{\log(1/\chi)}}
       \qquad(3\le\theta\le4).
\end{equation}
If $p_x(m)=\Pp((m,m+1)\in G_x)$, then every retained class satisfies
\begin{equation}\label{eq:upper-class-density}
 |\mathcal C|\sum_m p_x(m)^2\le N^{-\eta/2}
                  \qquad(x\in\mathcal C).
\end{equation}
All these estimates are uniform in the original array, first-column
assignment, and retained class.
\end{lemma}
\begin{proof}
A boundary is specified by its crossing position and digit and its
exceptional positions and digits; every other digit is the fixed
left or right default. For a word of length at most $T\log N$,
the number of ways to choose at most $q$ exceptional positions and
digits is at most $\sum_{j\le q}\binom{\lfloor T\log N\rfloor}j k^j$.
The bound $\binom mj\le(em/j)^j$ shows that the logarithm of this
quantity, divided by $\log N$, tends uniformly to zero as the
allowed exceptional fraction tends to zero. Lengths and crossing
times cost only powers of $\log N$, and the number of cost bins
is bounded for fixed $\chi$. Applying this to both boundaries and
the extension gives the class bound after choosing $\chi$ small.

For one row, independence of its digits gives probability at most
$(k\chi)^j$ for exceptional digits at $j$ prescribed positions.
A union bound, followed by the same binomial estimate, bounds the
expected discarded-row count. The boundary threshold contributes
logarithmic exponent at most $1-(\eta+7)+o_\chi(1)$; the extension
threshold contributes at most
$1-\sqrt{\log(1/\chi)}+o_\chi(1)$. Their sum is at most
$N^{-(\eta+4)}$ for sufficiently small $\chi$ and large $N$.

Within a class, the nonexceptional digits determine the word.
Their product tilted law, with digit probabilities proportional to
$(p_d^{(j)})^\theta$, assigns to $x$ the mass
$N^{-\theta S_x+B^G_\theta}$, where $B^G_\theta$ sums only over
nonexceptional positions. All good costs in the class lie in one
bin. Summing these tilted masses, whose total is at most one, yields
\eqref{eq:upper-class-entropy}: the bin costs at most $4(\eta/256)$,
and each removed coordinate costs at most $4\log k/\log N$.

Every internal edge can be charged to a distinct pair of equal rows.
The fixed-count falling-factorial bound therefore gives
\[
 \sum_m p_x(m)\le\tfrac12\mu_x^2N^{e_x},\qquad
 \sum_m p_x(m)^2\le\tfrac14 N^{4\ell_x+2e_x}.
\]
For a full word the latter sum is also at most $\mu_x$.
For a nonterminal block, condition as in
Lemma~\ref{lem:upper-smoothing}. Its internal edge offsets are known
under that reveal. For each offset, putting an edge at $m$ fixes
the lower endpoint. The atom bound thus bounds the maximal edge
probability on the good event by
$16\sqrt{(2|v|+\chi^{-1})/g}\,N^{-F_v}$ times the expected internal
edge count. Multiplying by the sum of edge probabilities, and
charging the bad event at most twice at each of the $n-1$ positions,
gives
\begin{align*}
 \sum_m p_x(m)^2
 &\le4\sqrt{(2|v|+\chi^{-1})/g}\,N^{4\ell_x-F_v+2e_x}\\
 &\quad+6(n-1)\exp\left\{-\frac{1-\log2}{2}g\min(L_v,R_v)\right\}.
\end{align*}

Put $q=\log(1/\chi)/\log N$ within this proof.
For every $\theta\in[3,4]$, the identity
\[
 4\ell_x-F_v+\theta S_x
 =\theta V_v+(4-\theta)(\ell_x-F_v)
                         +(3-\theta)F_v-\theta Q_x
\]
and $2F_v\ge V_v-q$ bound the right side by
$(\theta+3)\beta^{(N)}(M_v/\log N)/2+\delta+q/2$.
Using \eqref{eq:upper-class-entropy} and then taking the infimum
in \eqref{eq:upper-empirical-margin} bounds the large-block exponent
$H_{\mathcal C}+4\ell_x-F_v+2e_x$ by
\[
 -\eta+\delta+q/2+\eta/64+4\log k/\sqrt{\log(1/\chi)}.
\]
For a proper small block the parent stopping condition instead gives
$S_x\le V_v-F_v+q\le\beta^{(N)}(M_v/\log N)/2+3q/2$.
Its nonnegativity implies $\beta^{(N)}(M_v/\log N)\ge-3q$, whence
\[
 \theta S_x\le\tfrac{\theta+3}{2}\beta^{(N)}(M_v/\log N)+21q/2.
\]
Since $4\ell_x\le8\delta$, the small-block exponent is at most
$-\eta+8\delta+21q/2+\eta/64+4\log k/\sqrt{\log(1/\chi)}$.
For an immediate small block the class is a singleton and
$\ell_x+e_x\le\beta^{(N)}(K/\log N)\le-\eta$, because
$\psi_p(2)\ge a_p$. Its exponent is at most $-2\eta+4\delta$.
A terminal block contributes at most $N^{-\eta}$.

Choose $\chi$ also so that
$4\log k/\sqrt{\log(1/\chi)}\le\eta/256$, and then take $N$
large enough that $q\le\delta$. The large and proper-small
exponents above are at most $-3\eta/4$. The square-root prefactor
is at most $N^{\eta/16}$ eventually. In the large case, stability
and $\ell_x\ge2\delta$ imply $F_v>\delta$, so its failure term,
even multiplied by the total number $(K+1)k^K$ of possible prefixes,
is at most $N^{-\eta}$. Finally absorb the sum of the two terms
using $2\le N^{3\eta/16}$. This proves
\eqref{eq:upper-class-density} in all four cases.
\end{proof}

\subsection{Estimating the first moment}
% Lean: Localization/EndpointWindow.lean; Localization/RelevantParents.lean; Localization/Eventual.lean;
% FirstMoment/ActiveCount.lean, GraphAssembly.lean: actualGraphAssembly;
% FirstMoment/UniformGraph.lean: empiricalGraphDecay.
\begin{proof}[Proof of Proposition~\ref{prop:upper-first}]
We explain the localization needed to sum the class estimates.
Put $d=K-a-1$. Express the endpoints using the smaller side of a
genuine join and its cylinder count. Every one of these at most
$d+1$ terms has mean at most $\chi^{-1}N^{2F_v}$.
Lemma~\ref{lem:upper-atoms} and the union bound show that the
probability either endpoint differs from its mean by more than $w$
is at most
\[
 2(d+1)\exp\left\{-\frac{w^2}
                {4(d+1)^2(\chi^{-1}N^{2F_v}+w)}\right\}.
\]
For a proper stop, take $w=N^{F_v+\delta/4}$ and choose $N$ so that
$\log(1/\chi)/\log N\le\delta/2$. Then $F_v\ge\delta/2$,
and the last bound is at most $\exp\{-N^{\delta/8}\}$ uniformly.
Call a block relevant at $m$ when its interval contains $m$ with
probability at least $\exp\{-N^{\delta/16}\}$. This probability
exceeds the endpoint failure bound. Since the block lies in its
parent, some realization simultaneously places $m$ in the parent
and both parent endpoints within $w$ of their means. Thus $m$ lies
in the parent's mean interval enlarged by $w$ on each side.

The parent of a proper stop has mean length at least
$N^{F_v+\delta}>2w$. At a fixed depth the mean intervals are
ordered and have disjoint interiors. At most two such enlarged
intervals can contain $m$: if three did, the entire middle interval
would fit in an interval of length $2w$. Each parent has at most
$k$ children. Summing over depths gives at most $2kd$ relevant
proper blocks for each boundary descriptor, with at most one
additional immediate or terminal block. Thus at any edge the number of relevant retained
blocks is at most
$N^{\eta/64}(2k(K-a-1)+1)\le N^{\eta/32}$.
The union bound over all possible prefixes makes the total omitted
nonrelevant marginal at most $N^{-(\eta+4)}$ for large $N$.

For completeness, the finite summation behind this reduction is as
follows. Let $q_x(m)$ be the retained-block edge marginals and let
the remaining error marginal at $m$ be the discarded-row marginal
plus the total omitted marginal. At most $L$ retained marginals are
active at one edge. Cauchy--Schwarz and $(u+v)^2\le2u^2+2v^2$ give
\[
 \E|G\cap G'|
 \le4k^2 B+4(n-1)z^2+2L\sum_x\sum_m q_x(m)^2,
\]
where $B$ bounds the expected discarded-row count in every first
pile and $z$ bounds the omitted marginal. The factor $k^2$ follows
by summing the pile errors before squaring. This calculation uses
$\E|G\cap G'|=\sum_m\Pp(m\in G)^2$, which follows from independence.

Take $B=z=N^{-(\eta+4)}$ and $L\le N^{\eta/32}$.
Equation~\eqref{eq:upper-class-density}, summed over the at most
$N^{\eta/64}$ classes, bounds the double sum by
$N^{\eta/64-\eta/2}$. Hence, for sufficiently large $N$,
\[
 \E|G\cap G'|\le
 N^{-\eta}+2N^{3\eta/64-\eta/2}\le N^{-\eta/8}.
\]
All choices were made before $n,a,K$ and the column counts. Averaging
over the first-column assignments therefore preserves the estimate.
\end{proof}
% UPPER FIRST MOMENT END

% UPPER EXPONENTIAL BEGIN
\section{The exponential moment}\label{sec:upper-exponential}
\subsection{Preparatory estimates}
Fix $k\ge2$, $T>0$, $0<\eta\le1$, $0<g<1/k$ and
$0<\gamma\le1$. We consider a retained array of length at most
$T\log N$, with empirical margin at least $2\eta$. Its first column
is $g$-good, two further good columns occur within the next
$D=\lceil6/\gamma\rceil$ columns, and every consecutive window of
$\lceil\rho\log N\rceil+1$ columns contains at least
$\gamma\rho\log N-2$ good columns. All randomness below is in the
uniform assignments of these fixed column counts.

The short tail prefix has length $\min\{D,K-1\}$. The two good
columns in this prefix provide three
ordered short words $v_-<v_0<v_+$, each of tail mass at least
$g^2k^{-D}$. To see this, choose two ordered digits of mass at least
$g$ in each of these columns and a maximal-mass digit in every other
coordinate. Three of the four resulting words are strictly ordered
and have the required masses. Choose the lexicographically first
admissible triple; it depends only on the counts. Assign a full word
to the forward direction when its short tail prefix is at most $v_0$,
and to the backward direction when it is at least $v_0$.

In a first-column pile $p$ of size $n_p$, reveal the sorted words
from the left for the forward direction and from the right for the
backward direction. The natural history $h_i$ consists precisely of
the first $i$ words so revealed. The chronological order of their
coordinates is unchanged in the backward exploration. We never add
a global concentration event to a natural history.

Let $H$ be the following event on one matrix. For every pile and
later column, each digit count outside the pile differs from its
mean by at most $N^{2/3}$. For every prefix $x$ with
$\mu_x\ge N^{\eta/10}$, its count differs from $\mu_x$ by at most
$\mu_x^{2/3}$, and each later digit count within that prefix differs
from its count times the original digit proportion by at most
$\mu_x^{2/3}$. Finally, in every pile of size at least $N^{\eta/10}$,
the numbers of words strictly before and at or before $v_0$ differ
from their means by at most $g^2k^{-D}n_p/4$.

% Lean: Exploration/UniformConcentration.lean: uniformConcentration;
% Exploration/HistoryProbability.lean: historyGoodProbability;
% Exploration/Regularity.lean: historyGood_unionBound.
\begin{lemma}% Lean: formalization/upper_bound_v2/RiffleUpperV2/Exploration/UniformConcentration.lean :: uniformConcentration
% Lean: formalization/upper_bound_v2/RiffleUpperV2/Exploration/AllCylinders.lean :: allCylinderFailure; cylinder_half_concentration
% Lean: formalization/upper_bound_v2/RiffleUpperV2/Exploration/GuardConcentration.lean :: allCylinder_implies_concentration; allCylinder_guard_deviation
% Lean: formalization/upper_bound_v2/RiffleUpperV2/Exploration/Regularity.lean :: historyGood_unionBound
\label{lem:upper-history-good}
There exist $A,c_H>0$ and $N_0$, depending only on the fixed
parameters, such that for $N\ge N_0$,
\begin{equation}\label{eq:upper-H}
 \Pp(H^c)\le AN^A\exp\{-c_H N^{\eta/30}\}.
\end{equation}
For any fixed pile, direction and reveal time, the probability of
histories $h_i$ for which
\begin{equation}\label{eq:upper-history-good}
 \Pp(H^c\mid h_i)\le
          \sqrt{AN^A\exp\{-c_H N^{\eta/30}\}}
\end{equation}
fails is at most the right side of \eqref{eq:upper-history-good}.
The probability that it fails at any nonterminal positive reveal
time in any pile or direction is at most $2N$ times that bound.
\end{lemma}
\begin{proof}
Use the stronger simultaneous event obtained by prescribing an
arbitrary set of allowed digits in every column. If such a cylinder
has mean $\mu$, require its count to lie within
$\max\{\mu,N^{\eta/10}\}^{2/3}/2$ of its mean.
The transform proof of Lemma~\ref{lem:upper-atoms} applies equally
to allowed-digit sets. Substituting this tolerance in its tail
bound shows that each test fails with probability at most
$2\exp\{-N^{\eta/30}/24\}$.
There are exactly $(2^k)^K$ tests. Their simultaneous failure
probability is therefore at most
$2(2^k)^K\exp\{-N^{\eta/30}/24\}$.

This stronger event implies $H$ deterministically. Outside-pile
digit counts are allowed-digit cylinders of mean at most $N$.
For a prefix of mean $\mu_x\ge N^{\eta/10}$, both its count and
the count obtained by prescribing one later digit have errors at
most $\mu_x^{2/3}/2$. Subtract the original digit proportion times
the first error from the second; their sum is at most
$\mu_x^{2/3}$. Each guard is a sum of at most $k^D$ short-cylinder
counts. In a pile with $n_p\ge N^{\eta/10}$, all their means are
at most $n_p$, so its error is at most $k^D n_p^{2/3}/2$.
For large $N$,
$2k^D/(g^2k^{-D})\le N^{\eta/30}\le n_p^{1/3}$,
which makes this at most the required $g^2k^{-D}n_p/4$.
Finally $(2^k)^K\le N^{kT\log2}$. This proves
\eqref{eq:upper-H}, with the permissible explicit choices
$A=2+kT\log2$ and $c_H=1/24$.

The expectation of $\Pp(H^c\mid h_i)$ is $\Pp(H^c)$, by summing
over the finite history fibers. Markov's inequality proves the
single-history assertion. There are $2\sum_p n_p=2N$ pile-direction
reveal indices, so the final assertion follows by the union bound.
\end{proof}

Choose $0<\tau<g$ and write
$\alpha=\log(1/g)/\log(1/\tau)$. We take $\tau$ so small that
\begin{equation}\label{eq:upper-alpha}
 \alpha<\tfrac12,\qquad
 \alpha\le\frac{\eta/10}{4(1+\eta/10)},\qquad
 \alpha\le\frac{\eta}{40(1+\eta)}.
\end{equation}
This choice precedes $\rho,N$, the array, and the history. Denote by
$\lambda_x$ the product probability of a prefix $x$, so
$\mu_x=N\lambda_x$.

% Lean: Reservoir/PrefixCount.lean: prefixCount; Reservoir/ExtraCopies.lean: extraCurrentCopies;
% Reservoir/TailCollisions.lean: tailCollisions; Reservoir/BackwardCounts.lean.
\begin{lemma}% Lean: formalization/upper_bound_v2/RiffleUpperV2/Reservoir/Denominator.lean :: history_denominator_lower_uniform
% Lean: formalization/upper_bound_v2/RiffleUpperV2/Reservoir/ProductUpper.lean :: product_power_envelope; coordinate_power_envelope; eventual_logarithmic_error
% Lean: formalization/upper_bound_v2/RiffleUpperV2/Reservoir/PrefixCount.lean :: prefixCount
% Lean: formalization/upper_bound_v2/RiffleUpperV2/Reservoir/ExtraCopies.lean :: extraCurrentCopies
% Lean: formalization/upper_bound_v2/RiffleUpperV2/Reservoir/TailCollisions.lean :: tailCollisions
\label{lem:upper-reservoir}
There exist $C>0$ and $N_0$, depending only on the fixed parameters,
with the following property for $N\ge N_0$.
Fix a positive-probability natural history ending at word $s$ in its
assigned direction, and suppose $\Pp(H^c\mid h_i)\le\varepsilon$,
where $\varepsilon\ge0$. Every prefix cylinder $x$ strictly beyond $s$ in this
direction and contained in its first pile satisfies
\begin{align}
 \E[C_x\mid h_i]&\le CN\lambda_x^{1-\alpha}+N\varepsilon,
                                           \label{eq:upper-reservoir-count}\\
 \E\left[\sum_{w\succeq x}\binom{C_w}{2}\,\middle|\,h_i\right]
 &\le CN^2\lambda_x^{2(1-\alpha)}
       \prod_{j>|x|}\sum_d(p_d^{(j)})^{2(1-\alpha)}
                  +N^2\varepsilon.          \label{eq:upper-reservoir-pairs}
\end{align}
If $t$ copies of the current word $s$ have been observed, then also
\begin{equation}\label{eq:upper-extra-copies}
 \E[(C_s-t)_+\mid h_i]
       \le CN\lambda_s^{1-\alpha}+N\varepsilon.
\end{equation}
\end{lemma}
\begin{proof}
We give the forward proof. Mark one or two labeled rows of the
initial pile and expose all other pile rows. Keep all outside-pile
rows unexposed. In each later column, if $m\in\{1,2\}$ rows are
marked, the remaining population has size $N-n_p+m$. Conditional
on the exposed values, its assignments are uniform with the exact
remaining counts and independent between columns. This follows
by counting completions of each column; the number of completions
is the corresponding multinomial coefficient.

Deleting marked rows which lie strictly after the observed prefix
leaves that prefix valid. Completing the marked rows inside the
right reserve $v_+$ therefore preserves the history. On every
exposed fiber meeting $H$, its residual reserve counts differ from
their means by at most $N^{2/3}+m$. Since $N-n_p\ge gN$ and
the reserve depth is fixed, the product of the relevant falling-
factorial ratios is at least $(g^2k^{-D})^m/2$ for large $N$.
This gives a positive, uniform lower bound for the history denominator
on such a fiber.

For the numerator, a prescribed digit of mass $p\ge\tau$ has
residual proportion at most $p(1+CN^{-1/3})$, by the outside-pile
concentration bound. A digit of mass $p<\tau$ has residual
proportion at most $p/g$, because the remaining population has
size at least $gN$. The identity $\tau^\alpha=g$ gives
$p/g\le p^{1-\alpha}$ in the latter case; in the former,
$p\le p^{1-\alpha}$. For $m=1,2$, the exact sampling factor is
no larger than the $m$th power of that residual proportion.
Multiplying over at most $T\log N$ columns costs at most $2^m$
for large $N$. The numerator is consequently bounded by
$2^m$ times the prescribed product mass to power $m(1-\alpha)$.

Divide by the reserve lower bound and then average the exposed
information back down to the natural history. Summing the one-row
bound over labeled rows gives \eqref{eq:upper-reservoir-count}.
Summing the two-row bound over pairs and their common full tail
gives \eqref{eq:upper-reservoir-pairs}; the full-tail sum factorizes
into the displayed product. The bad part is charged once, before
this summation: a count is at most $N$ and the total number of
pairs at most $N^2$, so its costs are $N\varepsilon$ and
$N^2\varepsilon$, respectively. In particular there is no error
factor counting terminal words.

For the current word, mark only a removable copy: after deleting it,
the $t$ observed copies must remain. The number of removable copies
dominates $(C_s-t)_+$. Complete it in the right reserve as before,
and apply the same one-row calculation. This proves
\eqref{eq:upper-extra-copies}, including the possible edge from the
last observed copy to the first unobserved copy. For the backward
direction use the left reserve and the reversed spatial inequalities.
The columns themselves retain their chronological order.
\end{proof}

% Lean: Reservoir/ConditionalFutureCylinderLaw.lean: conditionalFutureCylinderLaw,
% conditionalFutureCylinderLaw_unlabeled; Conditional/Jensen.lean;
% Conditional/GraphSquares.lean: fineBlockSquares_le_firstMoment.
\begin{lemma}% Lean: formalization/upper_bound_v2/RiffleUpperV2/Reservoir/ConditionalFutureCylinderLaw.lean :: conditionalFutureCylinderLaw; conditionalFutureCylinderLaw_unlabeled
\label{lem:upper-induced-law}
Let a finite family of disjoint prefix cylinders lie strictly beyond
the last word of a natural history. Refine the history by revealing
their labeled row sets, their rank intervals, and their counts of
each digit in each later column, together with the matrix outside
their interiors. Conditional on this information, the internal
graph of any cylinder has the law obtained from independent uniform
fixed-count tail columns with its induced counts, translated by its
known lower rank. The same statement holds after averaging out the
labeled row sets while retaining the induced counts and intervals.
\end{lemma}
\begin{proof}
Inside one specified cylinder all its prefix coordinates are fixed.
The natural history concerns words strictly outside it, so it places
no further order condition on its interior. After the outside matrix
and inside column counts are fixed, every compatible assignment of
one tail column inside has the same number of full completions.
Counting independently in each column therefore gives the uniform
product of its fixed-count column laws. Sorting inside the cylinder
sorts precisely these tail words; its lower rank supplies the
deterministic translation. Permuting the labels within the fixed
row set does not change that graph law. Hence every labeled atom
with the same counts and interval has the same conditional graph
law, and averaging these atoms proves the unlabeled assertion.
\end{proof}

\subsection{The exploration}
For independent matrices, let $X_i^{p,\pm}$ count the common
directional edges in pile $p$ with both endpoints among its first
$i$ revealed vertices. Use the identically zero process if
$n_p<N^{\eta/10}$. Each process starts at zero and has increments
zero or one. Its pair history consists of the two natural histories
in the same pile and direction. The remaining count is
$X_{n_p}^{p,\pm}-X_i^{p,\pm}$; at time $i$ it includes the edge
crossing from the last revealed vertex to the next vertex.
A pair history is good if each of its two histories satisfies
\eqref{eq:upper-history-good}; declare time zero and terminal time
good without further restrictions.

% Lean: Contracts/V1/Conditional.lean: C; Conditional/Main.lean: remainingCommonEdges;
% Conditional/OneHistory.lean: oneHistory_uniform; Conditional/Parameters.lean.
\begin{proposition}% Lean: formalization/upper_bound_v2/RiffleUpperV2/Conditional/Main.lean :: remainingCommonEdges
% Lean: formalization/upper_bound_v2/RiffleUpperV2/Conditional/OneHistory.lean :: oneHistory_uniform
% Lean: formalization/upper_bound_v2/RiffleUpperV2/Conditional/RegimeAveraging.lean :: regimeBlockAverage_uniform
% Lean: formalization/upper_bound_v2/RiffleUpperV2/Conditional/PreparedFamily.lean :: preparedFamily_exists
\label{prop:upper-conditional}
Under the fixed parameters and retained-array hypotheses of this
section, there is $\rho_0>0$ such that every fixed
$0<\rho\le\rho_0$ admits $A,c_H,C>0$ and $N_0\ge2$ for which
\eqref{eq:upper-H} holds and, on every positive-probability good
pair-history fiber in a pile of size at least $N^{\eta/10}$,
\begin{equation}\label{eq:upper-conditional}
 \E[X_{n_p}^{p,\pm}-X_i^{p,\pm}\mid h_i,h'_i]
                          \le CN^{-3\eta/80}
                         \qquad(0\le i\le n_p).
\end{equation}
At $i=n_p$ the remaining count is identically zero. All constants
are chosen before the array, pile, direction, and history.
\end{proposition}
\begin{proof}
Choose $\tau$ according to \eqref{eq:upper-alpha}, and take
$\rho\log k\le3\eta/160$. Fix one positive history with current
word $s$. If $s$ is outside its assigned direction, no directional
edge remains. Otherwise partition the strict future in that pile
into prefix cylinders. Following the two boundary paths of the
lexicographic interval gives at most $2kK$ disjoint cylinders.
If at least $\lceil\rho\log N\rceil+1$ columns remain after a
prefix, refine it to just before a good column in the next such
window. The integer-window condition and the three-index argument
used in Lemma~\ref{lem:upper-restart} provide that column. Each
refinement adds at most $\lceil\rho\log N\rceil+1$ coordinates.
The resulting disjoint family $\mathcal B$ therefore has size at
most
\begin{equation}\label{eq:upper-future-count}
 2k^3T(\log N)N^{\rho\log k}.
\end{equation}
Each member either has a good next column or has fewer than that
many remaining columns. In the latter case its maximal tail mass
is at least $k^{-\lceil\rho\log N\rceil-1}$, whereas its mean
times that tail mass is at most $N^{-\eta}$. Its mean is thus
less than $N^{-\eta/10}$ for large $N$.

Let $q_e(h_i)$ be the conditional probability of a remaining
directional edge at position $e$. Refine the history only for this
estimate, using the family information in
Lemma~\ref{lem:upper-induced-law}. Conditional Jensen gives
\[
 \sum_e q_e(h_i)^2
 \le\E\left[\sum_e q_e(h_i,Z)^2\,\middle|\,h_i\right].
\]
The refined edge slots split into the remaining current-word run
and the disjoint cylinder interiors. An equal-word edge cannot
cross between two distinct cylinders. The current-run contribution
is bounded by \eqref{eq:upper-extra-copies}. For the final cylinders,
after refinement, use three regimes.

If $\mu_x\le N^{-\eta/10}$, the sum of squared edge probabilities
is at most the expected vertex count. By
\eqref{eq:upper-reservoir-count} it is at most
\[
 CN^{-\eta/10+\alpha(1+\eta/10)}+N\varepsilon
                   \le CN^{-3\eta/40}+N\varepsilon.
\]
If $N^{-\eta/10}<\mu_x\le N^{\eta/10}$, it is at most the
expected number of equal-word pairs. Write
$q_* =\prod_{j>|x|}\max_d p_d^{(j)}$.
The terminal margin gives $\mu_xq_*\le N^{-\eta}$, hence
$q_*\le N^{-\eta+\eta/10}$. Since $2(1-\alpha)>1$,
\[
 \prod_{j>|x|}\sum_d(p_d^{(j)})^{2(1-\alpha)}
                         \le q_*^{1-2\alpha}.
\]
Equation~\eqref{eq:upper-reservoir-pairs} bounds this regime by
\[
 CN^{-\eta+3\eta/10+2\alpha(1+\eta-2\eta/10)}+N^2\varepsilon
                         \le CN^{-3\eta/5}+N^2\varepsilon.
\]

Finally, suppose $\mu_x>N^{\eta/10}$. Such a cylinder has a good
next column. On $H$, its size $m$ and later digit counts satisfy
\[
 \tfrac12\mu_x\le m\le2\mu_x,\qquad
 \max_{j,d}|p_{d,\mathrm{ind}}^{(j)}-p_d^{(j)}|
                            \le2N^{-\eta/30}.
\]
For two simplex vectors at sup-distance at most $z$, the scalar
cost bounds are
$|a_p-a_q|\le kz$ and
$|\psi_p(\theta)-\psi_q(\theta)|\le4k^4z$ for $2\le\theta\le4$.
Indeed, their maxima are at least $1/k$, their power sums are at
least $k^{-3}$, the logarithm is Lipschitz above these lower bounds,
and $|x^\theta-y^\theta|\le4|x-y|$ on $[0,1]$.
Thus replacing the suffix population by the induced one loses at
most
\[
 \frac72\left(\frac{\log2}{\log N}+2kT N^{-\eta/30}\right)
                         +24k^4T N^{-\eta/30}
\]
in its empirical expression. The maximal-prefix inequality
$\mu_x\le N\prod_{j\le|x|}\max_d p_d^{(j)}$ supplies the required
bound on its initial population exponent. The original margin
$2\eta$ therefore leaves margin $\eta$ in the induced array, and
its first column is $g/2$-good, for large $N$.
Lemma~\ref{lem:upper-induced-law} and
Proposition~\ref{prop:upper-first} bound its conditional
self-intersection by $N^{-\eta/8}$ on these good family atoms.
Their complement has conditional probability at most $\varepsilon$,
since the natural history is good and $H$ implies the needed
conditions. Its contribution is at most $N\varepsilon$.

The current word has mean at most $N^{-\eta}$, so
\eqref{eq:upper-extra-copies} bounds its excess by
$CN^{-\eta+\alpha(1+\eta)}+N\varepsilon$.
Combining the regimes, using \eqref{eq:upper-future-count}, and
taking $\varepsilon=\sqrt{AN^A e^{-c_HN^{\eta/30}}}$, gives
\[
 \sum_e q_e(h_i)^2\le
 CN^{-3\eta/40}
 +C(\log N)N^{\rho\log k}
       \bigl(N^{-3\eta/40}+N^{-3\eta/5}+N^{-\eta/8}
                                      +N^2\varepsilon\bigr)
 \le CN^{-3\eta/80}.
\]
The last step uses $\rho\log k\le3\eta/160$ and
$\log N\le N^{3\eta/160}$ eventually; the stretched exponential
absorbs every polynomial factor.

The two matrices remain independent when their separate histories
are fixed. Thus the expected remaining common count is
$\sum_e q_e(h_i)q_e(h'_i)$, bounded by half the sum of their two
self-intersections. This proves \eqref{eq:upper-conditional} at
positive nonterminal times. At time zero the remaining directional
count is bounded by the total common-edge count, so
Proposition~\ref{prop:upper-first} applies directly. The terminal
count is zero by definition. The backward proof uses the same
argument with the spatial order reversed.
\end{proof}

\subsection{Proof of the exponential estimate}
% Lean: Exponential/Stopping.lean; Exponential/Moments.lean;
% Exponential/Excess.lean: adaptedExcess.
\begin{lemma}% Lean: formalization/upper_bound_v2/RiffleUpperV2/Exponential/Excess.lean :: adaptedExcess
% Lean: formalization/upper_bound_v2/RiffleUpperV2/Exponential/Moments.lean :: stoppedTail; binomialMoment; stoppedBinomialMoment
\label{lem:upper-adapted}
On a finite probability space let $(\mathcal F_i)_{0\le i\le m}$
be an increasing filtration, let $X_0=0$ be an adapted process with
increments in $\{0,1\}$, and let $B_i\in\mathcal F_i$.
Suppose that for every $i<m$, on every positive-probability
$\mathcal F_i$-atom contained in $B_i$,
\[
 \E[X_m-X_i\mid\mathcal F_i]\le u,
 \qquad u\ge0.
\]
For $t\ge0$ with $(e^t-1)u<1$ and any event
$C\subseteq\bigcap_{i=0}^m B_i$,
\begin{equation}\label{eq:upper-adapted}
 \E[(e^{tX_m}-1)\mathbf1_C]
                  \le\frac{(e^t-1)u}{1-(e^t-1)u}.
\end{equation}
\end{lemma}
\begin{proof}
Set $D_i=\bigcap_{h=0}^iB_h$ and
$Z_i=\sum_{j=1}^i(X_j-X_{j-1})\mathbf1_{D_j}$.
Then $Z$ is adapted, has unit increments, and
$Z_m-Z_i\le\mathbf1_{D_i}(X_m-X_i)$. Therefore its conditional
expected remainder is at most $u$ on every positive-probability
history atom, without a good-event restriction.

For every integer $r\ge0$ we claim
$\E[\binom{Z_m-Z_i}{r}\mid\mathcal F_i]\le u^r$.
The claim is immediate for $r=0$. For $r\ge1$, the pathwise identity
\[
 \binom{Z_m-Z_i}{r}=
 \sum_{j=i+1}^m(Z_j-Z_{j-1})\binom{Z_m-Z_j}{r-1}
\]
counts an $r$-tuple of future successes by its first success.
Condition each summand on $\mathcal F_j$, apply the induction
hypothesis to the second factor, then condition on $\mathcal F_i$.
The conditional first-moment bound proves the claim.
The binomial theorem now gives
\[
 \E(e^{tZ_m}-1)
 \le\sum_{r\ge1}((e^t-1)u)^r
       =\frac{(e^t-1)u}{1-(e^t-1)u}.
\]
On $C$ no increment was discarded, so $Z_m=X_m$; outside $C$
the exponential excess is nonnegative. This proves the result.
\end{proof}

% Lean: Exponential/FiniteTV.lean: finiteExplorationTV_of_majorization,
% finiteExplorationTV; Exponential/Aggregation.lean: finiteExplorationExcess.
\begin{proposition}% Lean: formalization/upper_bound_v2/RiffleUpperV2/Exponential/FiniteTV.lean :: finiteExplorationTV; finiteExplorationTV_of_majorization
\label{prop:upper-exponential}
Let $S$ be an event on one matrix whose graph components have size
at most $B\ge2$. Suppose $d\ge1$ separate explorations on the pair
of matrices satisfy Lemma~\ref{lem:upper-adapted} with common bound
$u$, and that on $S\times S$ all their good events hold and
$|G\cap G'|\le\sum_{a=1}^d X_{m_a}^{(a)}$.
Put $b=\binom B{\lfloor B/2\rfloor}$. If $(b^d-1)u<1$, then
\[
 \dTV(P_N^K,\nuUnif)
 \le\frac12\sqrt{\frac{(b^d-1)u}{1-(b^d-1)u}}+\Pp(S^c).
\]
\end{proposition}
\begin{proof}
For nonnegative $x_1,\ldots,x_d$, convexity of the exponential gives
$b^{\sum_a x_a}-1\le d^{-1}\sum_a(b^{d x_a}-1)$.
Apply this on $S\times S$ and then apply
Lemma~\ref{lem:upper-adapted} separately to each exploration with
$t=d\log b$ and $C=S\times S$. It bounds the expected truncated
excess $b^{|G\cap G'|}-1$ by $(b^d-1)u/(1-(b^d-1)u)$.
Lemmas~\ref{lem:upper-overlap} and \ref{lem:upper-truncation}
then give the asserted total variation bound. No common filtration
for the separate explorations is required.
\end{proof}

% Lean: Exploration/UniformRegularity.lean, DirectionalCover.lean, SmallPiles.lean;
% EmpiricalUpper/ExplorationTV.lean: concreteExplorationTV.
We finish by checking the truncation used in the proof of
Theorem~\ref{thm:upper-empirical}. Let $\mathcal R$ be the event
that $H$ holds, every positive nonterminal history in every pile
and direction satisfies \eqref{eq:upper-history-good}, and every
pile smaller than $N^{\eta/10}$ is edge-free. The expected number
of equal-word pairs in such a pile is at most
$n_p\mu_*/2$. Summing over at most $k$ small piles gives
\[
 \Pp(\mathcal R^c)\le AN^A e^{-c_HN^{\eta/30}}
      +2N\sqrt{AN^A e^{-c_HN^{\eta/30}}}
      +\frac{k}{2}N^{\eta/10-2\eta}.
\]
The guard concentration ensures the directional cover for two
matrices: their middle-cylinder rank intervals overlap, because
their common expected gap is at least $g^2k^{-D}n_p$ whereas the
two endpoint errors total at most $g^2k^{-D}n_p/2$.
Thus an edge common to both matrices is assigned to a common
direction in its pile. Small piles contribute no edges on
$\mathcal R$. Hence
$|G\cap G'|\le\sum_{p,\pm}X_{n_p}^{p,\pm}$ on
$\mathcal R\times\mathcal R$.
The natural pair histories refine with time, the edge processes
are adapted and have unit increments, and their good events are
history measurable. Proposition~\ref{prop:upper-conditional}
provides every conditional remainder bound, including time zero.
Intersecting $\mathcal R$ with the bounded-component event therefore
verifies all hypotheses of Proposition~\ref{prop:upper-exponential}.
This completes the upper-bound argument.
% UPPER EXPONENTIAL END



\section{Lower bound}\label{sec:lower-bound}\label{sec:lower}
Throughout this section $k\ge2$, $T>0$, and $K=\lfloor T\log N\rfloor$, unless a statement allows arbitrary $K$. As above, $P_N^K$ denotes the law of the inverse shuffled permutation and $\nuUnif$ the uniform law on $\mathfrak S_N$. Inversion preserves total variation from uniform, so the same conclusions apply to the chronological forward shuffle. We construct events whose probabilities tend to one under $P_N^K$ and to zero under $\nuUnif$. Every assertion of uniform convergence below means that, after the displayed fixed parameters and an accuracy $\varepsilon>0$ are chosen, one threshold $N_0$ works for every admissible cut array and deterministic family in the statement.
\subsection{Increasing sequences}\subsubsection{Prefix counts}

For a word \(w=(w_1,\ldots,w_m)\) of length \(m\le K\), let
\[
 \lambda_w=\prod_{j=1}^m p^{(j)}_{w_j},\qquad
 t_w=\sum_{\substack{v\in[k]_0^m\\v<_{\rm lex}w}}\lambda_v,
 \qquad C(w)=\#\{i:s_i|_m=w\}.
\]
Here \(s_i|_m\) is the length-\(m\) prefix of the sorted row \(s_i\).
Define the zero-based cuts
\[
 L_w=\#\{i:s_i|_m<_{\rm lex}w\},\qquad
 U_w=\#\{i:s_i|_m\le_{\rm lex}w\}.
\]
Thus the actual ranks occupied by this prefix are
\(I(B_w)=(L_w,U_w]\cap\mathbb Z\), and \(U_w-L_w=C(w)\).
Its deterministic probability interval is
\(J_w=[t_w,t_w+\lambda_w)\). These definitions also apply to the empty
word, whose mass is one and whose count is \(N\).

% Lean source: formalization/lower_bound_v2/RiffleLowerV2/Thinning/Process.lean :: RiffleLowerV2.Thinning.successiveThinning; RiffleLowerV2.Thinning.thinningTails
% Lean source: formalization/lower_bound_v2/RiffleLowerV2/Thinning/Compat/HypergeometricConcentration.lean :: RiffleLowerV2.Thinning.Compat.HypergeomParams.expect_pow_le_all
% Lean source: formalization/lower_bound_v2/RiffleLowerV2/Thinning/Main.lean :: RiffleLowerV2.Thinning.residualPrefixTransform; RiffleLowerV2.Thinning.residualPrefixTails
\begin{lemma}[Successive thinning transforms and tails]
% Lean: formalization/lower_bound_v2/RiffleLowerV2/Thinning/Process.lean :: RiffleLowerV2.Thinning.successiveThinning; RiffleLowerV2.Thinning.thinningTails
% Lean: formalization/lower_bound_v2/RiffleLowerV2/Thinning/Compat/HypergeometricConcentration.lean :: RiffleLowerV2.Thinning.Compat.HypergeomParams.expect_pow_le_all
% Lean: formalization/lower_bound_v2/RiffleLowerV2/Thinning/Main.lean :: RiffleLowerV2.Thinning.residualPrefixTransform; RiffleLowerV2.Thinning.residualPrefixTails
\label{lem:successive-thinning-tail}
Let \(M\ge1\) and let \(P_1,\ldots,P_m\) be independent uniform
subsets of a fixed set of \(M\) labels, with deterministic sizes
\(m_1,\ldots,m_m\). Put
\[
 X=|P_1\cap\cdots\cap P_m|,\qquad
 q_j=m_j/M,\qquad Q=\prod_jq_j,\qquad \mu=MQ.
\]
For an empty intersection use \(X=M\) and \(Q=1\). For every \(s>0\),
\begin{equation}
\label{eq:thinning-transform}
 \E s^X\le(1-Q+Qs)^M\le \exp(\mu(s-1)).
\end{equation}
Consequently, for \(t>0\),
\begin{equation}
\label{eq:thinning-bernstein}
 \Pp(|X-\mu|\ge t)\le
 2\exp\left(-\frac{t^2}{4(\mu+t)}\right),
\end{equation}
and, since $(1-\log2)/2>0$,
\begin{equation}
\label{eq:thinning-lower-tail}
 \Pp(X\le\mu/2)\le e^{-\frac{1-\log2}{2}\mu},
 \qquad \Pp(X<2)\le2e^{-\mu/2}.
\end{equation}
These statements apply to the count of any prescribed word in the shuffle
matrix. They also hold conditionally after any admissible fixed prescriptions
on \(b<N\) labelled rows: delete those rows, replace \(M\) by \(N-b\),
and subtract the prescribed occurrences of each digit from its column count.
\end{lemma}

\begin{proof}
First let $V$ be the intersection of a fixed $x$-element set with a
uniform $m$-element subset of an $M$-element population. Counting prescribed
marked labels gives
\[
 \E\binom Vr=\binom xr\frac{(m)_r}{(M)_r}.
\]
For $s=1+u\ge1$, expand $(1+u)^V$ and use
$(m)_r/(M)_r\le(m/M)^r$. All terms are nonnegative, and hence
\begin{equation}\label{eq:one-step-thinning-transform}
 \E s^V\le(1-m/M+(m/M)s)^x.
\end{equation}
For $0<s\le1$, interchange the roles of the marked set and the sample.
The complement count $x-V$ is hypergeometric with sample size $x$ and
$M-m$ marked positions. Apply the preceding bound to $s^{-1}\ge1$ and
multiply by $s^x$. This gives exactly the right side of
\eqref{eq:one-step-thinning-transform}. The cases $x=0$ and $m=0$ are
included. Here $(n)_r=n(n-1)\cdots(n-r+1)$, with $(n)_0=1$.
Let \(X_j=|P_1\cap\cdots\cap P_j|\), with \(X_0=M\).
Conditionally on the first \(j-1\) subsets,
\(X_j\sim\operatorname{Hyp}(M,X_{j-1},m_j)\).
Starting from \(z_m=s\), set
\(z_{j-1}=1-q_j+q_jz_j\). Every \(z_j\) is positive, and
\(z_0=1-Q+Qs\). Iterating \eqref{eq:one-step-thinning-transform}
gives \(\E s^X\le z_0^M\). The inequality \(1+v\le e^v\)
gives the second bound in \eqref{eq:thinning-transform}.

For \(0\le h\le1\), both
\(e^h-1-h\le h^2\) and \(e^{-h}-1+h\le h^2\).
Substitution of \(s=e^{\pm h}\) into
\eqref{eq:thinning-transform} therefore yields
\(\E e^{\pm h(X-\mu)}\le e^{\mu h^2}\).
Markov's inequality with \(h=t/(2(\mu+t))\) bounds each tail by
\(\exp[-t^2/(4(\mu+t))]\), proving
\eqref{eq:thinning-bernstein}.
Taking \(s=1/2\) in \eqref{eq:thinning-transform} and applying
Markov's inequality at the thresholds \(2^{-\mu/2}\) and \(1/2\)
proves the two bounds in \eqref{eq:thinning-lower-tail}.

In the shuffle matrix each target-digit pile is a uniform fixed-size
subset of the row labels, independently across columns. After prescribed
symbols have been fixed on selected rows, exchangeability makes each
residual column uniform with its residual counts. The conditioning event
factors over columns, so these residual columns remain independent.
The same argument therefore applies with exactly the residual parameters
specified in the statement.
\end{proof}

% Lean source: formalization/lower_bound_v2/RiffleLowerV2/Thinning/Prefixes.lean :: RiffleLowerV2.Thinning.allPrefixConcentration
% Lean source: formalization/lower_bound_v2/RiffleLowerV2/Thinning/UniformError.lean :: RiffleLowerV2.Thinning.prefixFailureUniform
\begin{lemma}[Uniform prefix-count concentration]
% Lean: formalization/lower_bound_v2/RiffleLowerV2/Thinning/Prefixes.lean :: RiffleLowerV2.Thinning.allPrefixConcentration
% Lean: formalization/lower_bound_v2/RiffleLowerV2/Thinning/UniformError.lean :: RiffleLowerV2.Thinning.prefixFailureUniform
\label{cor:iterated-hypergeom}
For every fixed $\varepsilon>0$, the probability that the following inequality fails for any prefix is at most
\[
 2(K+1)k^K\exp(-N^\varepsilon/8)=o(1).
\]
Thus, with probability $1-o(1)$,
simultaneously for every word \(w\) of length at most \(K\),
\begin{equation}
\label{eq:all-prefix-counts}
 |C(w)-N\lambda_w|\le
 N^\varepsilon(\sqrt{N\lambda_w}+1).
\end{equation}
The convergence is uniform over all deterministic column counts.
\end{lemma}

\begin{proof}
Set \(A=N^\varepsilon\), \(\mu=N\lambda_w\), and
\(t=A(\sqrt\mu+1)\). For \(A\ge1\),
\(t^2/(\mu+t)\ge A/2\).
Lemma~\ref{lem:successive-thinning-tail} bounds the failure probability
for this word by \(2e^{-N^\varepsilon/8}\).
There are at most \((K+1)k^{K}\le (T\log N+1)N^{T\log k}\)
words of the indicated lengths. A union bound proves
\eqref{eq:all-prefix-counts}, including words of zero mass, whose counts
are identically zero.
\end{proof}

\subsubsection{The largest word}

% Lean source: formalization/lower_bound_v2/RiffleLowerV2/EmpiricalLower/LargestFinite.lean :: RiffleLowerV2.EmpiricalLower.finiteLargestWordTV
% Lean source: formalization/lower_bound_v2/RiffleLowerV2/EmpiricalLower/LargestAsymptotic.lean :: RiffleLowerV2.EmpiricalLower.largestMean_uniform
% Lean source: formalization/lower_bound_v2/RiffleLowerV2/Thinning/MaximalWord.lean :: RiffleLowerV2.Thinning.maximalMean_eq_rpow_beta
\begin{proposition}[Finite largest-word bound]
% Lean: formalization/lower_bound_v2/RiffleLowerV2/EmpiricalLower/LargestFinite.lean :: RiffleLowerV2.EmpiricalLower.finiteLargestWordTV
% Lean: formalization/lower_bound_v2/RiffleLowerV2/EmpiricalLower/LargestAsymptotic.lean :: RiffleLowerV2.EmpiricalLower.largestMean_uniform
% Lean: formalization/lower_bound_v2/RiffleLowerV2/Thinning/MaximalWord.lean :: RiffleLowerV2.Thinning.maximalMean_eq_rpow_beta
\label{prop:largest-prefix}
Let
\(\mu_*=N\prod_{j=1}^{K}\max_\ell p^{(j)}_\ell\).
For every integer \(2\le m\le N\) with \(m\le\mu_*/2\),
\begin{equation}
\label{eq:finite-largest-word-tv}
 \dTV(P_N^{K},\nuUnif)
 \ge1-e^{-\frac{1-\log2}{2}\mu_*}-\frac{N-m+1}{m!}.
\end{equation}
The exact identity $\mu_*=N^{\beta^{(N)}(S_N)}$ holds, where $S_N=K/\log N$. In particular, if \(\mu_*\ge N^b\) for a fixed \(b>0\), the
total-variation distance tends to one, uniformly over these column counts.
\end{proposition}

\begin{proof}
Choose a maximizing digit in every column, breaking ties by the smallest
digit, and let \(w_*\) be the resulting full word. Its count has mean
\(\mu_*\). By \eqref{eq:thinning-lower-tail} it is at least \(m\)
with probability at least \(1-e^{-\frac{1-\log2}{2}\mu_*}\).
Equal full words form one component of the shuffle graph. By
the component-sorting representation, the inverse shuffled permutation
therefore has an increasing run of length \(m\) with at least this
probability. Under a uniform permutation each interval of \(m\) consecutive
positions is increasing with probability \(1/m!\). A union bound over
its \(N-m+1\) possible starting positions proves
\eqref{eq:finite-largest-word-tv}. For the last assertion take
\(m=\lfloor N^{b/2}\rfloor\); the hypotheses force \(b\le1\), and
for all sufficiently large \(N\) this choice satisfies the required
inequalities and \((N-m+1)/m!\to0\).
\end{proof}
\subsection{Cold spots}

For $H\subseteq[N]$, let $E(H)$ be the path edges with both endpoints in $H$, and let $\partial H$ be its boundary in the infinite integer path, including the edges from $0$ to $1$ and from $N$ to $N+1$ when present. For a permutation \(\sigma\in\mathfrak S_N\) and a set \(H\subseteq[N]\),
write
\[
  X_H(\sigma)=\sum_{(i,i+1)\in E(H)}
                 \mathbf 1_{\{\sigma(i)<\sigma(i+1)\}}.
\]
If \(G\) is a subgraph of the path on \([N]\), list its nonempty
components from left to right as \(C_1,\ldots,C_q\).  The permutation
\(\pi_G\) is obtained by sorting the values of \(\pi\) increasingly
on each \(C_j\), and we put \(X_{G,H}(\pi)=X_H(\pi_G)\).
All probability and expectation statements with subscript \(\pi\) below
refer to a uniform permutation, with \(G\) and \(H\) fixed.

% Lean source: formalization/lower_bound_v2/RiffleLowerV2/Ascent/RelativeOrder/Law.lean :: RiffleLowerV2.Ascent.relativeOrderIndependence
% Lean source: formalization/lower_bound_v2/RiffleLowerV2/Ascent/RelativeOrder/Boundary.lean :: RiffleLowerV2.Ascent.boundaryAscentProbability
\begin{lemma}[Independent relative orders]
% Lean: formalization/lower_bound_v2/RiffleLowerV2/Ascent/RelativeOrder/Law.lean :: RiffleLowerV2.Ascent.relativeOrderIndependence
% Lean: formalization/lower_bound_v2/RiffleLowerV2/Ascent/RelativeOrder/Boundary.lean :: RiffleLowerV2.Ascent.boundaryAscentProbability
\label{lem:relative-order-independence}
Let \(A_1,\ldots,A_m\) be pairwise disjoint sets of positions in
\([N]\).  Under a uniform permutation \(\pi\), the relative orders of
the values on these sets are mutually independent and uniform.  In
particular, if \(A\) and \(B\) are disjoint nonempty sets of sizes
\(a\) and \(b\), then
\[
  \Pp_\pi\{\max\pi(A)<\min\pi(B)\}
  =\binom{a+b}{a}^{-1}.
\]
\end{lemma}

\begin{proof}
Put \(d_j=|A_j|\) and \(d_0=N-\sum_{j=1}^m d_j\).  Fix one
relative order on each \(A_j\), using increasing position order to
enumerate its entries.  There are
\(N!/(d_0!\prod_jd_j!)\) ways to assign disjoint value sets of the
prescribed sizes to these position sets and their complement.  Each
assignment admits exactly one ordering on each \(A_j\) with the
prescribed relative order and \(d_0!\) orderings on the complement.
Thus precisely \(N!/\prod_jd_j!\) permutations realize the prescribed
orders.  Dividing by \(N!\) proves the product uniform law.  On
\(A\cup B\), exactly \(a!b!\) of the \((a+b)!\) relative orders
put all values on \(A\) before those on \(B\), proving the last claim.
\end{proof}

% Lean source: formalization/lower_bound_v2/RiffleLowerV2/Ascent/Concentration.lean :: RiffleLowerV2.Ascent.componentConcentration; RiffleLowerV2.Ascent.componentConcentrationTwoSided
\begin{lemma}[Component-sorting concentration]
% Lean: formalization/lower_bound_v2/RiffleLowerV2/Ascent/Concentration.lean :: RiffleLowerV2.Ascent.componentConcentration; RiffleLowerV2.Ascent.componentConcentrationTwoSided
\label{lem:permutation-bd}
Let \(G\) be a deterministic path subgraph, let \(H\subseteq[N]\),
and let \(M>0\) satisfy \(|H|\le M\).  For every \(t>0\),
\begin{equation}
\label{eq:component-ascent-one-sided}
  \Pp_\pi\{X_{G,H}-\E_\pi X_{G,H}\ge t\},\quad
  \Pp_\pi\{X_{G,H}-\E_\pi X_{G,H}\le-t\}
  \le \exp(-t^2/M).
\end{equation}
Consequently,
\[
  \Pp_\pi\{|X_{G,H}-\E_\pi X_{G,H}|\ge t\}
  \le 2\exp(-t^2/M).
\]
No condition on \(\partial H\) is needed for these estimates.
\end{lemma}

\begin{proof}
Every edge of \(E(G)\cap E(H)\) is an ascent after component
sorting.  Let \(J\subseteq\{1,\ldots,q-1\}\) consist of the
indices for which the edge joining \(C_j\) to \(C_{j+1}\) belongs
to \(E(H)\), and put \(B=|J|\).  Then
\[
  X_{G,H}=|E(G)\cap E(H)|+\sum_{j\in J}Y_j,
  \qquad
  Y_j=\mathbf 1_{\{\max\pi(C_j)<\min\pi(C_{j+1})\}}.
\]
If \(B=0\), this variable is constant and both probabilities in
\eqref{eq:component-ascent-one-sided} vanish.  Assume \(B>0\).
Separate \(J\) into its even and odd component indices.  Within
either part, the sets \(C_j\cup C_{j+1}\) are disjoint, so
Lemma~\ref{lem:relative-order-independence} makes the corresponding
Bernoulli variables mutually independent.  Their means need not equal
\(1/2\).

For completeness, every Bernoulli variable \(Y\) satisfies
\[
  \E\exp\{\lambda(Y-\E Y)\}\le\exp(\lambda^2/8)
  \qquad(\lambda\in\mathbb R).
\]
Indeed the logarithm of the left side has value and first derivative
zero at the origin, and its second derivative is the variance of a
Bernoulli variable under exponential tilting, hence is at most \(1/4\).
Taylor's theorem gives the asserted bound, including the constant
Bernoulli cases.  Let \(S_0,S_1\) be the centered sums over the two
parity classes.  Cauchy--Schwarz and the independence within each class
give
\[
  \E_\pi e^{\lambda(S_0+S_1)}
  \le
  \bigl(\E_\pi e^{2\lambda S_0}
        \E_\pi e^{2\lambda S_1}\bigr)^{1/2}
  \le \exp(\lambda^2B/4).
\]
Markov's inequality, with \(\lambda=2t/B\), bounds either one-sided
tail by \(\exp(-t^2/B)\).  Finally,
\(B\le|E(H)|\le|H|\le M\), which proves
\eqref{eq:component-ascent-one-sided}.  Adding its two estimates proves
the two-sided bound.
\end{proof}

% Lean source: formalization/lower_bound_v2/RiffleLowerV2/Ascent/Expectation.lean :: RiffleLowerV2.Ascent.componentExpectation
\begin{lemma}[Finite component expectation comparison]
% Lean: formalization/lower_bound_v2/RiffleLowerV2/Ascent/Expectation.lean :: RiffleLowerV2.Ascent.componentExpectation
\label{lem:component-ascent-expectation}
For every deterministic path subgraph \(G\) and set \(H\subseteq[N]\),
put \(e=|E(H)|\), \(r=|E(G)\cap E(H)|\), and
\(d=|\partial H|\), with the boundary taken in the integer path.
Then
\begin{equation}
\label{eq:component-ascent-expectation}
  \E_\pi X_{G,H}\ge \frac e2+\frac r6-d.
\end{equation}
\end{lemma}

\begin{proof}
Call a component \(C\) of \(G\) partial if it meets both \(H\)
and its complement in \([N]\).  Every partial component contains an
edge of \(\partial H\), and different components contain different
such edges.  Thus there are at most \(d\) partial components.
Call an edge of \(E(H)\setminus E(G)\) bad if at least one of its
two incident components is partial.  A path component is incident to
at most two intercomponent edges, so there are at most \(2d\) bad
edges.  Each has ascent probability at least zero and therefore loses
at most \(1/2\) against the baseline \(1/2\).

The two components incident to any remaining intercomponent edge are
wholly contained in \(H\).  If their sizes are \(a,b\ge1\), its
ascent probability is \(\binom{a+b}{a}^{-1}\), by
Lemma~\ref{lem:relative-order-independence}, and
\begin{equation}
\label{eq:component-boundary-deficit}
  \frac12-\binom{a+b}{a}^{-1}
  \le \frac{(a-1)+(b-1)}6.
\end{equation}
For \(a+b\ge5\), the right side is at least \(1/2\).
For \(a+b=2\) both sides are zero; for \(a+b=3\) both sides are
\(1/6\); and for \(a+b=4\) the left side is either \(1/4\)
or \(1/3\), while the right side is \(1/3\).
Each component wholly contained in \(H\) is charged at most twice
in \eqref{eq:component-boundary-deficit}.  Its \(|C|-1\) internal
edges belong to \(E(G)\cap E(H)\).  The total deficit from these
remaining intercomponent edges is therefore at most \(r/3\).
All \(r\) internal graph edges, including those in partial components,
are retained and contribute excess \(r/2\).  Combining the internal
excess, the deficit at most \(r/3\), and the bad-edge deficit at most
\(d\), gives \eqref{eq:component-ascent-expectation}.
\end{proof}

% Lean source: formalization/lower_bound_v2/RiffleLowerV2/Ascent/Concentration.lean :: RiffleLowerV2.Ascent.uniformScan
% Lean source: formalization/lower_bound_v2/RiffleLowerV2/Ascent/Asymptotic.lean :: RiffleLowerV2.Ascent.uniform_null_scan_eventually
\begin{lemma}[Uniform translate-scan ascent bound]
% Lean: formalization/lower_bound_v2/RiffleLowerV2/Ascent/Concentration.lean :: RiffleLowerV2.Ascent.uniformScan
% Lean: formalization/lower_bound_v2/RiffleLowerV2/Ascent/Asymptotic.lean :: RiffleLowerV2.Ascent.uniform_null_scan_eventually
\label{lem:uniform-ascent}
Let \(H_0\subset\mathbb Z\) be finite with \(M=|H_0|>0\),
and put \(H_q=(q+H_0)\cap[N]\) for \(0\le q\le N\).
For every \(s>0\), the event
\[
  \mathcal A(s)=
  \{\sigma:\exists q\in\{0,\ldots,N\},\
          X_{H_q}(\sigma)>|E(H_q)|/2+s\}
\]
satisfies
\begin{equation}
\label{eq:uniform-ascent-scan-finite}
  \nuUnif(\mathcal A(s))\le (N+1)\exp(-s^2/M).
\end{equation}
In particular, if \(M\ge N^a\) for a fixed \(a>0\), then for
every fixed \(\delta>0\),
\[
  \nuUnif\bigl(\mathcal A(M^{1/2+\delta/2})\bigr)\longrightarrow0.
\]
\end{lemma}

\begin{proof}
A comparison between two distinct positions of a uniform permutation
has mean \(1/2\), so \(\E_\pi X_{H_q}=|E(H_q)|/2\).
Apply Lemma~\ref{lem:permutation-bd} to the edgeless graph and
\(H_q\), for which \(|H_q|\le M\), and take a union bound over
the \(N+1\) choices of \(q\).  This proves
\eqref{eq:uniform-ascent-scan-finite}.  With
\(s=M^{1/2+\delta/2}\), its right side is at most
\((N+1)\exp(-N^{a\delta})\), which tends to zero.
\end{proof}

% Lean source: formalization/lower_bound_v2/RiffleLowerV2/Ascent/Test.lean :: RiffleLowerV2.Ascent.componentComparison
% Lean source: formalization/lower_bound_v2/RiffleLowerV2/Ascent/Quantitative.lean :: RiffleLowerV2.Ascent.quantitative_component_failure
\begin{lemma}[Lalley component comparison]
% Lean: formalization/lower_bound_v2/RiffleLowerV2/Ascent/Test.lean :: RiffleLowerV2.Ascent.componentComparison
% Lean: formalization/lower_bound_v2/RiffleLowerV2/Ascent/Quantitative.lean :: RiffleLowerV2.Ascent.quantitative_component_failure
\label{lem:lalley-comparison}
Let \(M>0\), \(|H|\le M\), and \(G\) be a deterministic
path subgraph.  Write \(e=|E(H)|\),
\(r=|E(G)\cap E(H)|\), and \(d=|\partial H|\).
If \(s>0\) and \(\Delta=r/6-d-s>0\), then
\begin{equation}
\label{eq:lalley-comparison-finite}
  \Pp_\pi\{X_{G,H}\le e/2+s\}
  \le\exp(-\Delta^2/M).
\end{equation}
Consequently, for fixed \(D\ge0\), \(\delta>0\), and any such
family with \(M\to\infty\), \(d\le D\sqrt M\), and
\(r\ge M^{1/2+\delta}\),
\[
  \Pp_\pi\{X_{G,H}>e/2+M^{1/2+\delta/2}\}\longrightarrow1
\]
uniformly over \(N,G,H\).
\end{lemma}

\begin{proof}
Lemma~\ref{lem:component-ascent-expectation} gives
\(\E_\pi X_{G,H}-(e/2+s)\ge\Delta\).
The event in \eqref{eq:lalley-comparison-finite} is therefore
contained in the lower-tail event from
\eqref{eq:component-ascent-one-sided}, with \(t=\Delta\).
For the last assertion take \(s=M^{1/2+\delta/2}\).
For all sufficiently large \(M\), depending only on \(D,\delta\),
\[
  D+M^{\delta/2}\le M^\delta/12,
  \qquad
  \Delta\ge M^{1/2+\delta}/12.
\]
Thus the failure probability is at most
\(\exp(-M^{2\delta}/144)\), uniformly as claimed.
\end{proof}

% Lean source: formalization/lower_bound_v2/RiffleLowerV2/Ascent/Test.lean :: RiffleLowerV2.Ascent.conditionalLalleyTest; RiffleLowerV2.Ascent.conditionalLalleyTendsto
\begin{lemma}[Lalley cold-spot test]
% Lean: formalization/lower_bound_v2/RiffleLowerV2/Ascent/Test.lean :: RiffleLowerV2.Ascent.conditionalLalleyTest; RiffleLowerV2.Ascent.conditionalLalleyTendsto
\label{lem:lalley}
Let \(H_0\subset\mathbb Z\) be deterministic and finite, put
\(M=|H_0|\), and let \(H_q=(q+H_0)\cap[N]\).
Suppose that, for constants \(a,\delta>0\) and \(D\ge0\)
independent of \(N\),
\[
  M\ge N^a,\qquad |\partial H_q|\le D\sqrt M
  \quad(0\le q\le N),
\]
and that the shuffle graph satisfies
\[
  \Pp\{\exists q\in\{0,\ldots,N\}:\
         |E(G)\cap E(H_q)|\ge M^{1/2+\delta}\}\longrightarrow1.
\]
Then \(\dTV(P_N^{K},\nuUnif)\longrightarrow1\).
\end{lemma}

\begin{proof}
Use the observable event
\(\mathcal A=\mathcal A(M^{1/2+\delta/2})\) from
Lemma~\ref{lem:uniform-ascent}.  That lemma gives
\(\nuUnif(\mathcal A)\to0\).
Use the representation in the component-sorting representation
to sample the shuffle graph and an independent uniform permutation
\(\pi\), and then form \(\pi_G\).
On the event that a translate has at least
\(M^{1/2+\delta}\) internal graph edges, choose its smallest index
\(q\).  This choice depends only on \(G\).
Conditional on \(G\), Lemma~\ref{lem:lalley-comparison} bounds the
probability that this translate fails the ascent test by
\(\exp(-M^{2\delta}/144)\), for all sufficiently large \(N\).
The bound is uniform in \(G\) and the chosen translate.  Hence
\[
  P_N^K(\mathcal A)
  \ge
  \Pp\{\exists q:\ |E(G)\cap E(H_q)|\ge M^{1/2+\delta}\}
     -\exp(-M^{2\delta}/144)
  \longrightarrow1.
\]
Finally,
\(\dTV(P_N^{K},\nuUnif)
\ge P_N^K(\mathcal A)-\nuUnif(\mathcal A)\to1\).
\end{proof}

\subsection{Construction of the cold spots}
\subsubsection{Empirical costs}

Write $L=\log N$ and $S_N=K/L$.  On
$((j-1)/L,j/L]$, $1\le j\le K$, use the actual profile $p^{(j)}$.
For $0\le r\le s\le S_N$, define
\[
 (A^{(N)},B_\alpha^{(N)},J_\alpha^{(N)},H_\alpha^{(N)})(r,s)
 =\int_r^s(a_{p_N(t)},\psi_{p_N(t)}(\alpha),
                 I_{p_N(t)}(\alpha),H_{p_N(t)}(\alpha))\,dt,
\]
where $p_N(t)=p^{(j)}$ on the displayed column interval.
Here
\[
 p^{[\alpha]}_\ell=\frac{p_\ell^\alpha}{\sum_a p_a^\alpha},\qquad
 I_p(\alpha)=\sum_{\ell:p_\ell>0}p^{[\alpha]}_\ell\log(1/p_\ell),\qquad
 H_p(\alpha)=-\sum_{\ell:p_\ell>0}p^{[\alpha]}_\ell\log p^{[\alpha]}_\ell.
\]
These are integrals of bounded step functions.  In particular, at grid
endpoints they are the normalized finite sums of the actual column costs.
Set
\[
 \beta^{(N)}(s)=1-A^{(N)}(0,s),\qquad
 F_\alpha^{(N)}(s,u)=\frac{\alpha+3}{2}\beta^{(N)}(s)
                 -B_\alpha^{(N)}(s,u)-2B_2^{(N)}(u,S_N).
\]
Set $\Lambda_1^{(N)}=\sup_{0\le s\le u\le S_N}
\inf_{3\le\alpha\le4}F_\alpha^{(N)}(s,u)$.
No distributional assumption on the order of the columns is made here.

For $\alpha\in[2,4]$, the one-column functions satisfy
\begin{equation}\label{eq:repair-global-one-column-bounds}
 0\le a_p\le\log k,\quad
 0\le I_p(\alpha)\le\log k,\quad
 0\le H_p(\alpha)\le I_p(\alpha),\quad
 0\le\psi_p(\alpha)\le(\alpha-1)\log k.
\end{equation}
To verify these assertions, apply the finite Gibbs inequality to
$q=p^{[\alpha]}$ and the weights $p_\ell^\zeta$, restricting all sums to
the positive coordinates. It gives
\[
 H_p(\alpha)\le\zeta I_p(\alpha)-\psi_p(\zeta),\qquad
 H_p(\alpha)=\alpha I_p(\alpha)-\psi_p(\alpha).
\]
Apply the inequality twice, with $\alpha$ and $\zeta$ interchanged. For
$1\le\alpha<\zeta$ their sum gives
$(\zeta-\alpha)(I_p(\zeta)-I_p(\alpha))\le0$.
Thus $I_p$ is decreasing, and $I_p(1)=H(p)\le\log k$; the latter
entropy bound follows from the same Gibbs inequality with the uniform
probability vector. Taking $\zeta=1$ gives $H_p(\alpha)\le I_p(\alpha)$.
Taking $q=p$ in Gibbs gives
$\psi_p(\alpha)\le(\alpha-1)H(p)\le(\alpha-1)\log k$.
Nonnegativity follows from $p_\ell\le1$, and $a_p\le I_p(\alpha)$
follows by averaging $\log(1/p_\ell)\ge-\log\max_a p_a$.
All finite sums are continuous in $\alpha$ on $[2,4]$ and have positive
power-sum denominators. Zero coordinates have zero tilted weight.
Finally, $\sum p_\ell^2\ge(\max p_\ell)^2$ and
$\sum p_\ell^4\le(\sum p_\ell^2)^2$ give
\begin{equation}\label{eq:repair-two-comparisons}
 \psi_p(2)\le2a_p\le2I_p(\alpha),\qquad
 \psi_p(4)\ge2\psi_p(2).
\end{equation}Thus every empirical endpoint displacement by $r$ changes an integrated
$I$ or $H$ by at most $r\log k$, uniformly in the cuts and $N$.


\subsubsection{A finite cold witness}

% Lean source: formalization/lower_bound_v2/RiffleLowerV2/Witness/Main.lean :: RiffleLowerV2.Witness.coldWitness
% Lean source: formalization/lower_bound_v2/RiffleLowerV2/Witness/Continuous.lean :: RiffleLowerV2.Witness.continuousWitness
% Lean source: formalization/lower_bound_v2/RiffleLowerV2/Witness/Finite.lean :: RiffleLowerV2.Witness.roundWitness
\begin{lemma}[Finite cold witness]% Lean: formalization/lower_bound_v2/RiffleLowerV2/Witness/Main.lean :: RiffleLowerV2.Witness.coldWitness
% Lean: formalization/lower_bound_v2/RiffleLowerV2/Witness/Continuous.lean :: RiffleLowerV2.Witness.continuousWitness
% Lean: formalization/lower_bound_v2/RiffleLowerV2/Witness/Finite.lean :: RiffleLowerV2.Witness.roundWitness
\label{lem:cold-witness}
Fix $\eta,\rho,\chi_g,\gamma>0$ and assume the local good-column condition on $[0,T]$ with these last three parameters.
Suppose $\Lambda_1^{(N)}\ge\eta$ and
$\beta^{(N)}(S_N)<\min\{\eta/100000,1/1000\}$.
Choose a bridge length $2\rho$ with
$2\rho\log k\le\eta/100000$, where the local good-column hypothesis holds
on every interval of length $\rho$ contained in $[0,T]$. For all sufficiently large $N$ there
are grid endpoints $s=m_s/L$, $u=m_u/L$, a tilt $\theta\in[3,4]$, and an
integer bridge length $b_N=\lceil 2\rho L\rceil$ such that
$0\le m_s\le m_u\le K-b_N$, the bridge columns
$m_u<j\le m_u+b_N$ contain two good columns, and, writing
$b=\beta^{(N)}(s)$, $J=J_\theta^{(N)}(s,u)$,
$H=H_\theta^{(N)}(s,u)$, $J_t=J_2^{(N)}(u,S_N)$,
$H_t=H_2^{(N)}(u,S_N)$,
\begin{gather}
 F_\theta^{(N)}(s,u)\ge\eta/2,\qquad
 J\le b/2-3\eta/1000,\label{eq:repair-witness-upper}\\
 \theta=3\quad\hbox{or}\quad J\ge b/2-\eta/200,
 \label{eq:repair-witness-lower}\\
 h:=b-J+H\ge\eta/1000,\qquad H\le b-J-\eta/1000,
 \label{eq:repair-witness-geometry}\\
 \Gamma:=2b-2J-2J_t+H+H_t,\qquad
 \Gamma-h/2\ge\eta/1000,\label{eq:repair-witness-collision}\\
 b-J-J_t\le\beta^{(N)}(S_N)<\min\{\eta/100000,1/1000\}.
 \label{eq:repair-witness-mean}
\end{gather}
The constants are uniform in the finite cut sequence.
\end{lemma}

\begin{proof}
Choose $s_0\le u_0$ with
$F_\alpha^{(N)}(s_0,u_0)\ge7\eta/8$ for every $\alpha\in[3,4]$.
Write $b_0=\beta^{(N)}(s_0)$.  Nonnegativity of the $B$ terms gives
$b_0\ge\eta/4$, so $b_0/2-\eta/250>0$.
If $J_3^{(N)}(s_0,u_0)\le b_0/2$, choose $\theta=3$.
Otherwise, if $J_4^{(N)}(s_0,u_0)\ge b_0/2$, choose $\theta=4$.
In the remaining case the continuity and monotonicity in $\alpha$ give
$\theta\in(3,4)$ with $J_\theta^{(N)}(s_0,u_0)=b_0/2$.
This selection does not interchange a supremum and an infimum.

If $J_\theta^{(N)}(s_0,u_0)\le b_0/2-\eta/250$, leave $u_1=u_0$;
this case can occur only for $\theta=3$.  Otherwise choose
$u_1\in[s_0,u_0]$ with
$J_\theta^{(N)}(s_0,u_1)=b_0/2-\eta/250$.
Existence follows from continuity and the value zero at $s_0$.
In the $\theta=3$ or interior cases the decrease of $J$ is at most
$\eta/250$.  By \eqref{eq:repair-two-comparisons},
\[
 F_\theta^{(N)}(s_0,u_1)-F_\theta^{(N)}(s_0,u_0)
 =B_\theta^{(N)}(u_1,u_0)-2B_2^{(N)}(u_1,u_0)
 \ge-4J_\theta^{(N)}(u_1,u_0)\ge-2\eta/125.
\]
For $\theta=4$ the same difference is nonnegative, even if the
backward displacement is large.  Hence in every case
$F_\theta^{(N)}(s_0,u_1)\ge7\eta/8-2\eta/125$.

Set
\[
 m_s=\lfloor s_0L\rfloor,\qquad
 m_u=\min\{\lfloor u_1L\rfloor,K-b_N\}.
\]
These integers are ordered for large $N$.  To check the only nontrivial
case, observe that
\[
 (S_N-s_0)\log k\ge A^{(N)}(s_0,S_N)
 =b_0-\beta^{(N)}(S_N)>\eta/4-\min\{\eta/100000,1/1000\}.
\]
Consequently $s_0<S_N-2\rho-2/L$ for large $N$, and clipping $m_u$ to
$K-b_N$ cannot move it below $m_s$.
The endpoint errors satisfy $0\le s_0-s<1/L$ and
$0\le u_1-u\le 2\rho+2/L$.
Equation~\eqref{eq:repair-global-one-column-bounds} bounds the resulting
change of $F$ by $10\rho\log k+20\log k/L$ and that of
$J-b/2$ by $2\rho\log k+4\log k/L$.
These errors are smaller than $\eta/10000$ for large $N$.
The retained slack proves \eqref{eq:repair-witness-upper} and
\eqref{eq:repair-witness-lower}.
For the two good columns, use the interior interval
$[u+\rho/2,u+3\rho/2]$. It lies within the bridge and within $[0,T]$.
Every column time in this interval is strictly above $u$ and at most
$u+b_N/L$, so it belongs to the bridge. The local good-column condition
provides at least $\gamma\rho L>3$ good columns for large $N$; choose two
of them in increasing order.
The exact identities $H=\theta J-B_\theta^{(N)}(s,u)$ and
$H_t=2J_t-B_2^{(N)}(u,S_N)$ imply
\[
 2\Gamma-h=F_\theta^{(N)}(s,u)+(\theta-3)(J-b/2)
 \ge\eta/2-\eta/200.
\]
This is greater than $\eta/500$.
Also $H\le J$ and $J\le b/2-3\eta/1000$ give
$H\le b-J-3\eta/500$ and $h\ge b-J\ge b/2+3\eta/1000$.
Finally $J\ge A^{(N)}(s,u)$ and $J_t\ge A^{(N)}(u,S_N)$ give
\eqref{eq:repair-witness-mean}.
\end{proof}

\subsubsection{Finite type classes}

The construction below uses only the actual columns.  In particular it
needs no temporal approximation for the tilted observables $I$ and $H$.
Fix a $\xi$-regular simplex partition with regularity constant $C\ge0$ and representatives
$p_R$, where $0<\xi<1/(4k)$.  There are $m=m(\xi)$ cells, independent of $N$.
On a chosen set $I$ of column positions, let $I_R$ be the positions in cell
$R$ and $n_R=|I_R|$.  The positions retain their chronological order.
For $\alpha\in[2,4]$ put
\[
 G_R=\{\ell:p_R(\ell)>2\xi\},\qquad
 q_R(\ell)=\frac{p_R(\ell)^\alpha\mathbf1_{G_R}(\ell)}
                 {\sum_{a\in G_R}p_R(a)^\alpha}.
\]
Choose integers $m_{R,\ell}\ge0$ supported on $G_R$, with
$\sum_\ell m_{R,\ell}=n_R$ and
$|m_{R,\ell}-n_Rq_R(\ell)|\le1$; floor the counts and distribute the
remaining units deterministically.  Let $\mathcal C^0_I$ be the words on
$I$ with exactly these counts in each cell.  This class is nonempty.
For every word write $\lambda_x=\prod_{j\in I}p^{(j)}(x_j)$.
Partition $\mathcal C^0_I$ into the integer mass bins
\[
 \{x:j\le-\log\lambda_x<j+1\},\qquad j\in\mathbb Z_{\ge0}.
\]
Select a bin of maximal cardinality, breaking ties by the smallest index, and
call it $\mathcal C_I$.  If its index is $j_I$, define
\[
 J_I^*=j_I/L,\qquad H_I^*=L^{-1}\log|\mathcal C_I|.
\]
These quantities may depend on $N$; that dependence is intentional.

% Lean source: formalization/lower_bound_v2/RiffleLowerV2/Types/Main.lean :: RiffleLowerV2.Types.roundedType
% Lean source: formalization/lower_bound_v2/RiffleLowerV2/Types/Truncation.lean :: RiffleLowerV2.Types.restrictedCost_truncation_error
\begin{lemma}[Finite type estimates]% Lean: formalization/lower_bound_v2/RiffleLowerV2/Types/Main.lean :: RiffleLowerV2.Types.roundedType
% Lean: formalization/lower_bound_v2/RiffleLowerV2/Types/Truncation.lean :: RiffleLowerV2.Types.restrictedCost_truncation_error
\label{lem:type-estimates}
For $|I|\le TL$ and $\alpha\in[2,4]$, uniformly over partitions with the same fixed number $m$ of cells, the cuts, the chosen
position set, and the selected word,
\begin{gather}
 e^{-1}N^{-J_I^*}<\lambda_x\le N^{-J_I^*},\qquad
 |\mathcal C_I|=N^{H_I^*},\label{eq:repair-bin-exact}\\
 \left|J_I^*-L^{-1}\sum_{j\in I}I_{p^{(j)}}(\alpha)\right|
 +\left|H_I^*-L^{-1}\sum_{j\in I}H_{p^{(j)}}(\alpha)\right|
 \le O_{k,C}(T\xi\log(1/\xi))
       +O_{k,m,\xi,T}\!\left(\frac{\log\log N+1}{\log N}\right).
 \label{eq:repair-types-error}
\end{gather}
Every digit occurring in a selected word has actual column probability
at least $\xi$.  Here the first implied constant depends only on $k,C$, and the second only on $k,m,\xi,T$; these constants and the threshold in $N$ precede the partition, cut array, position set, and tilt. The empty position set is allowed.
\end{lemma}

\begin{proof}
Let $r=p_R$ and let $p$ belong to the same cell. On retained coordinates,
$p_\ell\ge\xi$ and $|\log p_\ell-\log r_\ell|\le C\xi$; on omitted
coordinates, $p_\ell\le3\xi$. The power-sum denominator is at least
$k^{-4}$. The omitted power mass is at most $9k\xi^2$. Moreover,
$p_\ell^\alpha\log(1/p_\ell)\le p_\ell$ for $\alpha\ge2$,
by $\log x\le x-1$, so the omitted cost numerator is at most $3k\xi$.
After division by the full power sum these bounds become $9k^5\xi^2$
and $3k^5\xi$. The retained cost is at most $\log(1/\xi)$. Splitting
the full cost into its retained and omitted parts therefore bounds the
truncation error by $18k^5\xi\log(1/\xi)$.

On retained coordinates, the normalized $\alpha$-power weights at $p$
and $r$ have ratios between $e^{-8C\xi}$ and $e^{8C\xi}$. Comparing
these weights and their digit costs gives
\[
 \left|\sum_\ell q_R(\ell)\log(1/r_\ell)-I_p(\alpha)\right|
 \le O_{k,C}(\xi\log(1/\xi)).
\]
The identity expressing tilted entropy as tilt times cost plus the
logarithm of the power sum, and the same retained/omitted decomposition,
give the corresponding entropy estimate. These bounds are uniform on
$0<\xi<1/(4k)$; this interval also ensures that the retained set is
nonempty, since one coordinate of the representative is at least $1/k$.
Put $\bar J=L^{-1}\sum_R n_R\sum_\ell q_R(\ell)\log(1/p_R(\ell))$
and $\bar H=L^{-1}\sum_R n_RH(q_R)$.
The rounding of counts changes the logarithmic product at representatives
by at most $km\log(1/(2\xi))$.
Replacing representative coordinates by actual coordinates changes it by
at most $C\xi|I|$.  Consequently every $x\in\mathcal C^0_I$ satisfies
\[
 |{-L^{-1}\log\lambda_x}-\bar J|
 \le CT\xi+km\log(1/(2\xi))/L.
\]
The multinomial formula gives
$|\mathcal C^0_I|=\prod_Rn_R!/\prod_\ell m_{R,\ell}!$.
The elementary factorial estimate
$\log(n!)=n\log n-n+O(\log(n+1)+1)$ and the entropy continuity bound
\[
 |n_RH((m_{R,\ell}/n_R)_\ell)-n_RH(q_R)|
 \le O_k(\log(n_R+1)+1)
\]
show the asserted cardinality estimate: if $f(t)=t\log t$, with $f(0)=0$,
then $|f(x)-f(y)|\le\log(n+1)+2$ for $x,y\in[0,n]$ with
$|x-y|\le1$. This follows by applying $\log t\le t-1$ to the ratio of the two positive arguments and using $|t\log t|\le1/e$ on $[0,1]$ when an argument is small. Apply it to
$x=m_{R,\ell}$ and $y=n_Rq_R(\ell)$ and sum over the coordinates;
the empty-cell contribution is zero. Consequently,
$|L^{-1}\log|\mathcal C^0_I|-\bar H|
 =O_{k,m,T}((\log\log N+1)/L)$.
Since every allowed actual digit is at least $\xi$, there are at most
$2+TL\log(1/\xi)$ nonempty integer mass bins.
Keeping a largest bin therefore changes the logarithmic cardinality by
$O_{\xi,T}(\log\log N+1)$ only.  Its index divided by $L$ lies within
$1/L$ of the logarithmic mass of each selected word.
Combining these estimates proves both assertions.
\end{proof}

\subsubsection{Cold sets and bridges}

For a grid endpoint $s=m_s/\log N$, let $a_s$ be the length-$m_s$ word
which uses a maximizing digit in every column, with ties broken by the
smallest digit. For any word on a specified consecutive set of columns,
$\lambda_x$ denotes its product probability on those columns.
For $x$ on $(m_s,m_u]$, put
\[
 t_x^{\rm rel}=\sum_{\substack{x'\in[k]_0^{m_u-m_s}\\x'<_{\rm lex}x}}
                \lambda_{x'},\qquad
 J_x^{\rm rel}=[t_x^{\rm rel},t_x^{\rm rel}+\lambda_x).
\]

% Lean source: formalization/lower_bound_v2/RiffleLowerV2/Localization/Geometry.lean :: RiffleLowerV2.Localization.deterministicColdGeometry
% Lean source: formalization/lower_bound_v2/RiffleLowerV2/Localization/ColdSet.lean :: RiffleLowerV2.Localization.coldRelativeSet_estimates
\begin{lemma}[Finite cold geometry and mass bounds]
% Lean: formalization/lower_bound_v2/RiffleLowerV2/Localization/Geometry.lean :: RiffleLowerV2.Localization.deterministicColdGeometry
% Lean: formalization/lower_bound_v2/RiffleLowerV2/Localization/ColdSet.lean :: RiffleLowerV2.Localization.coldRelativeSet_estimates
\label{lem:cold-geometry}
Fix $\eta,\rho,\chi_g>0$, subject to the
bridge budget $2\rho\log k\le\min\{\eta/100000,1/1000\}/100$. There exist $0<\xi<1/(4k)$ and a $\xi$-regular finite partition with regularity constant $1$, fixed before $N$. For every sufficiently large $N$, every cut array, and every choice of $s,u,\theta,b_N$ satisfying all the conclusions of Lemma~\ref{lem:cold-witness}, there are deterministic prefix
and tail families $\mathcal P_N,\mathcal T_N$, a bridge word $v_b$ with
reserve words $v_b^-,v_b^+$, and a deterministic set $H_{{\rm rel},N}$,
as constructed in the proof. The classes $\mathcal P_N,\mathcal T_N$
are obtained from Lemma~\ref{lem:type-estimates} on the column sets
$(m_s,m_u]$ with tilt $\theta$ and $(m_u+b_N,K]$ with tilt $2$,
respectively; denote their mass-bin exponents by $J_p^*,H_p^*$ and
$J_t^*,H_t^*$. Set
\[
 H_{{\rm rel},N}=\mathbb Z\cap
    \bigcup_{x\in\mathcal P_N}N\lambda_{a_s}J_x^{\rm rel},\qquad
 M_N=|H_{{\rm rel},N}|,\qquad H_q=(q+H_{{\rm rel},N})\cap[N].
\]
Define the real exponents
\[
 \ell_p=b-J_p^*,\quad h_N=\ell_p+H_p^*,\quad
 \ell_b=-\log_N\lambda_{v_b},\quad
 z_N=b-J_p^*-J_t^*-\ell_b,\quad
 g_N=2z_N+H_p^*+H_t^*.
\]
For all sufficiently large $N$,
\begin{gather}
 \ell_p\ge b/2+\eta/500,\qquad H_p^*\le\ell_p-\eta/1000,\qquad
 h_N\ge\eta/2000,\label{eq:cold-effective-geometry}\\
 e^{-2}N^{h_N}\le M_N\le2N^{h_N},\qquad
 N^{\eta/3000}\le M_N\le N,\qquad
 |\partial H_q|\le4\sqrt{M_N}\quad(0\le q\le N),
 \label{eq:cold-geometric-conclusions}\\
 g_N-h_N/2\ge\eta/2000,\qquad z_N\le2\min\{\eta/100000,1/1000\},
 \label{eq:cold-effective-collision}
\end{gather}
The full-word family $\mathcal Y_N=\{a_sxv_bz:x\in\mathcal P_N,z\in\mathcal T_N\}$
satisfies, uniformly in its words,
\begin{equation}
\label{eq:cold-full-word-masses}
 |\mathcal Y_N|=N^{H_p^*+H_t^*},\qquad
 e^{-2}N^{z_N}<N\lambda_y\le N^{z_N}.
\end{equation}
The reserves satisfy $v_b^-<_{\rm lex}v_b<_{\rm lex}v_b^+$ and
\begin{equation}
\label{eq:cold-reserve-means}
 \min_{x\in\mathcal P_N,\,\diamond\in\{-,+\}}
       N\lambda_{a_sxv_b^\diamond}\ge N^{b/2+\eta/1000}.
\end{equation}
Every digit selected by a full word or a reserve has column probability
at least $\min\{\xi,\chi_g,1/k\}>0$. Both type approximations have error at most $\min\{\eta/100000,1/1000\}/1000$. The bridge words differ only at their distinguished columns as described below, and each has mass at least $\chi_g^2k^{-b_N}$. Their central bridge exponent satisfies
\[
 0\le\ell_b\le2\rho\log k+
 \frac{\log k+2\log(1/\chi_g)}{\log N}.
\]
All constants and vanishing errors are uniform
over the finite cut sequences under consideration.
\end{lemma}

\begin{proof}

Use the finite witness, and apply Lemma~\ref{lem:type-estimates} to
$I=(m_s,m_u]$ with $\alpha=\theta$.
Call the resulting class $\mathcal P_N$, with exponents $J_p^*,H_p^*$.
The ambient maximal prefix $a_s$ has exactly
$N\lambda_{a_s}=N^b$ because $s$ is a grid endpoint.
Set
\[
 \ell_p=b-J_p^*,\qquad h_N=\ell_p+H_p^*.
\]
Choose $\xi$, a regular partition with constant $1$, and then $N$ so that each type error in \eqref{eq:repair-types-error} is at most $\min\{\eta/100000,1/1000\}/1000$. The type estimates choose these parameters uniformly before the cut array and the two selected column intervals.
Then
\begin{equation}\label{eq:repair-finite-geometric-buffers}
 \ell_p\ge b/2+\eta/500,\qquad
 H_p^*\le\ell_p-\eta/1000,\qquad h_N\ge\eta/2000.
\end{equation}
For each $x\in\mathcal P_N$ use its exact relative probability interval
$J_x^{\rm rel}$, and put
\[
 H_{{\rm rel},N}=\mathbb Z\cap\bigcup_{x\in\mathcal P_N}
                         N\lambda_{a_s}J_x^{\rm rel}.
\]
The intervals are disjoint and each has length between
$e^{-1}N^{\ell_p}$ and $N^{\ell_p}$.  Their integer counts differ from their
lengths by at most one. Since $N^{\ell_p}\ge e^2$ for sufficiently large $N$, summing the interval counts gives exactly $e^{-2}N^{h_N}\le M_N\le2N^{h_N}$. Therefore, uniformly in $N$,
\[
 |H_{{\rm rel},N}|=N^{h_N+o(1)}\ge N^{\eta/3000},\qquad
 |\partial H_{{\rm rel},N}|\le2N^{H_p^*}
                         =o(|H_{{\rm rel},N}|^{1/2}).
\]
This does not assert that every probability interval is a distinct connected
component: adjacent selected intervals may merge, which only reduces the
boundary.  Translations clipped to $[N]$ have no larger cardinality and
have at most two additional boundary edges. The gap $H_p^*\le\ell_p-\eta/1000$ gives $H_p^*\le h_N/2-\eta/2000$. Hence, for all sufficiently large $N$, $2N^{H_p^*}+2\le4e^{-1}N^{h_N/2}\le4\sqrt{M_N}$, proving the stated coefficient $4$.

Choose two good bridge columns $g_L<g_R$ and define the word $v_b$ as follows.
At $g_L$ use the larger of two good digits; at $g_R$ use the smaller.
Use maximal digits at all other bridge positions.
For the reserve words $v_b^-,v_b^+$, change respectively the digit at
$g_L$ to the smaller and the digit at $g_R$ to the larger, keeping all
other digits equal to the central bridge word.
Then $v_b^-<_{\rm lex}v_b<v_b^+$, and each bridge word has probability at least
$\chi_g^2k^{-b_N}$, where $\chi_g$ is the fixed good-column threshold.
Writing $\ell_b=-L^{-1}\log\lambda_{v_b}$, we obtain
\[
 0\le\ell_b\le 2\rho\log k+(\log k+2\log(1/\chi_g))/L.
\]
The same estimate applies to both reserve words.  In particular every
reserve prefix $a_sxv_b^\pm$ has mean at least
$N^{b/2+\eta/1000}$ for large $N$, after the bridge budget is imposed.
This is the deterministic input for the later concentration/localization
argument; it makes no assertion of independence between reserve counts.

Let $\mathcal R=(m_u+b_N,K]$ and apply
Lemma~\ref{lem:type-estimates} there with $\alpha=2$.
Write $\mathcal T_N,J_t^*,H_t^*$ for the selected tail class and exponents.
The bridge deletion changes the empirical $J_2,H_2$ by at most
$(2\rho+1/L)\log k$ each.
The full selected family is
\[
 \mathcal Y_N=\{a_sxv_bz:x\in\mathcal P_N, z\in\mathcal T_N\}.
\]
For $y\in\mathcal Y_N$ put $\mu_y=N\lambda_y$ and define
\[
 z_N=b-J_p^*-J_t^*-\ell_b,\qquad
 g_N=2z_N+H_p^*+H_t^*.
\]
The mass bins give the exact uniform bounds
\[
 e^{-2}N^{z_N}<\mu_y\le N^{z_N},\qquad
 |\mathcal Y_N|=N^{H_p^*+H_t^*},\qquad
 \sum_{y\in\mathcal Y_N}\mu_y^2=N^{g_N+o(1)}.
\]
Write
\[
 \delta_{p,J}=J_p^*-J,\quad \delta_{p,H}=H_p^*-H,\quad
 \delta_{t,J}=J_t^*-(J_t-d_J),\quad
 \delta_{t,H}=H_t^*-(H_t-d_H),
\]
where $d_J=J_2^{(N)}(u,u+b_N/L)$ and
$d_H=H_2^{(N)}(u,u+b_N/L)$ are the deleted bridge integrals.
Each $\delta$ has absolute value at most $\min\{\eta/100000,1/1000\}/1000$ and
$0\le d_J,d_H\le d:=(2\rho+1/L)\log k$.
For all sufficiently large $N$, our budgets give
$d,\ell_b\le\min\{\eta/100000,1/1000\}/50$. Direct subtraction now yields
\begin{align*}
 |h_N-h|&\le\frac1{500}\min\{\eta/100000,1/1000\},\\
 g_N-h_N/2&\ge\Gamma-h/2
   -\frac7{1000}\min\{\eta/100000,1/1000\}-d-2\ell_b
   \ge\eta/2000.
\end{align*}
and
\begin{align*}
 z_N&\le b-J-J_t+\frac1{500}\min\{\eta/100000,1/1000\}+d\\
    &<\frac{501}{500}\min\{\eta/100000,1/1000\}+d\\
    &\le2\min\{\eta/100000,1/1000\}.
\end{align*}
These prove \eqref{eq:cold-effective-collision}. The reserve mass lower
bound proved above gives \eqref{eq:cold-reserve-means}.
All selected digits have probability at least
$c=\min\{\xi,\chi_g,1/k\}>0$. Finally
$H_{{\rm rel},N}\subset\mathbb Z\cap[0,N\lambda_{a_s})\subset[0,N)\cap\mathbb Z$,
so $M_N\le N$. The bounds already established give
\eqref{eq:cold-geometric-conclusions}, including its uniform clipped-boundary
estimate, and complete the proof.
\end{proof}

\subsubsection{Localization}

% Lean source: formalization/lower_bound_v2/RiffleLowerV2/Localization/Endpoints.lean :: RiffleLowerV2.Localization.relativeEndpoint
\begin{lemma}[Finite relative endpoint estimate]
% Lean: formalization/lower_bound_v2/RiffleLowerV2/Localization/Endpoints.lean :: RiffleLowerV2.Localization.relativeEndpoint
\label{lem:finite-relative-endpoints}
Let \(a\) be a deterministic prefix of length \(l\), let
\(l\le m\le K\), and put \(\mu_a=N\lambda_a\).
For a suffix \(x\) on columns \(l+1,\ldots,m\), let
\(\lambda_x\) be its product mass and \(t_x\) the total product
mass of the length-\((m-l)\) suffixes strictly smaller than \(x\).
Define \(L_x^{\mathrm{rel}},U_x^{\mathrm{rel}}\) by counting rows
with prefix \(a\) whose length-\((m-l)\) suffixes are respectively
smaller than, or at most, \(x\).
If \(A>0\) and
\[
 |C(v)-N\lambda_v|\le A(\sqrt{N\lambda_v}+1)
 \quad\hbox{for every prefix \(v\) of length at most \(K\)},
\]
then simultaneously for all such \(a,x\),
\begin{equation}
\label{eq:relative-endpoint-finite-error}
\begin{split}
 &|L_x^{\mathrm{rel}}-\mu_at_x|
  +|U_x^{\mathrm{rel}}-\mu_a(t_x+\lambda_x)|\\
 &\hspace{6mm}\le
 A\bigl[(2\sqrt{kK}+1)\sqrt{\mu_a}+2kK+1\bigr].
\end{split}
\end{equation}
\end{lemma}

\begin{proof}
Partition the preceding and following suffixes by the first column
where they differ from \(x\).  The resulting exact identities are
\begin{align}
\label{eq:relative-left-side-exact}
 L_x^{\mathrm{rel}}-\mu_at_x
 &=\sum_{r=l+1}^m\sum_{d<x_r}
   \bigl[C(ax_{l+1}\cdots x_{r-1}d)
          -N\lambda_{ax_{l+1}\cdots x_{r-1}d}\bigr],\\
\label{eq:relative-right-side-exact}
 U_x^{\mathrm{rel}}-\mu_a(t_x+\lambda_x)
 &=C(a)-\mu_a
  -\sum_{r=l+1}^m\sum_{d>x_r}
   \bigl[C(ax_{l+1}\cdots x_{r-1}d)
          -N\lambda_{ax_{l+1}\cdots x_{r-1}d}\bigr].
\end{align}
Each side contains at most \(kK\) disjoint subcylinders of the
ambient cylinder.  Their means \(\mu_j\) therefore have sum at
most \(\mu_a\).  Cauchy--Schwarz bounds either side sum by
\[
 A\sum_j(\sqrt{\mu_j}+1)
 \le A\bigl[\sqrt{kK\mu_a}+kK\bigr].
\]
The additional ambient fluctuation in
\eqref{eq:relative-right-side-exact} is at most
\(A(\sqrt{\mu_a}+1)\).  Adding these bounds proves
\eqref{eq:relative-endpoint-finite-error}, including empty suffixes
and zero or small means.
\end{proof}

% Lean source: formalization/lower_bound_v2/RiffleLowerV2/Localization/Events.lean :: RiffleLowerV2.Localization.localization; RiffleLowerV2.Localization.coldLocalization_of_prefix_event
\begin{lemma}[Uniform cold-word localization]
% Lean: formalization/lower_bound_v2/RiffleLowerV2/Localization/Events.lean :: RiffleLowerV2.Localization.localization; RiffleLowerV2.Localization.coldLocalization_of_prefix_event
\label{lem:localization}
Use the finite witness, the selected prefix family \(\mathcal P_N\),
the full-word family \(\mathcal Y_N\), and the cold set
\(H_{{\rm rel},N}\).  Put
\[
 H_q=(q+H_{{\rm rel},N})\cap[N],\qquad q_*=L_{a_s},\qquad
 b=\beta^{(N)}(s),\qquad N\lambda_{a_s}=N^b.
\]
With probability \(1-o(1)\), simultaneously for every
\(y\in\mathcal Y_N\), all rows carrying \(y\) have sorted
positions in \(H_{q_*}\).  On this event each distinct full word
that occurs at least twice supplies a distinct edge of
\(E(G)\cap E(H_{q_*})\).  The probability bound is uniform over
the permitted finite witnesses and cut arrays.
\end{lemma}

\begin{proof}
Take \(\varepsilon=\eta/4000\) in
Lemma~\ref{cor:iterated-hypergeom}.  On its simultaneous
prefix-count event, Lemma~\ref{lem:finite-relative-endpoints} bounds
both relative endpoint errors for every \(x\in\mathcal P_N\)
by
\[
 e_N=N^{\eta/4000}
       \bigl[(2\sqrt{kK}+1)N^{b/2}+2kK+1\bigr].
\]
Since \(J_p^*\ge0\), the inequality
\(\ell_p=b-J_p^*\ge b/2+\eta/500\) in
\eqref{eq:repair-finite-geometric-buffers} implies \(b\ge\eta/250\).
As \(K\le T\log N\), uniformly for all these witnesses,
\begin{equation}
\label{eq:repair-endpoint-error-concrete}
 e_N\le N^{b/2+\eta/2000}
 \quad\hbox{for all sufficiently large \(N\)}.
\end{equation}

The reserve words \(v_b^-\) and \(v_b^+\) differ from the central
bridge \(v_b\) at only the first and second distinguished good
column, respectively.  At the distinguished columns their digit
probabilities are at least \(\chi_g\), and all other bridge digits
are maximal.  Hence their bridge products are at least
\(\chi_g^2k^{-b_N}\), where
\(b_N=\lceil 2\rho\log N\rceil\) is the bridge length.  By
\eqref{eq:repair-bin-exact} and
\eqref{eq:repair-finite-geometric-buffers}, every reserve prefix has
mean at least
\[
 \mu_x^\pm:=\E C(a_sxv_b^\pm)
 \ge e^{-1}\chi_g^2 k^{-1}
       N^{\ell_p-2\rho\log k}
 \ge N^{b/2+\eta/1000}
\]
for all sufficiently large \(N\), because
\(2\rho\log k\le\eta/100000\).  The fixed factor involving
\(\chi_g\) is absorbed by the remaining positive exponent gap.

The same simultaneous prefix event already controls every reserve count.
For a reserve mean $\mu\ge N^{b/2+\eta/1000}$, use $b\ge\eta/250$ and,
for sufficiently large $N$, $16\le N^{\eta/400}$ to obtain
\[
 N^{\eta/4000}(\sqrt\mu+1)\le\mu/2.
\]
Indeed $\mu\ge16N^{\eta/2000}$, so $\sqrt\mu\ge4N^{\eta/4000}$;
also $\sqrt\mu\ge1$. The prefix deviation bound gives $C(v)\ge\mu/2$.
The further inequality $4\le N^{\eta/2000}$ and the bound on $e_N$ yield
\begin{equation}\label{eq:repair-reserve-count-dominates}
 C(a_sxv_b^-),\ C(a_sxv_b^+)\ge\tfrac12N^{b/2+\eta/1000}
                           \ge e_N+2
 \qquad(x\in\mathcal P_N).
\end{equation}
Thus one prefix-count event controls both endpoints and all reserves,
with failure probability at most
$2(K+1)k^K\exp(-N^{\eta/4000}/8)$. Its threshold depends only on the
fixed parameters, since it is enough that
$4kT\log N+4\le N^{\eta/4000}$, $4\le N^{\eta/2000}$, and
$16\le N^{\eta/400}$.

Work on this event.For a fixed \(x\), write
\[
 \ell=N\lambda_{a_s}t_x^{\rm rel},\qquad
 r=N\lambda_{a_s}(t_x^{\rm rel}+\lambda_x).
\]
The relative endpoints differ from \(\ell,r\) by at most
\(e_N\).  Every one-based sorted row rank \(i\) carrying a word
\(a_sxv_bz\in\mathcal Y_N\) satisfies
\[
 q_*+L_x^{\mathrm{rel}}+C(a_sxv_b^-)+1
 \le i\le
 q_*+U_x^{\mathrm{rel}}-C(a_sxv_b^+),
\]
because \(v_b^-<_{\rm lex}v_b<_{\rm lex}v_b^+\).
Equation~\eqref{eq:repair-reserve-count-dominates} therefore gives
\(q_*+\ell<i<q_*+r\).  Thus
\[
 i\in q_*+(\mathbb Z\cap[\ell,r))\subseteq H_{q_*}.
\]
The two spare units in the reserve bound resolve the one-based-rank
and half-open-interval conventions.  No union bound over tail words
\(z\) and no estimate on the random translation \(q_*\) are needed.

Copies of a complete word are consecutive after lexicographic sorting.
If there are at least two, choose the edge between its first two copies.
The localization places both ranks in \(H_{q_*}\), and this pair
is an edge of \(G\).  Distinct words occupy disjoint constant-word
blocks, so their chosen edges are distinct.
\end{proof}
\subsubsection{Collision counts}

For a full word $w$, write $N_w=C(w)$ for its occupancy and
$\mu_w=N\lambda_w$ for its mean.

% Lean source: formalization/lower_bound_v2/RiffleLowerV2/Factorial/Exact.lean :: RiffleLowerV2.Factorial.exactJointFactorial
% Lean source: formalization/lower_bound_v2/RiffleLowerV2/Factorial/RelativeTargets.lean :: RiffleLowerV2.Factorial.relativeJointFactorial
\begin{lemma}[Fixed-column factorial moments]
% Lean: formalization/lower_bound_v2/RiffleLowerV2/Factorial/Exact.lean :: RiffleLowerV2.Factorial.exactJointFactorial
% Lean: formalization/lower_bound_v2/RiffleLowerV2/Factorial/RelativeTargets.lean :: RiffleLowerV2.Factorial.relativeJointFactorial
\label{lem:fixed-column-factorial}
For arbitrary $K$, distinct full words $w_1,\ldots,w_q$, and nonnegative
integers $r_i$ with $R=\sum_i r_i\le N$, let $m_{j,a}=\sum_{i:(w_i)_j=a}r_i$.
Then the exact identity is
\[
 \E\prod_i(N_{w_i})_{r_i}
   =(N)_R\prod_{j=1}^K\frac{\prod_a(n_a^{(j)})_{m_{j,a}}}{(N)_R}.
\]
In addition, suppose $K\le C\log N$ and every pile requested by the distinct full words
$w_1,\ldots,w_q$ has size at least $cN$, for fixed $c,C>0$.
For nonnegative integers $r_i$ with $R=\sum_i r_i\le N^{1/5}$,
\begin{equation}
\label{eq:exact-fixed-column-factorials}
 \E\prod_{i=1}^q(N_{w_i})_{r_i}
 =\prod_{i=1}^q\mu_{w_i}^{r_i}
    \left(1+O_{c,C}\left(\frac{(K+1)R^2}{N}\right)\right).
\end{equation}
The bound is uniform in the words and the column counts.
\end{lemma}

\begin{proof}
Expanding the falling factorials chooses ordered labelled rows for each
word. As the words are distinct, only disjoint choices contribute. There
are $(N)_R$ such choices. Let $m_{j,a}$ be the number of their requests for
digit $a$ in column $j$; thus $\sum_a m_{j,a}=R$.
Uniform sampling of each fixed-count column gives the exact formula
\[
 \E\prod_i(N_{w_i})_{r_i}
 =(N)_R\prod_{j=1}^K
                \frac{\prod_a(n_a^{(j)})_{m_{j,a}}}{(N)_R}.
\]
For $n\ge cN$ and $0\le v\le R\le N^{1/5}$,
\[
 \log\frac{(n)_v}{n^v}
 =\sum_{h=0}^{v-1}\log(1-h/n)=O_c(v^2/N),
\]
uniformly for sufficiently large $N$. Apply this also with $n=N$,
and compare the exact formula with
\[
 \prod_i\mu_{w_i}^{r_i}
   =N^R\prod_{j,a}(n_a^{(j)}/N)^{m_{j,a}}.
\]
Since $\sum_a m_{j,a}^2\le R^2$, the logarithm of their ratio is
$O_c((K+1)R^2/N)=o(1)$. Exponentiating proves
\eqref{eq:exact-fixed-column-factorials}. The case $R=0$ is equality.
\end{proof}

% Lean source: formalization/lower_bound_v2/RiffleLowerV2/Factorial/Central.lean :: RiffleLowerV2.Factorial.centralComparison; RiffleLowerV2.Factorial.centralComparisonOne
\begin{lemma}[Finite factorial comparison with a controlled remainder]
% Lean: formalization/lower_bound_v2/RiffleLowerV2/Factorial/Central.lean :: RiffleLowerV2.Factorial.centralComparison; RiffleLowerV2.Factorial.centralComparisonOne
\label{lem:central-window}
Let $X,Y$ be nonnegative integer-valued random variables on a finite probability space, let $\mu,\nu>0$,
and let $P_\mu,P_\nu$ be independent Poisson variables. Suppose that $m\ge2$
is even and, for $r,s\ge0$ with $r+s\le2m$,
\[
 \left|\E(X)_r(Y)_s-\mu^r\nu^s\right|
 \le\epsilon\mu^r\nu^s,\qquad 0\le\epsilon\le1.
\]
Set $f(t)=1-(1+t)e^{-t}$ and $d_m(t)=(m-1)t^m/m!$. Then
\begin{align*}
 \left|\Pp(X\ge2,Y\ge2)-f(\mu)f(\nu)\right|
 &\le \epsilon(1+\mu)(1+\nu)e^{\mu+\nu}\\
 &\quad+3\bigl(d_m(\mu)+d_m(\nu)+d_m(\mu)d_m(\nu)\bigr).
\end{align*}
If only the one-variable moments through order $m$ are assumed, then
\[
 \left|\Pp(X\ge2)-f(\mu)\right|
 \le\epsilon(1+\mu)e^\mu+3d_m(\mu).
\]
\end{lemma}

\begin{proof}
Define the even Bonferroni polynomial and its last term by
\[
 U_m(n)=\sum_{r=2}^{m}(-1)^r(r-1)\binom nr,
 \qquad D_m(n)=(m-1)\binom nm.
\]
The finite binomial identities give the pointwise inequalities
\begin{equation}
\label{eq:collision-bonferroni-sandwich}
 0\le U_m(n)-\mathbf 1_{\{n\ge2\}}\le D_m(n)
 \qquad(n\ge0).
\end{equation}
For completeness, put $S_j(n)=\sum_{r=2}^{j}(-1)^r(r-1)\binom nr$.
For $n>j\ge2$,
\[
 S_j(n)-1=(-1)^j
 \left(n\binom{n-2}{j-1}-\binom{n-1}{j}\right),
\]
and the expression in parentheses equals
$((j-1)n+1)\binom{n-2}{j-1}/j\ge0$.
For $2\le n\le j$, $S_j(n)=1$; for $n=0,1$, it is zero.
Thus $S_{m-1}(n)\le\mathbf 1_{\{n\ge2\}}\le S_m(n)$; the case $m=2$
uses $S_1=0$. Their difference is $D_m(n)$, proving
\eqref{eq:collision-bonferroni-sandwich}.

Writing $I_X=\mathbf 1_{\{X\ge2\}}$ and similarly for $Y$, the same inequality gives
\[
 0\le U_m(X)U_m(Y)-I_XI_Y
 \le D_m(X)+D_m(Y)+D_m(X)D_m(Y).
\]
Consequently the expected error is at most
$(1+\epsilon)(d_m(\mu)+d_m(\nu)+d_m(\mu)d_m(\nu))$.
The independent Poisson pair has the same bound with $\epsilon=0$.
The difference between the expectations of the two polynomial products is
at most
\[
 \epsilon\left(\sum_{r=2}^{m}\frac{(r-1)\mu^r}{r!}\right)
 \left(\sum_{s=2}^{m}\frac{(s-1)\nu^s}{s!}\right)
 \le\epsilon(1+\mu)(1+\nu)e^{\mu+\nu}.
\]
Adding the three bounds proves the assertion. The one-variable assertion
follows from the same sandwich. All expectations here are of polynomials of
total degree at most $2m$: in particular their values on $\{X+Y>2m\}$ have
been controlled, rather than discarded using only the probability of that event.
\end{proof}

% Lean source: formalization/lower_bound_v2/RiffleLowerV2/Occupancy/Main.lean :: RiffleLowerV2.Occupancy.uniformOccupancy; RiffleLowerV2.Occupancy.collisionChebyshev
\begin{lemma}[Uniform fixed-column collision probabilities]
% Lean: formalization/lower_bound_v2/RiffleLowerV2/Occupancy/Main.lean :: RiffleLowerV2.Occupancy.uniformOccupancy; RiffleLowerV2.Occupancy.collisionChebyshev
\label{lem:two-category-poisson}
Fix $c>0$ and $C>0$. For each $N$, let $K\le C\log N$, and take $K$
independent uniformly permuted fixed-count columns on the same $N$ labelled
rows. Let $\mathcal W$ be any deterministic family of distinct full words
such that every pile selected by any word in $\mathcal W$ has size at least
$cN$. Put
\[
 \lambda_w=\prod_{t=1}^{K}\frac{n_{w_t}^{(t)}}N,
 \quad \mu_w=N\lambda_w,
 \quad A_w=\{N_w\ge2\},\quad p_w=\Pp(A_w),
 \quad f(t)=1-(1+t)e^{-t}.
\]
Uniformly over all such data and words, as $N\to\infty$,
\begin{align}
 p_w&=f(\mu_w)\bigl(1+O_{c,C}(N^{-1/100})\bigr),
 \label{eq:uniform-collision-marginal}\\
 \Pp(A_w\cap A_{w'})
 &=p_wp_{w'}\bigl(1+O_{c,C}(N^{-1/100})\bigr),\qquad w\ne w'.
 \label{eq:uniform-collision-pair}
\end{align}
In particular $p_w\asymp\min\{\mu_w^2,1\}$. For
$Z=\sum_{w\in\mathcal W}\mathbf 1_{A_w}$,
\begin{equation}
\label{eq:uniform-collision-variance}
 \Var Z\le\E Z+O_{c,C}(N^{-1/100})(\E Z)^2.
\end{equation}
No upper or lower restriction on $\mu_w$ is required.
\end{lemma}

\begin{proof}
All constants below depend only on $c,C$. The exact labelled-row calculation
in Lemma~\ref{lem:fixed-column-factorial} gives, for distinct words $w,w'$
and $r+s\le R\le N^{1/5}$,
\begin{equation}
\label{eq:collision-uniform-factorials}
 \E(N_w)_r(N_{w'})_s
 =\mu_w^r\mu_{w'}^s\left(1+O\left(\frac{(K+1)R^2}{N}\right)\right).
\end{equation}
The same one-variable estimate holds in every model obtained by forcing
$b\le3$ labelled rows to carry a word in $\mathcal W$ and deleting these
rows: the residual population is $N-b$, the columns remain independent
uniform fixed-count columns, and every selected residual pile has size at
least $(c/2)(N-b)$ for all sufficiently large $N$.

For the occupancy $V$ of a fixed word in any of these original or residual
models, let $v=\E V$. Equation~\eqref{eq:thinning-transform} gives
\begin{equation}
\label{eq:collision-empty-singleton-tail}
 \E s^V\le e^{-(1-s)v}\quad(0<s\le1),\qquad
 \Pp(V<2)\le2e^{-v/2}.
\end{equation}
The latter bound is also \eqref{eq:thinning-lower-tail}.

Choose once and for all
\[
 \tau=N^{-1/100},\qquad B=\frac1{10}\log N,
 \qquad R=\lfloor N^{1/25}\rfloor,
 \qquad m=2\lfloor R/4\rfloor.
\]
For large $N$, $m\ge2$ is even and $2m\le R\le N^{1/5}$.
The relative factorial error in \eqref{eq:collision-uniform-factorials}
is bounded by
\[
 \epsilon=O_{c,C}((\log N)N^{-23/25}).
\]
If $0<\mu,\nu\le B$, Lemma~\ref{lem:central-window} therefore gives
\begin{equation}
\label{eq:collision-central-absolute}
 \left|\Pp(N_w\ge2,N_{w'}\ge2)-f(\mu)f(\nu)\right|
 =O_{c,C}(N^{-7/10}).
\end{equation}
Its one-variable version gives the same bound whenever $0<\mu\le2B$,
including the residual models on $N-b$ rows, $b\le3$.
To verify these errors explicitly, the moment-error term is at most
$C(\log N)^3N^{-18/25}=O(N^{-7/10})$ in the two-variable range,
and at most $C(\log N)^2N^{-18/25}$ in the one-variable range $\mu\le2B$.
Moreover,
\[
 d_m(t)\le m\exp(2t-m\log2)\le N^{-1},\qquad 0<t\le2B,
\]
for all sufficiently large $N$. Indeed $m\le N$, $m\log2\ge3\log N$, and $2t\le(2/5)\log N$ suffice for this estimate. The product of two such terms is at most $N^{-2}$. Changing $N$ to $N-b$ changes only the constants.

We first prove \eqref{eq:uniform-collision-marginal}.
If $\mu\le\tau$, the order-two/order-three Bonferroni inequalities are
\begin{equation}
\label{eq:collision-rare-bonferroni}
 \frac12\E(V)_2-\frac13\E(V)_3
 \le\Pp(V\ge2)\le\frac12\E(V)_2.
\end{equation}
The coefficient of $(V)_3$ is $1/3$.
Equation~\eqref{eq:collision-uniform-factorials} thus gives
$p_w=\mu^2(1+O(\tau+(K+1)/N))/2$.
Since $f(\mu)=\mu^2(1+O(\mu))/2$, the desired relative estimate follows,
uniformly even when $\mu$ is arbitrarily small.
If $\tau<\mu\le B$, then $f(\mu)\ge(e^{-1}/2)\min\{\mu^2,1\}\ge(e^{-1}/2)\tau^2$.
The one-variable central bound has relative error
$O(N^{-7/10}\tau^{-2})=O(N^{-17/25})$, again smaller than $O(\tau)$.
If $\mu>B$, \eqref{eq:collision-empty-singleton-tail} gives
$1-p_w\le2N^{-1/20}$, whereas
$1-f(\mu)=(1+\mu)e^{-\mu}\le(1+B)e^{-B}=O((\log N)N^{-1/10})$.
This proves \eqref{eq:uniform-collision-marginal} in all three regimes.
The same conclusion holds uniformly in the residual models, by the verified
population and pile-size bounds above.

This marginal result also provides the conditional stability needed when
one of the two categories is rare. Fix $b=2$ or $3$ distinct labelled rows
and force them to carry $w$. For $w'\ne w$, none of the forced rows counts
as a $w'$ row. The conditional mean of the remaining $w'$ count is
\[
 \nu_b=(N-b)\prod_{t=1}^{K}
 \frac{n_{w'_t}^{(t)}-b\mathbf 1_{\{w'_t=w_t\}}}{N-b},
 \qquad \frac{\nu_b}{\nu}=1+O((K+1)/N),\quad\nu=\mu_{w'}.
\]
All these conditioning events are positive for large $N$ because their
requested piles have size at least $cN$. The function $f$ satisfies
\[
 0\le\frac{tf'(t)}{f(t)}\le2\qquad(t>0),
\]
since $f(t)=\int_0^t u e^{-u}\,du\ge e^{-t}t^2/2$.
Equivalently, $f$ is increasing and $f(t)/t^2$ is decreasing. Thus
$0<x\le y$ implies $f(x)\le f(y)\le(y/x)^2f(x)$. Applying these two
inequalities to the original and residual means gives
$f(\nu_b)/f(\nu)=1+O((K+1)/N)$ uniformly for every $\nu>0$.Applying \eqref{eq:uniform-collision-marginal} to the original and residual
models proves
\begin{equation}
\label{eq:collision-conditional-stability}
 \Pp(A_{w'}\mid\text{the prescribed $b$ rows carry $w$})
 =p_{w'}(1+O(\tau)),\qquad b=2,3.
\end{equation}
The constants are independent of the row labels and both words.

To prove \eqref{eq:uniform-collision-pair}, first suppose
$\mu=\mu_w\le\tau$; no restriction is imposed on $\nu=\mu_{w'}$.
Summing \eqref{eq:collision-conditional-stability} over the ordered labelled
$b$-tuples gives
\[
 \E[(N_w)_b\mathbf 1_{A_{w'}}]
 =\mu^b p_{w'}(1+O(\tau)),\qquad b=2,3.
\]
Multiplying the pointwise Bonferroni sandwich
\eqref{eq:collision-rare-bonferroni} by $\mathbf 1_{A_{w'}}$ therefore yields
\[
 \Pp(A_w\cap A_{w'})=\frac{\mu^2}2p_{w'}(1+O(\tau))
 =p_wp_{w'}(1+O(\tau)).
\]
This includes the case where both means are arbitrarily small.

It remains to consider $\mu,\nu>\tau$. If, say, $\nu>B$, then
\eqref{eq:collision-empty-singleton-tail} implies
\[
 |\Pp(A_w\cap A_{w'})-p_wp_{w'}|
 \le2\Pp(A_{w'}^c)\le4N^{-1/20}.
\]
Since $p_w\ge(e^{-1}/4)\tau^2$ and $p_{w'}\ge1/2$ for large $N$, its relative
error is $O(N^{-3/100})=O(\tau)$.
Finally, if $\tau<\mu,\nu\le B$, use
\eqref{eq:collision-central-absolute} and the already proved marginal
relative estimates. Here $f(\mu)f(\nu)\ge(e^{-2}/4)\tau^4$, so the additional
central relative error is
$O(N^{-7/10}\tau^{-4})=O(N^{-33/50})$.
This proves the pair estimate uniformly in every case.

It follows that
\[
 \sum_{w\ne w'}|\operatorname{Cov}(\mathbf 1_{A_w},\mathbf 1_{A_{w'}})|
 \le C\tau\left(\sum_w p_w\right)^2.
\]
Adding $\sum_w\Var(\mathbf 1_{A_w})\le\E Z$ proves
\eqref{eq:uniform-collision-variance}.
\end{proof}

% Lean source: formalization/lower_bound_v2/RiffleLowerV2/EmpiricalLower/ManyEdges.lean :: RiffleLowerV2.EmpiricalLower.manyEdges
% Lean source: formalization/lower_bound_v2/RiffleLowerV2/EmpiricalLower/CollisionMean.lean :: RiffleLowerV2.EmpiricalLower.cold_collision_mean_lower
\begin{lemma}[Many localized collision edges]
% Lean: formalization/lower_bound_v2/RiffleLowerV2/EmpiricalLower/ManyEdges.lean :: RiffleLowerV2.EmpiricalLower.manyEdges
% Lean: formalization/lower_bound_v2/RiffleLowerV2/EmpiricalLower/CollisionMean.lean :: RiffleLowerV2.EmpiricalLower.cold_collision_mean_lower
\label{lem:many-edges}
For the finite witness and cold set of Lemma~\ref{lem:cold-geometry},
the following estimate holds:
\[
 \Pp\{|E(G)\cap E(H_{L_{a_s}})|\ge M_N^{1/2+\eta/5000}\}
 \longrightarrow1.
\]
The convergence is uniform over the permitted finite witnesses and cuts.
\end{lemma}

\begin{proof}
Let $A_y=\{N_y\ge2\}$ and $Z=\sum_{y\in\mathcal Y_N}\mathbf 1_{A_y}$.
Every selected pile has size at least $cN$ for a fixed $c>0$, by
Lemma~\ref{lem:cold-geometry}. Hence
Lemma~\ref{lem:two-category-poisson} applies to this whole deterministic
family. Its marginal bound and \eqref{eq:cold-full-word-masses} give
\[
 \Pp(A_y)\ge \frac{e^{-1}}4\min\{\mu_y^2,1\}
 \ge \frac{e^{-5}}4N^{2z_N-4\min\{\eta/100000,1/1000\}},\qquad
 \E Z\ge \frac{e^{-5}}4N^{g_N-4\min\{\eta/100000,1/1000\}},
\]
Here we used the upper bound on $z_N$ in \eqref{eq:cold-effective-collision}.
The other bound there gives
\[
 g_N-h_N/2-4\min\{\eta/100000,1/1000\}
 \ge\eta/2000-\eta/25000>\eta/3000.
\]
Together with \eqref{eq:cold-geometric-conclusions} this yields, for all
sufficiently large $N$,
\[
 \E Z\ge M_N^{1/2}N^{\eta/3000},\qquad \E Z\longrightarrow\infty.
\]
Equation~\eqref{eq:uniform-collision-variance} and Chebyshev's inequality
now imply
\[
 \Pp(Z<\tfrac12\E Z)
 \le\frac4{\E Z}+O(N^{-1/100})=o(1).
\]
In particular $Z\ge M_N^{1/2}N^{\eta/4000}$ with probability $1-o(1)$.
Intersect this event with that of Lemma~\ref{lem:localization}. Its
injective assignment of edges gives
$|E(G)\cap E(H_{L_{a_s}})|\ge Z$. Since $M_N\le N$, any fixed
$0<\delta<\eta/4000$, for example $\delta=\eta/5000$, proves the claim.
\end{proof}






\subsubsection{Completion of the empirical lower bound}

% Lean source: formalization/lower_bound_v2/RiffleLowerV2/EmpiricalLower/Cold.lean :: RiffleLowerV2.EmpiricalLower.coldOnly
% Lean source: formalization/lower_bound_v2/RiffleLowerV2/EmpiricalLower/Cold.lean :: RiffleLowerV2.EmpiricalLower.coldOnly
\begin{proposition}[Cold-spot lower bound]% Lean: formalization/lower_bound_v2/RiffleLowerV2/EmpiricalLower/Cold.lean :: RiffleLowerV2.EmpiricalLower.coldOnly
\label{prop:cold-lower}
Fix $\eta,\chi_g,\gamma,\rho>0$ with
\[
 2\rho\log k\le\frac1{100}\min\{\eta/100000,1/1000\}.
\]
For every $\varepsilon>0$, there exists $N_0\ge2$ such that, for every
$N\ge N_0$ and every cut array satisfying the local good-column condition
with parameters $\chi_g,\gamma,\rho$ on $[0,T]$, the inequalities
\[
 \Lambda_1^{(N)}\ge\eta,\qquad
 \beta^{(N)}(S_N)<\min\{\eta/100000,1/1000\}
\]
imply $\dTV(P_N^K,\nuUnif)\ge1-\varepsilon$.
\end{proposition}

\begin{proof}
Lemma~\ref{lem:cold-witness} supplies the finite witness, and
Lemma~\ref{lem:cold-geometry} supplies its deterministic cold set. Its
cardinality is at least $N^{\eta/3000}$, and every clipped translate has
boundary at most four times the square root of that cardinality.
Lemma~\ref{lem:many-edges} gives at least $M_N^{1/2+\eta/5000}$ internal
graph edges on the translate indexed by $L_{a_s}$, with probability tending
to one uniformly. Apply Lemma~\ref{lem:lalley} with $a=\eta/3000$,
$\delta=\eta/5000$, and $D=4$. The thresholds for the geometry, collisions,
localization, and ascent test are all fixed before the cut array. Taking
their maximum gives the asserted $N_0$.
\end{proof}

% Lean source: formalization/lower_bound_v2/RiffleLowerV2/EmpiricalLower/Main.lean :: RiffleLowerV2.EmpiricalLower.empiricalCombined
% Lean source: formalization/lower_bound_v2/RiffleLowerV2/EmpiricalLower/Main.lean :: RiffleLowerV2.EmpiricalLower.empiricalCombined


\begin{proof}[Proof of Theorem~\ref{thm:empirical-lower}]
The first alternative follows from Proposition~\ref{prop:largest-prefix}
and the exact maximal-word identity. In the second alternative, if
$\beta^{(N)}(S_N)\ge\min\{\eta_1/100000,1/1000\}$, apply the same
proposition with this fixed positive exponent. If the reverse strict
inequality holds, apply Proposition~\ref{prop:cold-lower} with
$\eta=\eta_1$. Take the maximum of the three thresholds. When the cold
parameters are not positive or fail the displayed restriction, only the
first alternative can occur, and its threshold alone suffices.
\end{proof}

\subsubsection{Passage to a macroscopic profile}

We now apply the profile transfer proved in Proposition~\ref{prop:profile-transfer}.
Its interval and margin bounds, \eqref{eq:interval-transfer} and
\eqref{eq:margin-transfer}, use the same selected finite partition in the
almost-profile condition and the scalar integrability hypotheses. In particular,
\[
 |\beta^{(N)}(S_N)-\Lambda_0(T)|\le e+\frac{\log k}{\log N},\qquad
 |\Lambda_1^{(N)}-\Lambda_1(T)|\le\frac{13e}{2}+\frac{14\log k}{\log N}.
\]
% Lean source: formalization/lower_bound_v2/RiffleLowerV2/MacroscopicLower/Main.lean :: RiffleLowerV2.MacroscopicLower.macroscopicIntegrable
% Lean source: formalization/lower_bound_v2/RiffleLowerV2/MacroscopicLower/Main.lean :: RiffleLowerV2.MacroscopicLower.macroscopicIntegrable_tendsto
% Lean source: formalization/lower_bound_v2/RiffleLowerV2/MacroscopicLower/Main.lean :: RiffleLowerV2.MacroscopicLower.macroscopicIntegrable; RiffleLowerV2.MacroscopicLower.macroscopicIntegrable_tendsto


\begin{proof}[Proof of Theorem~\ref{thm:macroscopic-lower}]
Choose
\[
 0<e<\min\left\{\frac1{10}\min\{\eta/200000,1/1000\},
                          \frac{\eta}{100}\right\}
\]
and apply Proposition~\ref{prop:profile-transfer} with a mesh upper bound
small enough for the displayed restriction on $\rho_0$. All these choices
precede the profile and the cut array. If
$\beta^{(N)}(S_N)\ge\min\{\eta/200000,1/1000\}$, the largest-word
proposition applies with this fixed positive exponent. Otherwise,
\eqref{eq:margin-transfer} gives $\Lambda_0(T)<\eta$ for
large $N$. The assumed maximum therefore implies $\Lambda_1(T)\ge\eta$.
Equation~\eqref{eq:margin-transfer} then gives
$\Lambda_1^{(N)}\ge\eta/2$ for large $N$. Proposition~\ref{prop:cold-lower}
with margin $\eta/2$ applies in this second case. Take the maximum of the
transfer and the two lower-bound thresholds. The inequalities
$0\le\dTV\le1$ give the asserted convergence for eventually admissible
sequences.
\end{proof}







\begin{thebibliography}{1}
\bibitem{SSW}
Mark Sellke, Jialu Shi, and Jiamin Wang.
\newblock Universality of cutoff for riffle shuffling.
\newblock \href{https://arxiv.org/abs/2510.22783v1}{arXiv:2510.22783v1}, 2025.
\end{thebibliography}
\end{document}
